\documentclass[12pt,a4paper]{article}
\usepackage[utf8]{inputenc}
\usepackage[T1]{fontenc}
\usepackage{amsmath,amssymb,amsthm}
\usepackage{hyperref}

\theoremstyle{definition}
\newtheorem{definition}{Definition}

\title{Chicago Notes -- MathGloss Definitions}
\author{Auto-generated for KnowTex Benchmark}
\date{}

\begin{document}
\maketitle

\begin{definition}[algebra of subsets]\label{def:algebra-of-subsets}
A \emph{field} or \emph{algebra} of subsets of a set $X$ is a pair $(X, \mathcal F)$ consisting of the set $X$ and a collection $\mathcal F$ of subset of $X$ such that
1. $\mathcal F$ is closed under complementation: $X\setminus F \in \mathcal F$ for all $F \in \mathcal F$;
2. $\mathcal F$ contains the empty set. 
3. $\mathcal F$ is closed under binary union (with induction and De Morgan's laws, this is euqivalent to being closed under binary intersection, finite union, and finite intersection.)
\end{definition}

\begin{definition}[automorphism]\label{def:automorphism}
An \emph{automorphism} is an isomorphism from an algebraic structure to itself.
\end{definition}

\begin{definition}[barycenter]\label{def:barycenter}
The \emph{barycenter} of an n-simplex $[v_0,\dots,v_n]$ is the point $b=\sum_i t_iv_i$ whose barycentric coordinates $t_i$ are all equal. Specifically, $t_i = \frac{1}{(n+1)}$ for each $i$.
\end{definition}

\begin{definition}[binary operation]\label{def:binary-operation}
A \emph{binary operation} $\cdot$ on a set $A$ is a function $A\times A\to A$.
\end{definition}

\begin{definition}[associative]\label{def:associative}
A binary operation~\ref{def:binary-operation} $\cdot$ on the set $A$ is \emph{associative} if for all $a,b,c \in A$, the following relationship holds: $$a\cdot (b\cdot c) = (a\cdot b)\cdot c.$$
\end{definition}

\begin{definition}[binary relation]\label{def:binary-relation}
A \emph{binary relation} $R$ on sets $X$ and $Y$ is a subset of $X \times Y$. If $(x,y) \in R$, then we say $xRy$ or "$x$ is related to $y$."
\end{definition}

\begin{definition}[commutative]\label{def:commutative}
A binary operation~\ref{def:binary-operation} $\cdot$ on the set $A$ is \emph{commutative} if for all $a,b\in A$, $$ab = ba.$$
\end{definition}

\begin{definition}[complex numbers]\label{def:complex-numbers}
The \emph{complex numbers} are given by ordered pairs of real numbers: $$\mathbb C = \{(a,b)\mid a,b\in \mathbb R\}$$ and endowed with two binary operations~\ref{def:binary-operation}: multiplication and addition. 

Addition is defined pairwise, while multiplication is defined as $$(a,b)*(c,d) = (ac-bd,bc+ad)$$ for all $(a,b),(c,d) \in \mathbb C$. 

We often write $a+bi$ for $(a,b)$ where $i^2=-1$.
\end{definition}

\begin{definition}[argument]\label{def:argument}
Let $z\in \mathbb C$ be a complex number~\ref{def:complex-numbers}. Then the \emph{argument} of $z$, denoted $\text{arg}(z)$, is the angle between the vector in $\mathbb R^2$ corresponding to $z$ and the positive real axis. Equivalently, it is any real number $\theta$ such that $z= r(\cos\theta+i\sin\theta) = re^{i\theta}$, via Euler's formula.
\end{definition}

\begin{definition}[complex conjugate]\label{def:complex-conjugate}
The \emph{complex conjugate} of a complex number~\ref{def:complex-numbers} $z = x+iy$ is $\overline z = x-iy$.
\end{definition}

\begin{definition}[endomorphism]\label{def:endomorphism}
An \emph{endomorphism} is a homomorphism from an object to itself.
\end{definition}

\begin{definition}[Euler characteristic of a graph]\label{def:euler-characteristic-of-a-graph}
The Euler characteristic of a graph is the difference $V-E$ between the number of vertices and the number of edges.
\end{definition}

\begin{definition}[holomorphic]\label{def:holomorphic}
Let $f:\mathbb C\to\mathbb C$, and write it as follows: $$f(z) = u(x,y) + iv(x,y)$$ where $z\in \mathbb C$ is such that $z = x+iy$ and $u,v:\mathbb R^2\to \mathbb R$. Then $f$ is \emph{complex differentiable} at $z_0\in \mathbb C$ with derivative $f'(z_0)$ if the limit $$f'(z_0) = \lim_{z\to z_0} \frac{f(z)-f(z_0)}{z-z_0}$$ exists.
\end{definition}

\begin{definition}[entire]\label{def:entire}
The function $f:\mathbb C\to \mathbb C$ is \emph{entire} if it is holomorphic~\ref{def:holomorphic} at each point in $\mathbb C$
\end{definition}

\begin{definition}[homogeneous relation]\label{def:homogeneous-relation}
A \emph{homogeneous relation} on a set $X$ is a binary relation~\ref{def:binary-relation} where the two sets being related are both $X$.
\end{definition}

\begin{definition}[image]\label{def:image}
The \emph{image} of a function $f:X\to Y$ is the set $$\text{im}(f) = \{T(x) \in Y \mid x \in X\}\subset Y.$$
\end{definition}

\begin{definition}[inconsistent]\label{def:inconsistent}
Let $\Gamma$ be a set of formulas in a language. We say that $\Gamma$ is \emph{inconsistent} if it is not consistent.
\end{definition}

\begin{definition}[injective]\label{def:injective}
A function $f:A\to B$ is \emph{injective} if for all $a,b \in A$, $a\neq b$ implies that $f(a)\neq f(b)$.
\end{definition}

\begin{definition}[interval]\label{def:interval}
The set $A\subset \mathbb R$ is an \emph{interval} if for all $a,b \in A$ with $a<b$ and $x\in \mathbb R$, $$a<x<b \implies x\in A.$$
\end{definition}

\begin{definition}[absolute continuity]\label{def:absolute-continuity}
A function $f:[a,b]\to \mathbb R^n$ is \emph{absolutely continuous} if for every $\varepsilon > 0$, there exists $\delta > 0$ such that if $\{(x_i,y_i)\}_{i=1}^k$ are pairwise disjoint subintervals~\ref{def:interval} of $[a,b]$ with total combined length less than $\delta$, then $$\sum_{i=1}^k \vert f(y_i) - f(x_i)\vert  < \varepsilon.$$
\end{definition}

\begin{definition}[complex differentiability]\label{def:complex-differentiability}
A function $f: [a,b] \to \mathbb C$ from a real interval~\ref{def:interval} to the complex numbers~\ref{def:complex-numbers} is \emph{differentiable} at $t$ if the limit of the difference quotient $$\lim_{h\to 0}\frac{f(t+h) - f(t)}{h}$$ exists. We call this limit $f'(t)$.
\end{definition}

\begin{definition}[inverse function]\label{def:inverse-function}
Let $f:A\to B$ be a function.

A \emph{left inverse} of $f$ is a function $g:B\to A$ such that $g\circ f(a) = a$ for all $a \in A$.

A \emph{right inverse} of $f$ is a function $g:B\to A$ such that $f\circ g(b) = b$ for all $b\in B$. 

An \emph{inverse} of $f$ is a function $f^{-1}$ that is both a left and a right inverse of $f$.
\end{definition}

\begin{definition}[Kronecker delta]\label{def:kronecker-delta}
The \emph{Kronecker delta} $\delta_{ij}$ is a function of the two variables $i$ and $j$ defined by $$\delta_{ij} = \begin{cases}0 & \text{if }i\neq j \\ 1 & \text{if } i=j.\end{cases}$$
\end{definition}

\begin{definition}[left identity element]\label{def:left-identity-element}
Let $X$ be a set with binary operation~\ref{def:binary-operation} $*$. An element $e\in X$ is a \emph{left identity} of $X$ if $x* a = a$ for all $a\in A$.
\end{definition}

\begin{definition}[locally monotonic]\label{def:locally-monotonic}
A function $f:A\to \mathbb R$ is \emph{locally monotonic} at $x \in A$ if there exists $r>0$ such that whenever $0 < y-x < r$, $f(y)\geq f(x)$ and whenever $0 < x-y < r$, $f(y)\leq f(x)$.
\end{definition}

\begin{definition}[magma]\label{def:magma}
A \emph{magma} is a set $X$ together with a binary operation~\ref{def:binary-operation} $*$ that is closed under $*$. That is, $a*b \in X$ for all $a,b \in X$.
\end{definition}

\begin{definition}[left coset]\label{def:left-coset}
Let $G$ be a magma~\ref{def:magma} with binary operation~\ref{def:binary-operation} $\cdot$ and sub-structure $H$. A \emph{left coset} of the $H$ for the element $x\in G$ is the set $x\cdot H = \{x\cdot h\mid h \in H\}$.
\end{definition}

\begin{definition}[matrix]\label{def:matrix}
An $m\times n$ (real or complex) \emph{matrix} $A$ is a collection of $mn$ \emph{entries} $(a_{ij})$ which are either real or complex numbers arranged in $m$ rows and $n$ columns as follows: $$\begin{bmatrix}a_{11} & a_{12} & \cdots & a_{1n} \\ a_{21} & a_{22} & \cdots & a_{2n} \\ \vdots&\vdots &\ddots &\vdots \\ a_{m1} & a_{m2} &\cdots & a_{mn}\end{bmatrix}$$
Two matrices of the same size may be added together entrywise.
\end{definition}

\begin{definition}[identity matrix]\label{def:identity-matrix}
The $n\times n$ \emph{identity matrix~\ref{def:matrix}} $I_n$ is the $n\times n$ matrix~\ref{def:matrix} whose entries $(a_{ij})$ are such that $a_{ij}=1$ when $i=j$ and $0$ otherwise. That is, $a_{ij} = \delta_{ij}$ where $\delta_{ij}$ is the Kronecker delta~\ref{def:kronecker-delta}.
\end{definition}

\begin{definition}[matrix multiplication]\label{def:matrix-multiplication}
For $A$ an $m\times n$ matrix~\ref{def:matrix} and $B$ an $n\times p$ matrix with entries $(a_{ij})$ and $(b_{k\ell})$ respectively, the \emph{product} $C=AB$ is the $m\times p$ matrix whose entries are given by $$c_{ij} = \sum_{k=1}^n a_{ik}b_{kj}$$ for $1 \leq i \leq m$ and $1 \leq j \leq m$. 

Note that this binary operation~\ref{def:binary-operation} is not commutative~\ref{def:commutative}, and $BA$ may not even be defined even though $AB$ is due to differing dimensions.
\end{definition}

\begin{definition}[inverse matrix]\label{def:inverse-matrix}
The \emph{inverse} of an $n\times n$ matrix~\ref{def:matrix} $A$ is an $n\times n$ matrix~\ref{def:matrix} $A^{-1}$ such that the products~\ref{def:matrix-multiplication} $$AA^{-1} = A^{-1}A = I_n$$ where $I_n$ is the $n\times n$ identity matrix~\ref{def:identity-matrix}.
\end{definition}

\begin{definition}[matrix transpose]\label{def:matrix-transpose}
The \emph{transpose} of an $m\times n$ matrix~\ref{def:matrix} $A = (a_{ij})$ is the $n\times m$ matrix~\ref{def:matrix} $A^T$ whose entries are given by $a_{ji}$.
\end{definition}

\begin{definition}[conjugate transpose]\label{def:conjugate-transpose}
The \emph{conjugate~\ref{def:complex-conjugate} matrix transpose~\ref{def:matrix-transpose}} of a complex $m\times n$ matrix~\ref{def:matrix} $A$ is the $n \times m$ matrix~\ref{def:matrix} $A^*$ given by taking the usual matrix transpose~\ref{def:matrix-transpose} $A^T$ and then setting $A^*$ to be the matrix~\ref{def:matrix} whose entries are the complex conjugates~\ref{def:complex-conjugate} of each entry of $A^T$.
\end{definition}

\begin{definition}[model]\label{def:model}
Let $\Gamma$ be a set of formulas in a language. A \emph{model} for $\Gamma$ is a truth function $t:\text{Form}(L) \to \{0,1\}$ such that $t(\gamma) = 1$ for all $\gamma\in \Gamma$. 

For any $\phi \in \text{Form}(L)$, we say that $\Gamma$ \emph{models} $\phi$ if for every model $t$ of $\Gamma$, $t(\phi) = 1$.
\end{definition}

\begin{definition}[ordinary generating function]\label{def:ordinary-generating-function}
The \emph{ordinary generating function} of a sequence $\{a_k\}_{k\in\mathbb N}$ is $$A(x) = \sum_{k=0}^\infty a_kx^k$$
\end{definition}

\begin{definition}[oriented edge]\label{def:oriented-edge}
An \emph{oriented edge} $k:I\to X$ in a graph $X$ where $I$ is the unit interval~\ref{def:interval} is the traversal of an edge in either the forward or the backward direction.
\end{definition}

\begin{definition}[edge path]\label{def:edge-path}
An \emph{edge path} in a graph $X$ is a finite composition of oriented edges~\ref{def:oriented-edge} $k_n$ with $k_{n+1}(0) = k_n(1)$.
\end{definition}

\begin{definition}[closed edge path]\label{def:closed-edge-path}
An edge path~\ref{def:edge-path} on a graph $X$ is \emph{closed} if it starts and ends at the same vertex.
\end{definition}

\begin{definition}[connected graph]\label{def:connected-graph}
A graph $X$ is \emph{connected} if there is an edge path~\ref{def:edge-path} connecting every vertex of $X$ to every other.
\end{definition}

\begin{definition}[p-adic integers]\label{def:p-adic-integers}
The \emph{$p$-adic integers} is the subset of Qp defined by $$\mathbb Z_p = \{x\in \mathbb Q_p\mid \vert x\vert _p \leq 1\}.$$
\end{definition}

\begin{definition}[partition of an integer]\label{def:partition-of-an-integer}
Let $n$ be a positive integer. A \emph{partition} of $n$ is a non-decreasing sequence of positive integers whose sum is $n$.
\end{definition}

\begin{definition}[power set]\label{def:power-set}
The \emph{power set} of a set $X$ is the connection of all subsets of $X$, including $X$ itself and the empty set. It is denoted $\mathcal P(X)$.
\end{definition}

\begin{definition}[outer measure]\label{def:outer-measure}
Let $\mathcal P(X)$ be the power set~\ref{def:power-set} of $X$. An \emph{outer measure} on $X$ is a function $\mu: \mathcal P(X) \to [0,\infty]$ such that
- $\mu(\emptyset) = 0$ 
- if $A, B \subset X$ with $A\subset B$, then $\mu(A) \leq\mu(B)$
- for arbitrary subsets $B_i \subset X$ for $i \in \mathbb N$, we have $\mu\left(\bigcup\limits_{i\in\mathbb N} B_i \right)\leq \sum\limits_{i\in\mathbb N}\mu(B_i)$.
\end{definition}

\begin{definition}[partition of a set]\label{def:partition-of-a-set}
A \emph{partition} of a set $X$ is a subset $\mathcal S$ of the power set~\ref{def:power-set} $\mathcal P(X)$ such that 
1. $A\in \mathcal S$ implies that $A\neq \emptyset$;
2. For every $x\in X$, there is a unique $A\in \mathcal S$ such that $x\in A$. (Equivalently, the union of the sets $A\in \mathcal S$ is $X$ and the intersection of any two elements of $\mathcal S$ is empty.)
\end{definition}

\begin{definition}[bounded variation]\label{def:bounded-variation}
A function $f: [a,b]\to \mathbb R$ is of \emph{bounded variation} if for every partition of a set~\ref{def:partition-of-a-set} $P =\{a= x_0 < <x_1 < \cdots < x_n = b\}$, we have $$\sup_P \sum_{i=1}^n \vert f(x_i)-f(x_{i-1})\vert  <\infty.$$
\end{definition}

\begin{definition}[rectangle]\label{def:rectangle}
A \emph{rectangle} $Q$ in $\mathbb R^n$ is a product of intervals~\ref{def:interval} in $\mathbb R$: $$Q = [a_1,b_1]\times \cdots\times [a_n,b_n].$$
\end{definition}

\begin{definition}[component interval]\label{def:component-interval}
A \emph{component interval} of a rectangle~\ref{def:rectangle} $Q \subset \mathbb R^n$ is one of the intervals $[a_i,b_i]$ that make up the product.
\end{definition}

\begin{definition}[mesh]\label{def:mesh}
Let $Q\subset\mathbb R^n$ be a rectangle~\ref{def:rectangle}. For any partition of a set~\ref{def:partition-of-a-set} $P$ of $Q$, the \emph{mesh} of $P$ is the maximum length of the component intervals~\ref{def:component-interval} of $P$.
\end{definition}

\begin{definition}[reduced edge path]\label{def:reduced-edge-path}
An edge path~\ref{def:edge-path} is \emph{reduced} if it is never the case that $k_{n+1}$ is $k_n$ with the opposite orientation.
\end{definition}

\begin{definition}[reflexive]\label{def:reflexive}
A binary relation~\ref{def:binary-relation} $R$ on a set $X$ (i.e. a subset of $X \times X$) is \emph{reflexive} if it relates every $x \in X$ to itself. That is, for all $x \in X$, $(x,x) \in R$.
\end{definition}

\begin{definition}[right identity element]\label{def:right-identity-element}
Let $X$ be a set with binary operation~\ref{def:binary-operation} $*$. An element $e\in X$ is a \emph{right identity} of $X$ if $a*x = a$ for all $a\in A$.
\end{definition}

\begin{definition}[identity element]\label{def:identity-element}
Let $X$ be a set together with a binary operation~\ref{def:binary-operation} $*$. An element $e \in X$ is an \emph{identity element} of $X$ if it is both a left~\ref{def:left-identity-element} and right identity~\ref{def:right-identity-element} of $X$ with respect to $*$.
\end{definition}

\begin{definition}[category]\label{def:category}
A \emph{category} $\mathcal C$ is a collection of \emph{objects} together with a set $\mathcal C(A,B)$ of \emph{morphisms} or maps between any two objects. In particular, there is always for each object $A$ in $\mathcal C$ an identity morphism $\text{id}_A \in\mathcal C(A,A)$. Finally, a category must obey the following composition law: $$\circ: \mathcal C(B,C)\times \mathcal C(A,B) \to \mathcal C(A,C)$$ for all triples $A$, $B$, $C$ of objects such that composition is associative~\ref{def:associative}, and identity morphisms  are identities~\ref{def:identity-element} for composition. That is, 
- $h\circ(g\circ f) = (h\circ g)\circ f$;
- $\text{id}\circ f = f$;
- $f\circ \text{id} = f$.
\end{definition}

\begin{definition}[category of sets]\label{def:category-of-sets}
The \emph{category of sets} is the category~\ref{def:category} whose objects are sets and whose morphisms are functions between sets.
\end{definition}

\begin{definition}[discrete category]\label{def:discrete-category}
A category~\ref{def:category} is \emph{discrete} if its morphisms are only the identity morphisms.
\end{definition}

\begin{definition}[epimorphism]\label{def:epimorphism}
Let $C$ be a category~\ref{def:category}. A morphism~\ref{def:category} $f\in \text{Hom}_C(A,B)$ is an \emph{epimorphism} if for all objects $Z$ of $C$ and all morphisms $\alpha,\beta\in \text{Hom}_C(B,Z)$, $$\alpha \circ f= \beta\circ f\implies \alpha = \beta.$$
\end{definition}

\begin{definition}[functor]\label{def:functor}
A covariant \emph{functor} $F:\mathcal C\to \mathcal D$ is a map between categories~\ref{def:category} $\mathcal C$ and $\mathcal D$. It must assign to each object $A$ of $\mathcal C$ an object $F(A)$ an object of $\mathcal D$ and to each morphism $f:A\to B$ of $\mathcal C$ a morphism $F(f):f(A) \to f(B)$ of $\mathcal D$ such that
- $F(\text{id}_A)) = \text{id}_{F(A)}$;
- $F(g\circ f) = F(g)\circ F(f)$.

A contravariant functor "reverses the direction of arrows;" that is, $F$ sends $f:A\to B$ to $F(f): f(B) \to f(A)$ and $F(g\circ f) = F(f) \circ F(g)$.
\end{definition}

\begin{definition}[left inverse element]\label{def:left-inverse-element}
Let $(X,*)$ be a magma~\ref{def:magma} with identity element~\ref{def:identity-element} $e$. An element $b\in X$ is a \emph{right inverse element} for $a\in X$ if $a * b=e$.
\end{definition}

\begin{definition}[monomorphism]\label{def:monomorphism}
Let $C$ be a category~\ref{def:category}. A morphism~\ref{def:category} $f\in \text{Hom}_C(A,B)$ is a \emph{monomorphism} if for all objects $Z$ of $C$ and all morphisms $\alpha,\beta\in \text{Hom}_C(Z,A)$, $$f\circ \alpha = f\circ \beta\implies \alpha = \beta.$$
\end{definition}

\begin{definition}[opposite category]\label{def:opposite-category}
The \emph{opposite category} $\mathcal C^{\text{op}}$ of a category~\ref{def:category} $\mathcal C$ is a category~\ref{def:category} consisting of the same objects as $\mathcal C$, but with morphisms "going in the opposite directions."
\end{definition}

\begin{definition}[right inverse element]\label{def:right-inverse-element}
Let $(X,*)$ be a magma~\ref{def:magma} with identity element~\ref{def:identity-element} $e$. An element $a\in X$ is a \emph{right inverse element} for $b\in X$ if $a * b=e$.
\end{definition}

\begin{definition}[inverse element]\label{def:inverse-element}
Let $(X,*)$ be a magma~\ref{def:magma} with identity element~\ref{def:identity-element} $e$. An element $a\in X$ is an \emph{inverse element} for $b\in X$ if it is both a right~\ref{def:right-inverse-element} and a left inverse element~\ref{def:left-inverse-element} for $b$.
\end{definition}

\begin{definition}[ring of subsets]\label{def:ring-of-subsets}
A nonempty family $R$ of sets is a \emph{ring} if it is closed under union and relative complement. That is, for all $A, B \in R$, 
- $A\cup B \in R$; 
- $A\setminus B \in R$.
\end{definition}

\begin{definition}[pre-measure]\label{def:pre-measure}
Let $R$ be a ring of subsets~\ref{def:ring-of-subsets} of $X$ and let $\mu_0:[0,\infty]$. Then $\mu_0$ is a \emph{pre-measure} if $\mu_0(\emptyset) = 0$ and for every sequence $\{A_n\}_{n\in\mathbb N}\subset R$ of pairwise disjoint subsets whose union is in $R$, $$\mu_0\left(\bigcup_{n\in\mathbb N}A_n\right) = \sum_{n\in\mathbb N} \mu_0(A_n).$$ That is, $\mu_0$ is σ-additive.
\end{definition}

\begin{definition}[row operations]\label{def:row-operations}
Let $A$ be a matrix~\ref{def:matrix} and let $A_i$ be the $i$th row of $A$. A \emph{row operation} on $A$ is any of the following operations:
1. Switch two rows
2. Multiply a row by a nonzero constant
3. Add one row to another
\end{definition}

\begin{definition}[elementary matrix]\label{def:elementary-matrix}
An \emph{elementary matrix} is the result of applying a row operation~\ref{def:row-operations} to the identity matrix~\ref{def:identity-matrix}.
\end{definition}

\begin{definition}[semigroup]\label{def:semigroup}
A \emph{semigroup} is a magma~\ref{def:magma} $(X,*)$ whose binary operation~\ref{def:binary-operation} $*$ is associative~\ref{def:associative}.
\end{definition}

\begin{definition}[monoid]\label{def:monoid}
A \emph{monoid} is a semigroup~\ref{def:semigroup} $(X,*)$ that has an identity element~\ref{def:identity-element} under $*$.
\end{definition}

\begin{definition}[group]\label{def:group}
A \emph{group} is a set $G$ together with a binary operation~\ref{def:binary-operation} $\cdot:G\times G \to G$ satisfying the following properties:
1. $\cdot$ is associative~\ref{def:associative};
2. there exists an identity element~\ref{def:identity-element} $e \in G$ such that $g\cdot e = g = e\cdot g$ for all $g \in G$.
3. for each $g \in G$ there exists an inverse element~\ref{def:inverse-element} $g^{-1}$ such that $g\cdot g^{-1} = e = g^{-1}\cdot g$.



Groups can be thought of as 
- associative~\ref{def:associative} magmas~\ref{def:magma}
- semigroups~\ref{def:semigroup} with an identity~\ref{def:identity-element} and with inverses~\ref{def:inverse-element}
- monoids~\ref{def:monoid} with inverses~\ref{def:inverse-element}
- categories~\ref{def:category} with a single object and every morphism an isomorphism
\end{definition}

\begin{definition}[abelian group]\label{def:abelian-group}
A group~\ref{def:group} $G$ is \emph{abelian} if its operation is commutative~\ref{def:commutative}.
\end{definition}

\begin{definition}[center of a group]\label{def:center-of-a-group}
The \emph{center} of a group~\ref{def:group} $G$ is $$Z(G) = \{z\in G\mid zg=gz \text{ for all } g\in G\}.$$
\end{definition}

\begin{definition}[centralizer of a subset]\label{def:centralizer-of-a-subset}
Let $G$ be a semigroup~\ref{def:semigroup} and let $S$ be a subset of $G$. The \emph{centralizer} of $S$ is $$C_G(S) = \{g\in G\mid gs=sg \text{ for all } s\in S\}.$$ When $G$ is a group~\ref{def:group}, we can also say $$C_G(S) = \{g\in G\mid gsg^{-1} =s \text{ for all } s\in S\}.$$
\end{definition}

\begin{definition}[circle group]\label{def:circle-group}
The \emph{circle group~\ref{def:group}}, usually denoted $\mathbb T$, is the complex $1$-torus. Equivalently, it is the multiplicative group~\ref{def:group} of complex numbers~\ref{def:complex-numbers} with complex absolute value $1$.
\end{definition}

\begin{definition}[divisible group]\label{def:divisible-group}
An abelian group~\ref{def:group} $(A,+)$ is \emph{divisible} if for every $d\in \mathbb N$ and for every $a\in A$, there exists $a'\in A$ such that $da' = a$.
\end{definition}

\begin{definition}[generating set of a group]\label{def:generating-set-of-a-group}
The \emph{generating set} of a group~\ref{def:group} $G$ is a subset $S$ of $G$ such that every element of $G$ can be expressed as a combination under the group operation~\ref{def:binary-operation} of finitely many elements of $S$ and their inverses.
\end{definition}

\begin{definition}[cyclic group]\label{def:cyclic-group}
A group~\ref{def:group} $G$ is \emph{cyclic} if it is generated~\ref{def:generating-set-of-a-group} by a single element.
\end{definition}

\begin{definition}[group word]\label{def:group-word}
Let $G$ be a group~\ref{def:group} and $S$ some subset of $G$. A \emph{group word in $S$} is an expression of the form $$s_1^{\ell_1}s_2^{\ell_2}\cdots s_n^{\ell_n}$$ where $s_i \in S$ and $\ell_i = \pm 1$.
\end{definition}

\begin{definition}[free group]\label{def:free-group}
A group~\ref{def:group} $G$ is \emph{free} if there exists a generating set~\ref{def:generating-set-of-a-group} $X$ of $G$ such that every nonempty reduced group word~\ref{def:group-word} in $X$ defines a nontrivial~\ref{def:identity-element} element of $G$.
\end{definition}

\begin{definition}[normalizer of a group]\label{def:normalizer-of-a-group}
Let $G$ be a semigroup~\ref{def:semigroup} and let $S$ be a subset of $G$. The \emph{normalizer} of $S$ is $$C_G(S) = \{g\in G\mid gS=Sg\}.$$ When $G$ is a group~\ref{def:group}, this may be written equivalently as $$C_G(S) = \{g\in G\mid gSg^{-1}\}.$$
\end{definition}

\begin{definition}[p-primary component of an abelian group]\label{def:p-primary-component-of-an-abelian-group}
Let $p$ be a prime and let $(A,+)$ be an abelian group~\ref{def:group}. The \emph{$p$-primary component} of $A$ is $$A_{(p)} = \{a\in A\mid \exists n\in\mathbb N \text{ st. }p^na= 0\}.$$
\end{definition}

\begin{definition}[primitive root of unity]\label{def:primitive-root-of-unity}
A root of unity $\zeta$ is \emph{primitive} if it generates~\ref{def:generating-set-of-a-group} the cyclic group~\ref{def:cyclic-group} of all $n$th roots of unity. Equivalently, it is \emph{primitive} if $\zeta^k\neq 1$ for any positive $k$ less than $n$.
\end{definition}

\begin{definition}[rank of a finitely generated abelian group]\label{def:rank-of-a-finitely-generated-abelian-group}
Let $G$ be a finitely generated~\ref{def:generating-set-of-a-group} abelian group~\ref{def:group}. By the  Fundamental Theorem of Finitely Generated Abelian Groups, $G$ may be uniquely written as a direct sum of cyclic groups~\ref{def:cyclic-group}. The \emph{rank} of $G$ is the number of $\mathbb Z$ summands in this expression.
\end{definition}

\begin{definition}[ring]\label{def:ring}
A (unital) \emph{ring} is a set $R$ together with two binary operations~\ref{def:binary-operation} $+$ and $\cdot$ such that
1. $R$ is an abelian group~\ref{def:group} under $+$;
2. $R$ is a monoid~\ref{def:monoid} under $\cdot$;
3. $\cdot$ distributes over $+$: $a\cdot (b+c)=(a\cdot b) + (a\cdot c)$ and $(b+c)\cdot a = (b\cdot a) + (c\cdot a)$.
\end{definition}

\begin{definition}[group ring]\label{def:group-ring}
Let $G$ be a group~\ref{def:group} and $R$ a ring~\ref{def:ring}. The \emph{group ring} $RG$ is the set of formal sums of elements in $G$ with coefficients in $R$. Addition is defined by components and multiplication is defined by $$(ag_i)(bg_j) = abg_k$$ where $g_ig_j = g_k$, and extended distributively.
\end{definition}

\begin{definition}[idempotent]\label{def:idempotent}
An element $x$ of a ring~\ref{def:ring} is \emph{idempotent} if $x^2 = x$.
\end{definition}

\begin{definition}[module over a ring]\label{def:module-over-a-ring}
Let $R$ be a ring~\ref{def:ring} with multiplicative identity~\ref{def:identity-element} $1$. A \emph{left $R$-module} $M$ is an abelian group~\ref{def:group} $(M, +)$ and another operation $\cdot : R\times M \to M$ such that fior all $r,s\in R$ and $x,y\in M$,
1. $r\cdot (x+y) = r\cdot x + r\cdot y$;
2. $(r+s)\dot x = r\cdot x + s\cdot x$;
3. $(rs)\cdot x = r\cdot (s\cdot x)$;
4. $1\cdot x = x$.

The operation $\cdot$ is called \emph{scalar multiplication}.
\end{definition}

\begin{definition}[direct sum of modules]\label{def:direct-sum-of-modules}
Let $R$ be a ring~\ref{def:ring} and $\{M_i\}_{i\in I}$ a family of left $R$-modules~\ref{def:module-over-a-ring} indexed by $I$. The \emph{direct sum} of the $M_i$, written $$\bigoplus_{i\in I} M_i,$$ is the set of all sequences $\{a_i\}_{i\in I}$ where $a_i \in M_i$ and $a_i=0$ for "cofinitely many" indices. That is, $a_i \neq 0$ for finitely many $i$.
\end{definition}

\begin{definition}[linearly independent]\label{def:linearly-independent}
Let $V$ be a module~\ref{def:module-over-a-ring} over the ring~\ref{def:ring} $R$. A subset $W \subset V$ is \emph{linearly independent} if for $v_i \in W$ and $\alpha_i \in R$, $$\sum_{i=1}^n \alpha_i v_i=0$$ only when $\alpha_i = 0$ for all $i$.
\end{definition}

\begin{definition}[nilpotent]\label{def:nilpotent}
An element $x$ of a ring~\ref{def:ring} $R$ is \emph{nilpotent} if there exists $n \in \mathbb Z$ such that $x^n =0$.
\end{definition}

\begin{definition}[polynomial ring]\label{def:polynomial-ring}
Let $R$ be a commutative~\ref{def:commutative} ring~\ref{def:ring} and $x$ a variable. Then $R[x]$ is the \emph{polynomial ring} of polynomials in $x$ with coefficients in $R$.

We can also define the \emph{polynomial ring} in multiple variables. Let $X = \{x_1,\dots,x_n\}$ be a finite collection of variables. Then $R[X] = R[x_1,\dots,x_n]$ is defined as follows:

To each function $I:X\to \mathbb N$, associate a monomial $$x^I = x_1^{I(x_1)}+\cdots + x_n^{I(x_n)}.$$ Then $$R[X]= \left\{ \sum_{I:X\to \mathbb N} a_IX^I \mid a_I\in R, a_I = 0 \text{ for all but finitely many } I \right\}.$$ Multiplication and addition are defined exactly how you would think.
\end{definition}

\begin{definition}[algebraic number]\label{def:algebraic-number}
An \emph{algebraic number} is a number that is a root of a nonzero polynomial~\ref{def:polynomial-ring} in one variable over $\mathbb Q$. It is often denoted by $\overline{\mathbb Q}$. The algebraic numbers form a field, but *not* a subfield (nor even a subset) of $\mathbb Q$.
\end{definition}

\begin{definition}[degree of polynomial]\label{def:degree-of-polynomial}
The \emph{degree} of a polynomial~\ref{def:polynomial-ring} is the highest power of the variable in its expression as a formal sum $a_nx^n + \cdots + a_0$ for $a_i \in R$.
\end{definition}

\begin{definition}[elementary symmetric polynomials]\label{def:elementary-symmetric-polynomials}
Let $R$ be a ring~\ref{def:ring} and let $R[x_1,\dots,x_n]$ be the polynomial ring~\ref{def:polynomial-ring} in $n$ variables. Then the \emph{elementary symmetric polynomials} are 
$$\begin{align*}e_1(x_1,\dots,x_n) &= \sum_{1\leq j\leq n} x_j \\ e_2(x_1,\dots,x_n) &= \sum_{1\leq j < k \leq n} x_jx_k \\e_3(x_1,\dots, x_n) &=\sum_{1\leq j < k < \ell\leq n} x_jx_kx_\ell\\&\quad\vdots\\e_n(x_1,\dots,x_n) &= x_1x_2\cdots x_n.\end{align*}$$ In general, $$e_k(x_1,\dots,x_n) = \sum_{1\leq j_1<j_2<\cdots < j_k \leq n} x_{j_1}\cdots x_{j_k}$$ for $1\leq k\leq n$. They are symmetric polynomials.
\end{definition}

\begin{definition}[power series]\label{def:power-series}
A \emph{power series} over a ring~\ref{def:ring} $R$ is a series of the form $$\sum_{n=0}^\infty a_n (x-c)^n$$ for $a_n, c\in R$.
\end{definition}

\begin{definition}[ring norm]\label{def:ring-norm}
A \emph{norm} $N$ on a ring~\ref{def:ring} $R$ is a function $N: R\to \mathbb N$ (here $\mathbb N$ is assumed to include $0$) such that $N(0) = 0$. Such a function is called \emph{positive} if it is nonzero on every element except for the additive identity~\ref{def:identity-element}.
\end{definition}

\begin{definition}[Euclidean domain]\label{def:euclidean-domain}
A ring~\ref{def:ring} $R$ is a \emph{Euclidean domain} if there exists a norm~\ref{def:ring-norm} $N$ on $R$ such that for all nonzero $a,b\in R$, there exist $q,r\in R$ such that $a = qb+r$ and either $r=0$ or $N(r) < N(b)$.
\end{definition}

\begin{definition}[Gaussian integers]\label{def:gaussian-integers}
The ring~\ref{def:ring} of \emph{Gaussian integers} is exactly $$\mathbb Z[i] = \{a+bi\mid a,b\in\mathbb Z, i^2=-1\}.$$
It is a Euclidean domain~\ref{def:euclidean-domain} with norm~\ref{def:ring-norm} $N(a+bi) = a^2+b^2 = (a+bi)(a-bi)$. This norm is multiplicative.
\end{definition}

\begin{definition}[rng]\label{def:rng}
A \emph{rng} is a set $R$ together with two binary operations~\ref{def:binary-operation} $+$ and $\cdot$ such that
1. $R$ is an abelian group~\ref{def:group} under $+$;
2. $R$ is a semigroup~\ref{def:semigroup} under $\cdot$;
3. $\cdot$ distributes over $+$: $a\cdot (b+c)=(a\cdot b) + (a\cdot c)$ and $(b+c)\cdot a = (b\cdot a) + (c\cdot a)$.

This is just a ring~\ref{def:ring} without a multiplicative identity~\ref{def:identity-element}.
\end{definition}

\begin{definition}[separable polynomial]\label{def:separable-polynomial}
A polynomial~\ref{def:polynomial-ring} is \emph{separable} if all of its roots are distinct.
\end{definition}

\begin{definition}[small category]\label{def:small-category}
A category~\ref{def:category} is \emph{small} if both its collection of objects and its collection of morphisms are sets, not proper classes
\end{definition}

\begin{definition}[strong differential basis]\label{def:strong-differential-basis}
The \emph{strong differential basis} is the differential basis consisting of rectangles containing $x$ that are parallel to all axes of $\mathbb R^n$.
\end{definition}

\begin{definition}[subgraph]\label{def:subgraph}
A \emph{subgraph} $A$ of a graph $X$ is a graph $A\subset X$ with $A^0\subset X^0$. That is, $A$ is a union of some of the vertices and edges of $X$.
\end{definition}

\begin{definition}[subgroup]\label{def:subgroup}
A subset $H\subset G$ of the group~\ref{def:group} $G$ is a \emph{subgroup} of $G$ if it is itself a group~\ref{def:group} under the operation~\ref{def:binary-operation} inherited from $G$.

The \emph{subgroup axioms} are
- closure under $*$;
- the identity element~\ref{def:identity-element} of $G$ is in $H$;
- for every $h\in H$, the inverse element~\ref{def:inverse-element} $h^{-1}$ is also in $H$.
\end{definition}

\begin{definition}[conjugate subgroups]\label{def:conjugate-subgroups}
Let $G$ be a group~\ref{def:group}. The subgroups~\ref{def:subgroup} $H_1$ and $H_2$ are \emph{conjugate} if there exists $a\in G$ such that $H_2= aHa^{-1}$.
\end{definition}

\begin{definition}[index of a subgroup]\label{def:index-of-a-subgroup}
Let $H$ be a subgroup~\ref{def:subgroup} of the group~\ref{def:group} $G$. The \emph{index} of $G$ is $[G:H] = \vert G\vert /\vert H\vert $. By Lagrange's theorem, this is always an integer.
\end{definition}

\begin{definition}[left ring ideal]\label{def:left-ring-ideal}
A \emph{left ideal} of a ring~\ref{def:ring} $R$ is a subset $I\subseteq R$ such that 
- $(I,+)$ is a subgroup~\ref{def:subgroup} of $(R,+)$;
- $I$ is closed under left multiplication in $R$ (i.e. for all $i\in R$ and all $r\in R$, $ri \in I$).
\end{definition}

\begin{definition}[maximal subgroup]\label{def:maximal-subgroup}
Let $G$ be a group~\ref{def:group}. A \emph{maximal} subgroup~\ref{def:subgroup} of $G$ is proper a subgroup~\ref{def:subgroup} of $G$ that is not contained in any other proper subgroup of $G$.
\end{definition}

\begin{definition}[Frattini subgroup]\label{def:frattini-subgroup}
Let $G$ be a group~\ref{def:group}. The \emph{Frattini subgroup~\ref{def:subgroup}} of $G$, denoted $\Phi(G)$, is the intersection of all of the maximal subgroups~\ref{def:maximal-subgroup} of $G$.
\end{definition}

\begin{definition}[right coset]\label{def:right-coset}
A \emph{right coset} of the subgroup~\ref{def:subgroup} $H$ of the group~\ref{def:group} $G$ for the element $x\in G$ is the set $Hx = \{hx\mid h \in H\}$.
\end{definition}

\begin{definition}[normal subgroup]\label{def:normal-subgroup}
A subgroup~\ref{def:subgroup} $H$ of the group~\ref{def:group} $G$ is \emph{normal} if for any $g \in H$, the left~\ref{def:left-coset} and right cosets~\ref{def:right-coset} of $H$ are equal. That is, if $gH = Hg$ for all $g \in G$. This is denoted $H\triangleleft G$.
\end{definition}

\begin{definition}[normal series]\label{def:normal-series}
Let $G$ be a group~\ref{def:group}. A \emph{normal series} is an ascending sequence of subgroups~\ref{def:subgroup} $$\{e\} = N_0\subseteq N_1\subseteq \cdots\subseteq N_{k-1} \subseteq N_k = G$$ such that $N_i$ is normal~\ref{def:normal-subgroup} in $G$ for all $0\leq i\leq k$.
\end{definition}

\begin{definition}[right ring ideal]\label{def:right-ring-ideal}
A \emph{right ideal}of a ring~\ref{def:ring} $R$ is a subset $I\subseteq R$ such that 
- $(I,+)$ is a subgroup~\ref{def:subgroup} of $(R,+)$;
- $I$ is closed under right multiplication in $R$ (i.e. for all $i\in R$ and all $r\in R$, $ir \in I$).
\end{definition}

\begin{definition}[ring ideal]\label{def:ring-ideal}
Let $R$ be a ring~\ref{def:ring}.  An \emph{ideal} $I$ of $R$ is a subset that is both a left ideal~\ref{def:left-ring-ideal} and a right ideal~\ref{def:right-ring-ideal}.
\end{definition}

\begin{definition}[generate a ring ideal]\label{def:generate-a-ring-ideal}
The ideal~\ref{def:ring-ideal} \emph{generated} by a subset $T$ of a ring~\ref{def:ring} $R$ is the smallest ideal~\ref{def:ring-ideal} containing $T$. 

Equivalently, $(T) = \bigcap\limits_{I\supseteq T} I$ where $I$ varies over all of the ideals~\ref{def:ring-ideal} in 
$R$ that contain $T$. 

Equivalently, if $T = \{t_i\}_{i=1}^n$, then $(T) = \left\{\sum\limits_{i=1}^n r_it_i\right\}$ , the set of all formal sums of elements of $T$ with coefficients in $R$.
\end{definition}

\begin{definition}[division in a ring]\label{def:division-in-a-ring}
Let $R$ be a commutative~\ref{def:commutative} ring~\ref{def:ring} and let $a,b \in R$. We say that $a$ \emph{divides} $b$ if there exists $r\in R$ such that $a = rb$. Equivalently, $a$ \emph{divides} $b$ if the ideal~\ref{def:ring-ideal} generated~\ref{def:generate-a-ring-ideal} by $b$ is contained in the ideal generated~\ref{def:generate-a-ring-ideal} by $a$.
\end{definition}

\begin{definition}[maximal ring ideal]\label{def:maximal-ring-ideal}
An ideal~\ref{def:ring-ideal} $I$ of the ring~\ref{def:ring} $R$ is \emph{maximal} if for any ideal~\ref{def:ring-ideal} $J$ such that $I\subseteq J$, either $J=I$ or $J=R$.
\end{definition}

\begin{definition}[prime ideal]\label{def:prime-ideal}
An ideal~\ref{def:ring-ideal} $\mathfrak p$ of a ring~\ref{def:ring} $R$ is \emph{prime} if it is proper proper and if $a,b\in R$ such that $ab \in I$, then $a\in \mathfrak p$ or $b\in \mathfrak p$.
\end{definition}

\begin{definition}[principal ideal]\label{def:principal-ideal}
An ideal~\ref{def:ring-ideal} $I$ of a ring~\ref{def:ring} $R$ is \emph{principal} if it can be generated~\ref{def:generate-a-ring-ideal} by a single element.
\end{definition}

\begin{definition}[greatest common divisor]\label{def:greatest-common-divisor}
Let $R$ be a commutative~\ref{def:commutative} ring~\ref{def:ring} and let $a,b\in R$. An element $d\in R$ is a \emph{greatest common divisor} of $a$ and $b$ if $d$ divides~\ref{def:division-in-a-ring} $a$ and $d$ divides~\ref{def:division-in-a-ring} $b$ and for every $d'\in R$ such that $d'$ divides~\ref{def:division-in-a-ring} $a$ and $d'$divides~\ref{def:division-in-a-ring} $b$, we have that $d'$ divides~\ref{def:division-in-a-ring} $d$. 

Equivalently, $d$ is a \emph{greatest common divisor} of $a$ and $b$ if the ideal~\ref{def:ring-ideal} generated~\ref{def:generate-a-ring-ideal} by $d$ is the smallest principal ideal~\ref{def:principal-ideal} that contains both $a$ and $b$. This ideal~\ref{def:ring-ideal} is unique.
\end{definition}

\begin{definition}[coprime]\label{def:coprime}
Two elements of a commutative~\ref{def:commutative} ring~\ref{def:ring} are \emph{coprime} if their  greatest common divisor~\ref{def:greatest-common-divisor} is $1$. Specifically, if it is the additive identity element~\ref{def:identity-element} of the ring~\ref{def:ring}.
\end{definition}

\begin{definition}[Euler's totient function]\label{def:euler's-totient-function}
\emph{Euler's totient function}, written $\phi(n)$, is a function $\mathbb N\to\mathbb N$ that gives the number of integers less than $n$ that are coprime~\ref{def:coprime} to $n$. That is, the number of integers less than $n$ whose greatest common divisor~\ref{def:greatest-common-divisor} with $n$ is 1.
\end{definition}

\begin{definition}[product of ring ideals]\label{def:product-of-ring-ideals}
Let $R$ be a commutative~\ref{def:commutative} ring~\ref{def:ring} with ideals~\ref{def:ring-ideal} $I$ and $J$. Then the \emph{product} of $I$ and $J$, defined by $$IJ= \{i_1j_1+\cdots + i_kj_j \mid i_\ell\in I,j_\ell\in J, \ell \in \mathbb N\}$$ is itself an ideal~\ref{def:ring-ideal}.
\end{definition}

\begin{definition}[proper ring ideal]\label{def:proper-ring-ideal}
A ring ideal~\ref{def:ring-ideal} is \emph{proper} if it is a proper subset of its parent ring~\ref{def:ring}.
\end{definition}

\begin{definition}[simple group]\label{def:simple-group}
A group~\ref{def:group} $G$ is \emph{simple} if it contains no normal subgroups~\ref{def:normal-subgroup} other than the subgroup~\ref{def:subgroup} containing the identity and itself.
\end{definition}

\begin{definition}[subinterval]\label{def:subinterval}
Let $P$ be a partition of a set~\ref{def:partition-of-a-set} of the interval $[a,b]$. Then arrange the elements of $P$ in increasing order: $$a = t_0 < t_1< \cdots< t_k = b.$$ The $[t_i, t_i+1]$ are the \emph{subintervals} of $[a,b]$ given by $P$.
\end{definition}

\begin{definition}[submodule]\label{def:submodule}
A \emph{submodule} $N$ of a module over a ring~\ref{def:module-over-a-ring} $M$ is a subset $N\subset M$ that forms a module over a ring~\ref{def:module-over-a-ring} under the same operations. It must form a group~\ref{def:group} under $+$.
\end{definition}

\begin{definition}[generate an R-module]\label{def:generate-an-r-module}
Let $R$ be a ring~\ref{def:ring}, $L$ an $R$-module~\ref{def:module-over-a-ring}, and $\{\ell_i\}_{i\in I}$ a subset of $L$. The submodule~\ref{def:submodule} \emph{generated} by the $\ell_i$ is the set $$\left\{\sum_{i=1}^n a_i\ell_i\mid a_i\in R, \text{ all but finitely many } a_i=0\right\}.$$ Equivalently, it is the smallest submodule~\ref{def:submodule} of $L$ containing the $\ell_i$.
\end{definition}

\begin{definition}[finitely generated R-module]\label{def:finitely-generated-r-module}
Let $R$ be a ring~\ref{def:ring}. An $R$-module~\ref{def:module-over-a-ring} $M$ is \emph{finitely generated} if it is generated~\ref{def:generate-an-r-module} by a finite subset of itself.
\end{definition}

\begin{definition}[free module]\label{def:free-module}
Let $R$ be a ring~\ref{def:ring}. An $R$-module~\ref{def:module-over-a-ring} $M$ is \emph{free} over a subset $S\subseteq M$ if $S$ is linearly independent~\ref{def:linearly-independent} and generates~\ref{def:generate-an-r-module} $M$. We say that $S$ forms a \emph{basis} for $M$.
\end{definition}

\begin{definition}[Noetherian module]\label{def:noetherian-module}
Let $R$ be a ring~\ref{def:ring} and let $M$ be an $R$-module~\ref{def:module-over-a-ring}. Then $M$ is \emph{Noetherian} (or \emph{satisfies the ascending chain condition on submodules}) if there are no infinite increasing chains of submodules~\ref{def:submodule}. That is, if whenever $M_1\subseteq M_2\subseteq M_3\subseteq \cdots$ is an increasing chain of submodules~\ref{def:submodule} of $M$, then there is a positive integer $m$ such that for all $k\geq m$, $M_k=M_m$.
\end{definition}

\begin{definition}[Noetherian ring]\label{def:noetherian-ring}
A ring~\ref{def:ring} $R$ is \emph{Noetherian} if it is Noetherian~\ref{def:noetherian-module} as a module~\ref{def:module-over-a-ring} over itself. That is, if there are no infinite increasing chains of left ideals~\ref{def:left-ring-ideal} of $R$.
\end{definition}

\begin{definition}[quotient of modules]\label{def:quotient-of-modules}
Let $R$ be a ring~\ref{def:ring}, let $M$ be an $R$-module~\ref{def:module-over-a-ring}, and let $N\subseteq N$ be a submodule~\ref{def:submodule} of $M$. Then the \emph{quotient module} $M/N$ is the set of cosets~\ref{def:left-coset} $\{m+N\mid m\in M\}$ with addition defined by $(m_1+N)+ (m_2+N) = (m_1+m_2) + N$ and scalar multiplication by $r(m_1 + N) = (rm_1)+N$.
\end{definition}

\begin{definition}[rank of a module]\label{def:rank-of-a-module}
Let $R$ be a ring~\ref{def:ring} and let $M$ be a free~\ref{def:free-module} $R$-module~\ref{def:module-over-a-ring}. Then the cardinality of a basis for $M$ is the \emph{rank} of $M$.

We can also define the \emph{rank} for modules~\ref{def:module-over-a-ring} that are not free~\ref{def:free-module}: it is the maximum number of linearly independent~\ref{def:linearly-independent} elements of $M$.
\end{definition}

\begin{definition}[simple module]\label{def:simple-module}
A module over a ring~\ref{def:module-over-a-ring} $M$ is \emph{simple} if it has no nonzero proper submodules~\ref{def:submodule}.
\end{definition}

\begin{definition}[composition series of a module]\label{def:composition-series-of-a-module}
Let $R$ be a ring~\ref{def:ring} and let $M$ be an $R$-module~\ref{def:module-over-a-ring}. A \emph{composition series} for $M$ is a chain of submodules~\ref{def:submodule} of $M$ $$M = F_0 \supset F_1 \supset \cdots \supset F_{n-1}\supset F_n =\{0\}$$ such that the quotient~\ref{def:quotient-of-modules} $F_i/F_{i+1}$ is simple~\ref{def:simple-module} for all $i$.
\end{definition}

\begin{definition}[finite length R-module]\label{def:finite-length-r-module}
Let $R$ be a ring~\ref{def:ring} and let $M$ be an $R$-module~\ref{def:module-over-a-ring}. We say that $M$ is of \emph{finite length} if it admits a composition series~\ref{def:composition-series-of-a-module} of finite length.
\end{definition}

\begin{definition}[semisimple module]\label{def:semisimple-module}
A module over a ring~\ref{def:module-over-a-ring} $M$ is \emph{semisimple} if it may be written as the direct sum~\ref{def:direct-sum-of-modules} of simple~\ref{def:simple-module} modules~\ref{def:module-over-a-ring}.
\end{definition}

\begin{definition}[subnormal series]\label{def:subnormal-series}
Let $G$ be a group~\ref{def:group}. A \emph{subnormal series} is an ascending sequence of subgroups~\ref{def:subgroup} $$\{e\} = N_0\subseteq N_1\subseteq \cdots\subseteq N_{k-1} \subseteq N_k = G$$ such that $N_i$ is normal~\ref{def:normal-subgroup} in $N_{i+1}$ for all $0\leq i\leq k$.
\end{definition}

\begin{definition}[subrectangle]\label{def:subrectangle}
Let $P$ be a partition of a set~\ref{def:partition-of-a-set} of the rectangle~\ref{def:rectangle} $Q$ in $\mathbb R^n$. The \emph{subrectangles} of $Q$ given by $P$ are the products of the subintervals~\ref{def:subinterval} of the component intervals~\ref{def:component-interval} given by each $P_i$ as in the definition of partition of a set~\ref{def:partition-of-a-set} of a rectangle~\ref{def:rectangle}.
\end{definition}

\begin{definition}[sum of modules]\label{def:sum-of-modules}
Let $R$ be a ring~\ref{def:ring} and let $M$ be an $R$-module~\ref{def:module-over-a-ring}. For $A$ and $B$ submodules~\ref{def:submodule} of $M$, define the \emph{sum} of $A$ and $B$ by $$A+B = \{a+b\mid a\in A,b\in B\}.$$ It is itself a submodule~\ref{def:submodule} of $M$, and in particular is the smallest one containing both $A$ and $B$.
\end{definition}

\begin{definition}[sum of ring ideals]\label{def:sum-of-ring-ideals}
Let $R$ be a ring~\ref{def:ring} with ideals~\ref{def:ring-ideal} $I$ and $J$. Then the \emph{sum} of $I$ and $J$, defined by $I+J= \{i+j \mid i\in I,j\in J\}$ is itself an ideal~\ref{def:ring-ideal}.
\end{definition}

\begin{definition}[comaximal ring ideals]\label{def:comaximal-ring-ideals}
Let $R$ be a ring~\ref{def:ring}. Two ideals~\ref{def:ring-ideal} $A$ and $B$ of $R$ are \emph{comaximal} if the sum~\ref{def:sum-of-ring-ideals} $A+B=R$.
\end{definition}

\begin{definition}[support]\label{def:support}
The \emph{support} of a function $f: X \to Y$ is the subset $S \subset X$ such that for all $x \in S$, $f(x) \neq 0$.
\end{definition}

\begin{definition}[surface]\label{def:surface}
A \emph{surface} is a two-dimensional manifold.
\end{definition}

\begin{definition}[surjective]\label{def:surjective}
A function $f:A\to B$ is \emph{surjective} if for all $b \in B$ there exists $a\in A$ such that $f(a)=b$.
\end{definition}

\begin{definition}[bijective]\label{def:bijective}
A function $f:A\to B$ is \emph{bijective} if it is both injective~\ref{def:injective} and surjective~\ref{def:surjective}.
\end{definition}

\begin{definition}[cardinal number]\label{def:cardinal-number}
Assuming the Axiom of Choice, the \emph{cardinality} of a set $X$ is the least ordinal number $\alpha$ such that there exists a bijection~\ref{def:bijective} between $X$ and $\alpha$. 

There exists a well-ordering on the cardinal numbers under the relation $\vert A\vert  \leq \vert B\vert $ if and only if there exists an injection~\ref{def:injective} from $A$ to $B$.
\end{definition}

\begin{definition}[countable]\label{def:countable}
A set $X$ is \emph{countable} if there exists a bijection~\ref{def:bijective} between $X$ and the set $\mathbb N$ of natural numbers. 

$\vert \mathbb N\vert = \aleph_0$ is the cardinal number~\ref{def:cardinal-number} of the naturals.
\end{definition}

\begin{definition}[group homomorphism]\label{def:group-homomorphism}
A \emph{group~\ref{def:group} homomorphism}is a function $f: G\to H$ between groups~\ref{def:group} $G$ and $H$ is a function that preserves the group operation~\ref{def:binary-operation}; that is $f(gh)=f(g)f(g)$ for all $g,h \in G$. 

A \emph{group isomorphism} is a homomorphism that is also bijective~\ref{def:bijective}.
\end{definition}

\begin{definition}[characteristic subgroup]\label{def:characteristic-subgroup}
A subgroup~\ref{def:subgroup} $H$ of a group~\ref{def:group} $G$ is \emph{characteristic in $G$} if every automorphism~\ref{def:group-homomorphism} of $G$ leaves $H$ invariant.
\end{definition}

\begin{definition}[inner automorphism of a group]\label{def:inner-automorphism-of-a-group}
Let $G$ be a group~\ref{def:group} and let $g\in G$. Define a homomorphism~\ref{def:group-homomorphism} $\varphi_g: G\to G$ by $\varphi_g(x) = gxg^{-1}$ for all $x\in G$. These $\varphi_g$ are in fact automorphisms~\ref{def:automorphism} and moreover $\varphi_g\circ\varphi_h = \varphi_{gh}$. The $\varphi_g$ are called \emph{inner automorphisms} of $G$.
\end{definition}

\begin{definition}[kernel of group homomorphism]\label{def:kernel-of-group-homomorphism}
Let $G$ and $H$ be groups~\ref{def:group}. The \emph{kernel} of a group homomorphism~\ref{def:group-homomorphism} $f:G\to H$ is the set of all elements $g\in G$ such that $f(g) = e_H$, i.e. the elements of $G$ that get mapped to the identity element~\ref{def:identity-element} in $H$.
\end{definition}

\begin{definition}[module homomorphism]\label{def:module-homomorphism}
Let $M$ and $N$ be modules~\ref{def:module-over-a-ring} over the ring~\ref{def:ring} $R$. A \emph{module homomorphism} between $M$ and $N$ is a function $f:M\to N$ such that for any $x,y \in M$ and $r\in R$, 
1. $f(x+y) = f(x) + f(y)$
2. $f(rx) = rf(x)$
where the operation in (2) is the scalar multiplication.

If a module homomorphism is bijective~\ref{def:bijective}, then it is an isomorphism of modules.
\end{definition}

\begin{definition}[direct sum of module homomorphisms]\label{def:direct-sum-of-module-homomorphisms}
Let $R$ be a ring~\ref{def:ring} and let $M$, $N$ and $L$ be $R$-modules~\ref{def:module-over-a-ring}. Then given homomorphisms~\ref{def:module-homomorphism} $f:M\to L$ and $g:N\to L$, there exists a unique homomorphism~\ref{def:module-homomorphism} $f\oplus g:M\oplus N\to L$ from the direct sum~\ref{def:direct-sum-of-modules} $M\oplus N$ to $L$ defined by $f\oplus g(m,n) = f(m) + g(n)$. 


https://www.wikidata.org/wiki/Q1142861](https://www.wikidata.org/wiki/Q13582243
https://www.wikidata.org/wiki/Q1142861)
\end{definition}

\begin{definition}[kernel of module homomorphism]\label{def:kernel-of-module-homomorphism}
Let $R$ be a ring~\ref{def:ring} and $M$ and $N$ $R$-modules~\ref{def:module-over-a-ring}. If $f:M\to N$ is a module homomorphism~\ref{def:module-homomorphism}, the \emph{kernel} of $f$ is the set $$\{m\in M\mid f(m) = e_N\}=\text{ker}(f)\subseteq M$$
\end{definition}

\begin{definition}[order of a group]\label{def:order-of-a-group}
The \emph{order} of a group~\ref{def:group} $G$ is the cardinal number~\ref{def:cardinal-number} of the underlying set of $G$.
\end{definition}

\begin{definition}[order of a group element]\label{def:order-of-a-group-element}
The \emph{order} of an element $g$ of the group~\ref{def:group} $G$ is the smallest integer $n \geq 1$ such that $g^n = e$, where $e$ is the group's~\ref{def:group} identity element~\ref{def:identity-element}

Equivalently, this is the cardinal number~\ref{def:cardinal-number} of the cyclic cyclic~\ref{def:cyclic-group} subgroup~\ref{def:subgroup} of $G$ generated~\ref{def:generating-set-of-a-group} by $g$ alone, if it exists.
\end{definition}

\begin{definition}[exponent of a finite group]\label{def:exponent-of-a-finite-group}
Let $A$ be a finite abelian group~\ref{def:group}. The \emph{exponent} of $A$ is the least common multiple of the orders~\ref{def:order-of-a-group-element} of all of the elements of $A$.
\end{definition}

\begin{definition}[p-group]\label{def:p-group}
Let $p$ be a prime number. A \emph{$p$-group~\ref{def:group}} is a group~\ref{def:group} whose order~\ref{def:order-of-a-group} is a power of $p$. 

Equivalently, a group is a $p$-group if the order~\ref{def:order-of-a-group-element} of every element in a $p$-group is a power of $p$.
\end{definition}

\begin{definition}[permutation]\label{def:permutation}
A \emph{permutation} is a bijection~\ref{def:bijective} from a set to itself.
\end{definition}

\begin{definition}[ring homomorphism]\label{def:ring-homomorphism}
A \emph{ring homomorphism} is a map $f:S\to R$ between rings~\ref{def:ring} $S$ and $R$ such that 
- $f:(S,+) \to (R,+)$ is a group homomorphism~\ref{def:group-homomorphism};
- $f(s_1s_2)= f(s_1)f(s_2)$ for all $s_1,s_2\in S$;
- $f(1) = 1$.

A \emph{ring isomorphism} is a homomorphism that happens to also be a bijection~\ref{def:bijective}.
\end{definition}

\begin{definition}[rule of product]\label{def:rule-of-product}
If there are $a$ ways of doing one thing and $b$ ways of doing another thing, then there are $ab$ ways of doing both things.  This can be extended inductively.

Set theoretically, this is the definition of the product of cardinal numbers: $$\vert S_1\vert \cdot\vert S_2\vert \cdots\vert S_n\vert  = \vert S_1\times S_2\times\cdots S_n\vert $$ where $\vert \cdot\vert $ is the cardinal number~\ref{def:cardinal-number} of the sets $S_i$ and $\times$ is the Cartesian product of sets. The sets need not be finite nor does the number of sets.
\end{definition}

\begin{definition}[rule of sum]\label{def:rule-of-sum}
If we have $a$ ways of doing one thing and $b$ ways of doing another thing but we cannot do both at the same time, then there are $a+b$ ways to choose to do one thing.

Set theoretically, the sum of the sizes of a finite collection of pairwise disjoint sets is the size~\ref{def:cardinal-number} of the union of the sets. That is,
$$\vert S_1\vert  + \vert S_2\vert  + \cdots + \vert S_n\vert  = \vert S_1\cup S_2\cup\cdots\cup S_n\vert .$$
\end{definition}

\begin{definition}[Sylow p-subgroup]\label{def:sylow-p-subgroup}
Let $p$ be a prime number. A \emph{Sylow $p$-subgroup~\ref{def:subgroup}} of a group~\ref{def:group} $G$ is a maximal p-subgroup~\ref{def:p-group} of $G$. That is, it is a p-group~\ref{def:p-group} that is a subgroup~\ref{def:subgroup} of $G$ and that is not a proper subgroup~\ref{def:subgroup} of any other p-subgroup of $G$.
\end{definition}

\begin{definition}[symmetric]\label{def:symmetric}
A binary relation~\ref{def:binary-relation} on the set $X$ is \emph{symmetric} if for all $a,b \in X$, $a$ is related to $b$ if and only if $b$ is related to $a$.
\end{definition}

\begin{definition}[symmetric difference]\label{def:symmetric-difference}
The \emph{symmetric difference} of two sets $A$ and $B$ is given by $$A\triangle B = (A\setminus B)\cup(B\setminus A) = (A\cup B)\setminus (A\cap B).$$ It is the set of elements that are in exactly one of $A$ or $B$.
\end{definition}

\begin{definition}[symmetric group]\label{def:symmetric-group}
Let $X$ be a set. The \emph{symmetric group} $S_X$ is the set of all bijections~\ref{def:bijective} from $X$ to itself.

When $X$ is finite with cardinal number~\ref{def:cardinal-number} $n$, it is denoted $S_n$ and consists of the set of permutations on $n$ letters.
\end{definition}

\begin{definition}[alternating group]\label{def:alternating-group}
The \emph{alternating group} $A_n$ is the subgroup~\ref{def:subgroup} of the symmetric group~\ref{def:symmetric-group} $S_n$ given by the kernel of the sign homomorphism. That is, it consists of all permutations whose sign is $1$.
\end{definition}

\begin{definition}[automorphism group]\label{def:automorphism-group}
The \emph{automorphism group} of a group~\ref{def:group} $G$ is the group~\ref{def:group} consisting of automorphisms~\ref{def:automorphism} of $G$. It is a subgroup~\ref{def:subgroup} of the symmetric group~\ref{def:symmetric-group} on the underlying set of $G$.
\end{definition}

\begin{definition}[group action]\label{def:group-action}
A \emph{group action} of a group~\ref{def:group} $G$ on a set $X$ is a map $G\times A\to A$ such that 
1. $g_1\cdot(g_2\cdot x)= (g_1g_2)\cdot x$ for all $g_1,g_2\in G$ and for all $x\in X$;
2. $e\cdot x = x$ for all $x\in X$.

Equivalently, we can describe a group action in terms of the homomorphism $\rho:G\to S_X$ (to the symmetric group~\ref{def:symmetric-group} on $X$) given by $\rho(g) = \sigma_g$ where $\sigma_g \in S_X$ is such that $x\mapsto g\cdot x$ under $\sigma_g$.
\end{definition}

\begin{definition}[equivariant map]\label{def:equivariant-map}
Let $G$ be a group~\ref{def:group} acting~\ref{def:group-action} on the sets $X$ and $Y$. A function $f:X\to Y$ is \emph{$G$-equivariant} if for all $x\in X$ and all $g\in G$, $$f(g\cdot x) = g\cdot f(x).$$ 
If the function $f$ is bijective~\ref{def:bijective}, then this function is called a \emph{$G$-isomorphism}.
\end{definition}

\begin{definition}[faithful group action]\label{def:faithful-group-action}
Let $G$ be a group~\ref{def:group} acting~\ref{def:group-action} on the set $X$. This action~\ref{def:group-action} is \emph{faithful} if the associated homomorphism~\ref{def:group-homomorphism} to the symmetric group~\ref{def:symmetric-group} $S_X$ has a trivial kernel~\ref{def:kernel-of-group-homomorphism}.
\end{definition}

\begin{definition}[G-invariant function]\label{def:g-invariant-function}
Let $G$ be a group~\ref{def:group} acting~\ref{def:group-action} on a the set $X$. A function $f:X\to Y$ is \emph{$G$-invariant} if for all $g\in G$ and all $x\in X$, $f(gx) = f(x)$.
\end{definition}

\begin{definition}[group action invariants]\label{def:group-action-invariants}
Let $G$ be a group~\ref{def:group} acting~\ref{def:group-action} on the set $X$. A point $x\in X$ is a \emph{fixed point} for or an \emph{invariant} of the $G$-action if $g\cdot x=x$ for every $g\in G$. The set of these points is denoted $X^G$.
\end{definition}

\begin{definition}[paradoxical set]\label{def:paradoxical-set}
Let $G$ be a group~\ref{def:group} acting~\ref{def:group-action} on a set $A$. Then $A$ is \emph{$G$-paradoxical} if there exist disjoint subsets $A_1,\dots, A_n, B_1,\dots, B_m\subset A$ and group elements $g_1,\dots,g_n,h_1,\dots,h_m\in G$ such that $$A=\bigcup_{i=1}^n g_i(A_i) \text{ and } A=\bigcup_{i=1}^m h_i(B_i).$$
\end{definition}

\begin{definition}[semidirect product]\label{def:semidirect-product}
Let $H$ and $N$ be groups~\ref{def:group} and let $\phi: H\to \text{Aut}(N)$ be a group homomorphism~\ref{def:group-homomorphism} from $H$ to the automorphism group~\ref{def:automorphism-group} of $N$. The \emph{semidirect product} $H\ltimes N$ is a group~\ref{def:group} on the Cartesian product $H\times N$ with operation $*$ given by $$(h_1,n_1)*(h_2,n_2) = (h_1h_2, n_1\phi(h_1)n_2).$$ The identity element~\ref{def:identity-element} of $H\ltimes N$ is the ordered pair consisting of the identity elements~\ref{def:identity-element} of $H$ and $N$, respectively. The inverse with respect to $*$ is given by $$(h,n)^{-1} = (h^{-1},\phi(h^{-1})n^{-1}).$$
\end{definition}

\begin{definition}[stabilizer under group action]\label{def:stabilizer-under-group-action}
Let $G$ be a group~\ref{def:group} acting~\ref{def:group-action} on the set $X$. The \emph{stabilizer} of a point $x\in X$, denoted $\text{Stab}(x)$, is the set of elements $g\in G$ such that $g\cdot x = x$.
\end{definition}

\begin{definition}[free group action]\label{def:free-group-action}
A group~\ref{def:group} $G$ acts~\ref{def:group-action} \emph{freely} on the set $X$ if it has trivial stabilizer under stabilizer~\ref{def:stabilizer-under-group-action}. That is, if $gx=x$ implies $g=e$ for all $x\in X$.
\end{definition}

\begin{definition}[torsion subgroup]\label{def:torsion-subgroup}
Let $A$ be an abelian group~\ref{def:group}. The \emph{torsion subgroup~\ref{def:subgroup}} $A_T$ of $A$ is the subgroup~\ref{def:subgroup} consisting of all elements of $A$ with finite order~\ref{def:order-of-a-group-element}.
\end{definition}

\begin{definition}[transitive]\label{def:transitive}
A binary relation~\ref{def:binary-relation} on a set $X$ is \emph{transitive} if for all $a,b,c\in X$, if $a$ is related to $b$ and $b$ is related to $c$, then $a$ is related to $c$.
\end{definition}

\begin{definition}[equivalence relation]\label{def:equivalence-relation}
An \emph{equivalence relation} $\sim$ on a set $X$ is a binary relation~\ref{def:binary-relation} that is reflexive~\ref{def:reflexive}, symmetric~\ref{def:symmetric}, and transitive~\ref{def:transitive}. That is, for all $a,b,c \in X$,
1. $a\sim a$.
2. $a \sim b$ if and only if $b \sim a$.
3. If $a \sim b$ and $b \sim c$, then $a\sim c$.
\end{definition}

\begin{definition}[equivalence class]\label{def:equivalence-class}
Let $\sim$ be an equivalence relation~\ref{def:equivalence-relation} on $X$. The \emph{equivalence class} of $a \in X$, denoted $[a]$, is the set $\{x \in X \mid a \sim X\}$.
\end{definition}

\begin{definition}[connected component]\label{def:connected-component}
The equivalence classes~\ref{def:equivalence-class} with respect to the connection equivalence relation~\ref{def:equivalence-relation} are called \emph{connected components} of $X$.
\end{definition}

\begin{definition}[partially ordered set]\label{def:partially-ordered-set}
A \emph{partially ordered set} is a set $X$ together with a homogeneous relation~\ref{def:homogeneous-relation} $\leq$ such that for all $x,y,z\in X$,
1. $\leq$ is reflexive~\ref{def:reflexive}, i.e. $x\leq x$;
2. $\leq$ is antisymmetric, i.e. if $x\leq y$ and $y\leq x$ then $x=y$;
3. $\leq$ is transitive~\ref{def:transitive}, i.e. if $x\leq y$ and $y\leq z$ then $x\leq z$.
\end{definition}

\begin{definition}[lower bound]\label{def:lower-bound}
A \emph{lower bound} of a subset $S$ of a poset~\ref{def:partially-ordered-set} $(X,\leq)$ is an element $a$ of $P$ such that for all $x\in S$, $a\leq x$.
\end{definition}

\begin{definition}[infimum]\label{def:infimum}
Let $(X,\leq)$ be a poset~\ref{def:partially-ordered-set} and let $S\subseteq X$. The \emph{infimum} of this subset, if it exists, is a lower bound~\ref{def:lower-bound} $a$ of $S$ in $X$ such that for all lower bounds~\ref{def:lower-bound} $y$ of $S$ in $X$, $y\leq a$. That is, the infimum is the greatest lower bound.
\end{definition}

\begin{definition}[quotient by equivalence relation]\label{def:quotient-by-equivalence-relation}
Let $\sim$ be an equivalence relation~\ref{def:equivalence-relation} on the set $S$. The \emph{quotient} of $S$ by $\sim$ (denoted $S/\sim$) is the set of equivalence classes~\ref{def:equivalence-class} of elements of $S$ with respect to $\sim$.
\end{definition}

\begin{definition}[quotient by normal subgroup]\label{def:quotient-by-normal-subgroup}
Let $G$ be a group~\ref{def:group} and $N$ a normal subgroup~\ref{def:normal-subgroup} of $G$. Define an equivalence relation~\ref{def:equivalence-relation} $\sim$ by setting $g\sim h$ if $g,h \in H$. Then the \emph{quotient of $G$ by $N$} (written $G/N$) is the quotient~\ref{def:quotient-by-equivalence-relation} by $\sim$ of $N$. todo
\end{definition}

\begin{definition}[composition series of a group]\label{def:composition-series-of-a-group}
A subnormal series~\ref{def:subnormal-series} for a group~\ref{def:group} $G$ is a \emph{composition series} when all of the quotients~\ref{def:quotient-by-normal-subgroup} (called \emph{composition factors}) $N_{i+1}/N_i$ are simple~\ref{def:simple-group}.
\end{definition}

\begin{definition}[quotient by subgroup]\label{def:quotient-by-subgroup}
Let $G$ be a group~\ref{def:group} with subgroup~\ref{def:subgroup} $H$. Then the \emph{quotient} $G/H$ is the  set of classes~\ref{def:equivalence-class} of cosets~\ref{def:left-coset} $gH$ for $g\in G$. The quotient is itself a subgroup~\ref{def:subgroup} exactly when $H$ is normal~\ref{def:normal-subgroup} in $G$. See the definition of the quotient~\ref{def:quotient-by-normal-subgroup} when the subgroup~\ref{def:subgroup} is normal.
\end{definition}

\begin{definition}[quotient ring]\label{def:quotient-ring}
Let $R$ be a ring~\ref{def:ring} and $I$ an ideal~\ref{def:ring-ideal} in $R$. Define an equivalence relation~\ref{def:equivalence-relation} $\sim$ on $R$ by $a\sim b$ when $a-b \in I$. The equivalence classes~\ref{def:equivalence-class} are given by $[a] = a+ I = \{a+r \mid r\in I\}$. The set of all such equivalence classes is denoted $R/I$, or the \emph{quotient of $R$ by $I$}.  The operations addition ($(a+I)+ (b+I) = (a+b)+I$) and multiplication ($(a+I)(b+I) = (ab)+I$) are well-defined. write_proof 

Moreover, this itself is a ring~\ref{def:ring} write_proof.

The quotient comes with a surjective~\ref{def:surjective} homomorphism~\ref{def:ring-homomorphism} $R\to R/I$ given by $r\mapsto r+I$.
\end{definition}

\begin{definition}[skeleton of a category]\label{def:skeleton-of-a-category}
The \emph{skeleton} $\text{sk}\mathcal C$ of a category~\ref{def:category} $\mathcal C$ is a subcategory with one object from each isomorphism class~\ref{def:equivalence-class} of objects of $\mathcal C$ such that the morphisms between two objects in $\text{sk}\mathcal C$ are all of the morphisms between these objects in $\mathcal C$.
\end{definition}

\begin{definition}[solvable group]\label{def:solvable-group}
A group~\ref{def:group} $G$ is \emph{solvable} if it admits a subnormal series~\ref{def:subnormal-series} whose quotients quotients~\ref{def:quotient-by-normal-subgroup} $N_{i+1}/N_i$ are all abelian.
\end{definition}

\begin{definition}[Stirling numbers of the second kind]\label{def:stirling-numbers-of-the-second-kind}
The \emph{Stirling numbers of the second kind} count the number of ways to partition a set of $n$ distinct objects into $k$ indistinct sets, leaving none empty. They are denoted $S(n,k)$.

They also count the number of different equivalence relations~\ref{def:equivalence-relation} with exactly $k$ equivalence classes~\ref{def:equivalence-class} that can be defined on a set of size $n$.
\end{definition}

\begin{definition}[supremum]\label{def:supremum}
Let $(X,\leq)$ be a poset~\ref{def:partially-ordered-set} and let $S\subseteq X$. The \emph{supremum} of this subset, if it exists, is an upper bound $b$ of $S$ in $X$such that for all upper bounds $x$ of $S$ in $X$, $b\leq x$. That is, the supremum is the least upper bound.
\end{definition}

\begin{definition}[lattice]\label{def:lattice}
A \emph{lattice} is a partially ordered set~\ref{def:partially-ordered-set} in which every pair of elements has a unique infimum~\ref{def:infimum} and a unique supremum~\ref{def:supremum}. 

If $(X, \leq)$ is a lattice, we call the supremum~\ref{def:supremum} of $x$ and $y$ the \emph{join} of $x$ and $y$ and denote it $x\lor y$. Similarly, we call the infimum~\ref{def:infimum} of $x$ and $y$ their \emph{meet} and denote it $x\land y$. These are in fact binary operations~\ref{def:binary-operation} on $X$. Moreover, they respect the partial order~\ref{def:partially-ordered-set} on $X$.  

We can also define a \emph{lattice} as a triple $(L, \lor,\land)$ where $L$ is a set and $\lor$ and $\land$ are two commutative~\ref{def:commutative}, associative~\ref{def:associative} binary operations~\ref{def:binary-operation} on $L$ satisfying the following \emph{absorption laws} for all $a,b \in L$:
1. $a\lor (a\land b) = a$;
2. $a\land (a\lor b) = a$.
\end{definition}

\begin{definition}[distributed lattice]\label{def:distributed-lattice}
A lattice~\ref{def:lattice} $(L,\lor,\land)$ is \emph{distributed} if it satisfies the following distributive laws for all $a,b,c\in L$:
- $a\lor (b\land c) = (a\lor b)\land (a\lor c)$;
- $a\land(b\lor c) = (a\land b)\lor (a\land c)$.
\end{definition}

\begin{definition}[lattice complement]\label{def:lattice-complement}
Let $(L,\leq)$ be a lattice~\ref{def:lattice}. A \emph{complement} of $a\in L$ is an element $b\in L$ such that the supremum~\ref{def:supremum} $\sup(a,b)$ is a maximum element of $L$ and the infimum~\ref{def:infimum} $\inf(a,b)$ is a minimum element of $L$. That is, if $y\leq \sup(a,b)$ and $\inf(a,b)\leq y$ for all $y\in L$.
\end{definition}

\begin{definition}[complemented lattice]\label{def:complemented-lattice}
A lattice~\ref{def:lattice} $(L,\leq)$ is said to be \emph{complemented} if it has a maximum and a minimum element and such that each element of $L$ admits a lattice complement~\ref{def:lattice-complement}.
\end{definition}

\begin{definition}[transitive group action]\label{def:transitive-group-action}
The action~\ref{def:group-action} of a group~\ref{def:group} $G$ on the set $X$ is \emph{transitive} if for every pair $x,y \in X$, there exists $g\in G$ such that $g\cdot x=y$.
\end{definition}

\begin{definition}[orbit]\label{def:orbit}
Let $G$ be a group~\ref{def:group} acting~\ref{def:group-action} on a set $X$. The \emph{orbit} of some $x \in X$  is the set of elements of $X$ to which $x$ can be moved by the elements of $G$. It is denoted $$G\cdot x = \{g\cdot x \mid g \in G\}.$$

The set of orbits of points of $X$ creates a partition of $X$ with associated equivalence relation~\ref{def:equivalence-relation} given by $x\sim y$ if and only if there exists $g \in G$ with $g\cdot x =y$, and the orbits are equivalence classes~\ref{def:equivalence-class}.

Equivalently, an \emph{orbit} of $x$ is the subset of $X$ containing $x$ on which the action of $G$ is transitive transitive~\ref{def:transitive-group-action}.
\end{definition}

\begin{definition}[conjugacy classes]\label{def:conjugacy-classes}
Let $G$ be a group~\ref{def:group}. Then $G$ acts~\ref{def:group-action} on itself by conjugation~\ref{def:inner-automorphism-of-a-group}: $$g\cdot x = gxg^{-1}.$$ The orbits~\ref{def:orbit} of this action are called \emph{conjugacy classes} of $G$.
\end{definition}

\begin{definition}[class function]\label{def:class-function}
A \emph{class function} on a group~\ref{def:group} $G$ is a function that is constant on the conjugacy classes~\ref{def:conjugacy-classes} of $G$.
\end{definition}

\begin{definition}[cycle]\label{def:cycle}
A \emph{cycle} is an element $\sigma$ of the symmetric group~\ref{def:symmetric-group} $S_X$ if the action~\ref{def:group-action} on $X$ of the subgroup~\ref{def:subgroup} generated~\ref{def:generating-set-of-a-group} by $\sigma$ has at most one orbit~\ref{def:orbit} with more than one element.
\end{definition}

\begin{definition}[cycle type]\label{def:cycle-type}
If $\sigma \in S_n$ is an element of the symmetric group~\ref{def:symmetric-group} on $n$ letters, then it is the product of disjoint cycles~\ref{def:cycle} of length $n_1\leq\cdots\leq n_r$ (including the $1$-cycles!) and moreover the sum of the $n_i$ is equal to $n$. The partition~\ref{def:partition-of-an-integer} of $n$ given by the $n_i$ is called the \emph{cycle type} of $\sigma$.
\end{definition}

\begin{definition}[transposition]\label{def:transposition}
The \emph{transposition} $\sigma_{ij}$ is a permutation~\ref{def:symmetric-group} in $S_n$ such that for some $i,j \in \{1,\dots ,n\}$, $$\sigma_{ij}(i) = j, \sigma_{ij}(j) = i, \text{ and } \sigma_{ij}(k) = k$$ for all $k\neq i,j$.

Equivalently, a \emph{transposition} is a 2-cycle~\ref{def:cycle}.
\end{definition}

\begin{definition}[tree]\label{def:tree}
A \emph{tree} is a connected graph~\ref{def:connected-graph} with no closed~\ref{def:closed-edge-path} reduced~\ref{def:reduced-edge-path} edge paths~\ref{def:edge-path}.
\end{definition}

\begin{definition}[maximal subtree]\label{def:maximal-subtree}
A subtree~\ref{def:subgraph} (a subgraph~\ref{def:subgraph} that happens to be a tree~\ref{def:tree}) of a graph $X$ is \emph{maximal} if it is contained in no strictly larger tree~\ref{def:tree}.
\end{definition}

\begin{definition}[uncountable]\label{def:uncountable}
A set is \emph{uncountable} if it is not countable~\ref{def:countable}; that is, if it has cardinal number~\ref{def:cardinal-number} greater than that of the natural numbers $\mathbb N$.
\end{definition}

\begin{definition}[unit of a ring]\label{def:unit-of-a-ring}
A \emph{unit} of a ring~\ref{def:ring} $R$ is an element $r\in R$ that has a multiplicative inverse~\ref{def:inverse-element}.
\end{definition}

\begin{definition}[division ring]\label{def:division-ring}
A \emph{division ring} is a ring~\ref{def:ring} $R$ in which every nonzero element is a unit~\ref{def:unit-of-a-ring}.
\end{definition}

\begin{definition}[field]\label{def:field}
A \emph{field} is a set $F$ together with two binary operations~\ref{def:binary-operation} addition ($+$) and multiplication ($\cdot$) such that
- Addition and multiplication are both associative~\ref{def:associative};
- Addition and multiplication are both commutative~\ref{def:commutative};
- There exist two distinct elements $0\neq 1$ in $F$ such that $0$ is an additive identity~\ref{def:identity-element} and $1$ is a multiplicative identity~\ref{def:identity-element};
- Every element of $F$ has an additive inverse~\ref{def:inverse-element};
- Every nonzero element of $F$ has a multiplicative inverse~\ref{def:inverse-element};
- Multiplication distributes over addition as follows: $a\cdot (b+c) = (a\cdot b) + (a\cdot c)$ for all $a,b,c\in F$. 

Equivalently, a \emph{field} is a commutative~\ref{def:commutative} division ring~\ref{def:division-ring}.
\end{definition}

\begin{definition}[affine transformations]\label{def:affine-transformations}
Let $F$ be a field~\ref{def:field}.  The \emph{affine transformations} of $F$$$\text{Aff}(F) = \{t_bm_a\mid b\in F, a\in F^\times\}$$ really combinations of scaling and translation: $x\mapsto ax+b$.
\end{definition}

\begin{definition}[algebraically closed]\label{def:algebraically-closed}
A field~\ref{def:field} $k$ is \emph{algebraically closed} if every non-constant polynomial~\ref{def:polynomial-ring} in $k[t]$ has a root in $k$.
\end{definition}

\begin{definition}[characteristic of a field]\label{def:characteristic-of-a-field}
Let $F$ be a field~\ref{def:field} with multiplicative identity~\ref{def:identity-element} $1_F$ and additive identity~\ref{def:identity-element} $0_F$. The \emph{characteristic} of $F$ is the smallest positive integer $p$ such that $p\cdot 1_F=0_F$ if such $p$ exists, and $0$ otherwise.
\end{definition}

\begin{definition}[discriminant]\label{def:discriminant}
The \emph{discriminant} of a collection $\{x_i\}_{i=1}^n$ of elements of a field~\ref{def:field} (or a ring~\ref{def:ring}?) to be $$D = \prod_{i<j} (x_i-x_j)^2.$$ The \emph{discriminant of a polynomial} is the discriminant of the roots of the polynomial.
\end{definition}

\begin{definition}[fixed field]\label{def:fixed-field}
Let $G$ be a group~\ref{def:group} of automorphisms~\ref{def:automorphism-group} of the field~\ref{def:field} $K$. Then the \emph{fixed field} of $K$ with respect to $G$ is $$K^G = \{\alpha\in K\mid \sigma(\alpha) = \alpha\forall \sigma\in G\}.$$
It is a field.
\end{definition}

\begin{definition}[fraction ring]\label{def:fraction-ring}
Let $R$ be a commutative~\ref{def:commutative} ring~\ref{def:ring} and let $S$ be a subset of $R$ that is closed under multiplication and contains the multiplicative identity~\ref{def:identity-element}. Then $$S^{-1}R = \left\{\left[\frac{r}{s}\right]\mid r\in R,s\in S\right\}$$  where equivalence of fractions is given by $\left[\frac{r_1}{s_1}\right] =\left[\frac{r_2}{s_2}\right]$ if and only if there exists $k \in S$ such that $kr_1s_2 = kr_2s_1$.
 
This is itself a ring~\ref{def:ring} under the usual operations we know and love for fractions of integers. 
 
If $S$ is exactly the nonzero elements of $R$, then this forms a field~\ref{def:field}.
 
 write_proof
\end{definition}

\begin{definition}[Frobenius endomorphism]\label{def:frobenius-endomorphism}
Let $F$ be a field~\ref{def:field} of characteristic~\ref{def:characteristic-of-a-field} $p$. Then the map $x\mapsto x^p$ is called the \emph{Frobenius endomorphism}.
\end{definition}

\begin{definition}[monic]\label{def:monic}
A \emph{monic polynomial} is a polynomial~\ref{def:polynomial-ring} whose leading coefficient is the unit~\ref{def:unit-of-a-ring} of the ring over which it is defined.
\end{definition}

\begin{definition}[algebraic integer]\label{def:algebraic-integer}
An \emph{algebraic integer} is a number that is the root of a monic~\ref{def:monic} polynomial~\ref{def:polynomial-ring} over $\mathbb Z$. The algebraic integers form a ring. Note that it is *not* a subring (nor even a subset) of $\mathbb Z$.
\end{definition}

\begin{definition}[perfect field]\label{def:perfect-field}
A field~\ref{def:field} $F$ is \emph{perfect} if it is of characteristic~\ref{def:characteristic-of-a-field} $p>0$ and every element in $p$ is a $p$th power or if it is of characteristic $0$.
\end{definition}

\begin{definition}[unital subring]\label{def:unital-subring}
A \emph{unital subring} of a ring~\ref{def:ring} $R$ is a subset $S$ of $R$ that is itself a ring~\ref{def:ring} under the binary operations~\ref{def:binary-operation} in $R$ with the same additive and multiplicative identities~\ref{def:identity-element}.
\end{definition}

\begin{definition}[unitary matrix]\label{def:unitary-matrix}
A complex $n\times n$ matrix~\ref{def:matrix} $U$ is \emph{unitary} if its conjugate transpose~\ref{def:conjugate-transpose} $U^*$ is equal to its inverse~\ref{def:inverse-matrix} $U^{-1}$.
\end{definition}

\begin{definition}[upper central series]\label{def:upper-central-series}
Let $G$ be a group~\ref{def:group}. The \emph{upper central series} for $G$ is an ascending sequence of subgroups~\ref{def:subgroup} of $G$ defined as follows: Let $Z_0(G) = \{e\}$ and let $Z_1(G)$ be the center~\ref{def:center-of-a-group} $Z(G)$. Then for each $i\geq 1$, define $Z_{i+1}(G)$ to be a subgroup~\ref{def:subgroup} of $G$ containing $Z_i(G)$ such that the quotient by normal quotient $Z_{i+1}(G)/Z_i(G) = Z(G/Z_i(G))$.  Note that because the center is a normal subgroup, these quotient by normal quotients are well-defined. Also note that this series may or may not reach $G$.
\end{definition}

\begin{definition}[nilpotent group]\label{def:nilpotent-group}
A group~\ref{def:group} $G$ is \emph{nilpotent} if its upper central series~\ref{def:upper-central-series} reaches $G$ after some finite number of steps. The number of steps is called the \emph{nilpotency class} of $G$.
\end{definition}

\begin{definition}[vector space]\label{def:vector-space}
A \emph{vector space} $V$ over the field~\ref{def:field} $F$ is a set together with two operations "addition" ($+: V\times V \to V$) and "scalar multiplication" ($\cdot: F\times V \to V$) satisfying the following eight axioms:
1. Addition is associative~\ref{def:associative}
2. Addition is commutative~\ref{def:commutative}
3. There exists a vector $0 \in V$ such that $0 + v = v$ for all $v \in V$. ($0$ is the additive identity in $V$)
4. For every $v \in V$, there exists $-v \in V$ such that $v + (-v) = 0$. ($-v$ is the additive inverse~\ref{def:inverse-element} of $v$)
5. $a(bv)= (ab)v$ for all $a,b \in F$ and $v \in V$.
6. $1v = v$ for all $v \in V$, where $1$ is the multiplicative identity in $F$.
7. $a(u+v) = au+av$ for all $a \in F$, $u,v \in V$.
8. $(a+b)v = av+bv$ for all $a,b \in F$, $v\in V$.

Axioms 7. and 8. are distributive laws.
\end{definition}

\begin{definition}[convex set]\label{def:convex-set}
Let $V$ be a vector space~\ref{def:vector-space}. A subset $K\subset V$ is \emph{convex} if for all $x,y \in K$ and $t \in [0,1]$, the vector $tx + (1-t)y$ is in $K$.
\end{definition}

\begin{definition}[geometrically independent]\label{def:geometrically-independent}
A collection $\{v_0,\dots,v_n\}$ of $n+1$ vectors~\ref{def:vector-space} is \emph{geometrically independent} if the $n$ vectors~\ref{def:vector-space} $v_1-v_0, v_2-v_0, \dots, v_n-v_0$ are linearly independent~\ref{def:linearly-independent}.
\end{definition}

\begin{definition}[linear transformation]\label{def:linear-transformation}
Let $W$ and $V$ be vector spaces~\ref{def:vector-space} over the same ground field $F$. A function $T: V \to W$ is a \emph{linear transformation} if for all $v,w \in V$ and $\alpha,\beta \in F$, $$T(\alpha v+\beta w) = \alpha T(v) + \beta T(w).$$

This is the correct notion for a homomorphism of vector spaces~\ref{def:vector-space}. If a linear transformation is bijective~\ref{def:bijective}, then it is an isomorphism of vector spaces~\ref{def:vector-space}.
\end{definition}

\begin{definition}[eigenvalue]\label{def:eigenvalue}
Let $T:V\to V$ be a linear transformation~\ref{def:linear-transformation} from the vector space~\ref{def:vector-space} $(V,F)$ to itself. Then $\lambda \in F$ is an \emph{eigenvalue} of $T$ if there exists $v\in V$ such that $T(v) = \lambda v$.
\end{definition}

\begin{definition}[eigenvector]\label{def:eigenvector}
Let $T:V\to V$ be a linear transformation~\ref{def:linear-transformation} from the vector space~\ref{def:vector-space} $(V,F)$ to itself. Then $v \in V$ is an \emph{eigenvector} of $T$ if there exists $\lambda\in F$ such that $T(v) = \lambda v$.
\end{definition}

\begin{definition}[kernel of linear transformation]\label{def:kernel-of-linear-transformation}
The \emph{kernel} of a linear transformation~\ref{def:linear-transformation} $T:V\to W$ is the set $$\text{ker}(T)= \{v\in V\mid T(v) =0\}.$$ It is a vector space~\ref{def:vector-space}.
\end{definition}

\begin{definition}[linearly dependent]\label{def:linearly-dependent}
Let $V$ be a vector space~\ref{def:vector-space}. A subset $W \subset V$ is \emph{linearly dependent} if it is not linearly independent~\ref{def:linearly-independent}.
\end{definition}

\begin{definition}[multilinear]\label{def:multilinear}
Let $V_1,\dots,V_n, W$ be vector spaces. A \emph{multilinear} map is a map $$f: V_1 \times\cdots\times V_n \to W$$ such that for all $1\leq i \leq n$ and $v_j \in V_j$ with $j \neq i$, the map $$f(v_1,\dots, v_{i-1}, -, v_{i+1}, \dots, v_n):V_i \to W$$ (which varies only in the $i$th component) given by $$V_i \ni v \mapsto f(v_1,\dots, v_{i-1}, v, v_{i+1}, \dots, v_n) \in W$$ is a linear transformation~\ref{def:linear-transformation}.
\end{definition}

\begin{definition}[algebra over a field]\label{def:algebra-over-a-field}
An \emph{algebra over a field~\ref{def:field}} $F$ is a vector space~\ref{def:vector-space} $U$ over $F$ with a bilinear~\ref{def:multilinear} binary operation~\ref{def:binary-operation} $U\times U \to U$ that is not necessarily associative~\ref{def:associative} but that follows left and right distributivity.

One can loosen this definition and only require $F$ to be a commutative~\ref{def:commutative} ring~\ref{def:ring}, but then you just get an algebra over a ring.
\end{definition}

\begin{definition}[algebra homomorphism]\label{def:algebra-homomorphism}
Let $A$ and $B$ be two algebras~\ref{def:algebra-over-a-field} over the same field~\ref{def:field} (or commutative~\ref{def:commutative} ring~\ref{def:ring}) $K$. Then an \emph{algebra homomorphism} between $A$ and $B$ is a map $f:A\to B$ such that 
1. $f(kx) = kf(x)$
2. $f(x+y) = f(x) + f(y)$
3. $f(xy) = f(x)f(y)$
for all $k\in K$ and $x,y\in A$.

If an algebra homomorphism is bijective~\ref{def:bijective}, then it is an isomorphism of algebras.
\end{definition}

\begin{definition}[associative algebra]\label{def:associative-algebra}
An \emph{associative algebra} is an algebra~\ref{def:algebra-over-a-field} $A$ whose binary operation~\ref{def:binary-operation} is associative~\ref{def:associative}.
\end{definition}

\begin{definition}[center of an algebra]\label{def:center-of-an-algebra}
The \emph{center} of an algebra~\ref{def:algebra-over-a-field} $A$ over the ring~\ref{def:ring} (or field~\ref{def:field}) $R$ is the set $Z(A) = \{z\in A \mid za=az \text{ for all } a\in A\}$.
\end{definition}

\begin{definition}[derivation]\label{def:derivation}
A \emph{derivation} of the algebra over a field~\ref{def:algebra-over-a-field} $U$ is a linear transformation~\ref{def:linear-transformation} $\delta$ satisfying the product rule: $$\delta(ab) = a\delta(b) + \delta(a) b.$$
\end{definition}

\begin{definition}[division algebra]\label{def:division-algebra}
A \emph{division algebra} $D$  is an algebra~\ref{def:algebra-over-a-field} over the field~\ref{def:field} $F$ such that for any element $a$ in $D$ and any non-zero element $b$ in $D$, there exists exactly one element $x\in D$ with $a=bx$ and exactly one element $y \in D$ such that $a=yb$.

If $D$ is an associative algebra~\ref{def:associative-algebra}, then we just require that $D$ have an identity element~\ref{def:identity-element} under the binary operation~\ref{def:binary-operation} and every nonzero element have a multiplicative inverse~\ref{def:inverse-element}.
\end{definition}

\begin{definition}[generate an associative algebra]\label{def:generate-an-associative-algebra}
Let $A$ be an algebra~\ref{def:algebra-over-a-field} over the field~\ref{def:field} $\mathbb k$. The elements $\{a_i\}_{i\in I}$ \emph{generate} $A$ if every element $\alpha\in A$ can be written (not necessarily uniquely) as a polynomial~\ref{def:polynomial-ring} in the $a_i$ with coefficients in $\mathbb k$.
\end{definition}

\begin{definition}[Hermitian inner product]\label{def:hermitian-inner-product}
A \emph{Hermitian inner product} $\langle\cdot,\cdot\rangle$ on a complex vector space~\ref{def:vector-space} $V$ is a bilinear~\ref{def:multilinear} form on $V$ that is antilinear in the second slot. That is,
1. $\langle u+v,w\rangle = \langle u,w\rangle +\langle v,w\rangle$;
2. $\langle u,v+w\rangle = \langle u,v\rangle +\langle u,w\rangle$;
3. $\langle \alpha u ,v\rangle = \alpha\langle u,v\rangle$ ;
4. $\langle u ,\alpha v\rangle = \overline\alpha\langle u,v\rangle$ ;
5. $\langle u,v\rangle = \overline{\langle v,u\rangle}$
6. $\langle u,u \rangle \geq 0$
7. $\langle u,u \rangle = 0$ if and only if $u=0$.
\end{definition}

\begin{definition}[inner product]\label{def:inner-product}
Let $(V,F)$ be a vector space~\ref{def:vector-space}. A function $\langle\cdot,\cdot\rangle: V\times V \to F$ is an \emph{inner product} on $V$ if 
1. it is bilinear~\ref{def:multilinear};
2. it is commutative~\ref{def:commutative};
3. $\langle v, v\rangle \geq 0$ for all $v \in V$ and $\langle v,v\rangle = 0$ if and only if $v=0$.
\end{definition}

\begin{definition}[adjoint of a linear transformation]\label{def:adjoint-of-a-linear-transformation}
Let $(V,F)$ and $(W,F)$ be vector spaces~\ref{def:vector-space} over the same field with respective inner products~\ref{def:inner-product} $\langle\cdot,\cdot\rangle_V$ and $\langle\cdot,\cdot\rangle_V$. The \emph{adjoint} of a linear transformation~\ref{def:linear-transformation} $T:V\to W$ is the linear transformation~\ref{def:linear-transformation} $T^*: W \to V$ such that $$\langle T(v), w\rangle_W = \langle v, T^*(w)\rangle_V$$ for all $v \in V$ and $w \in W$.
\end{definition}

\begin{definition}[Euclidean inner product]\label{def:euclidean-inner-product}
The \emph{Euclidean inner product~\ref{def:inner-product}} is an inner product~\ref{def:inner-product} on $\mathbb R^n$ given by $$\langle x,y\rangle = \sum_{i=1}^n x_iy_i.$$ The notation $\langle x,y\rangle$ notation is often written $x\cdot y$ instead.
\end{definition}

\begin{definition}[inner product space]\label{def:inner-product-space}
An \emph{inner product space} is a vector space~\ref{def:vector-space} $V$ together with an inner product~\ref{def:inner-product} $\langle\cdot,\cdot\rangle$ on $V$.
\end{definition}

\begin{definition}[Lie algebra]\label{def:lie-algebra}
Let $(L,F)$ be a vector space~\ref{def:binary-operation} with a binary operation~\ref{def:field} $L\times L \to L$ denoted $(x,y) \mapsto [xy]$. This operation is called the \emph{bracket} or \emph{commutator} of $x$ and $y$. The vector space together with the operation is a \emph{Lie algebra} over the field~\ref{def:multilinear} $F$ if 
1. The bracket operation is bilinear~\ref{def:vector-space};
2. $[xx] = 0$ for all $x \in L$;
3. $x[yz~\ref{def:binary-operation} + y[zx~\ref{def:field}+z[xy~\ref{def:multilinear} =0$ for all $x,y,z\in L$.

The last axiom is called the \emph{Jacobi identity}.
\end{definition}

\begin{definition}[abelian Lie algebra]\label{def:abelian-lie-algebra}
A Lie algebra~\ref{def:lie-algebra} $\mathcal g$ is \emph{abelian} if its bracket is always zero, i.e. if $[x,y] = 0$ for all $x,y\in \mathcal g$.
\end{definition}

\begin{definition}[centralizer of a Lie algebra]\label{def:centralizer-of-a-lie-algebra}
The \emph{centralizer} of a subset $X$ of a Lie algebra~\ref{def:lie-algebra} $L$ is $$C_L(X) = \{x\in L \mid [xX]=0\}.$$
\end{definition}

\begin{definition}[derived series of a Lie algebra]\label{def:derived-series-of-a-lie-algebra}
Let $L$ be a Lie algebra~\ref{def:lie-algebra}. Its \emph{derived series} is given by the derived algebras $$\begin{align*}L^{(0)} &= L\\ L^{(1)}&= [L,L] \\ L^{(2)} &= [L^{(1)},L^{(1)}]\\ &\text{ }\text{ }\vdots \\ L^{(i)}&=[L^{(i-1)},L^{(i-1)}] \\ &\text{ }\text{ }\vdots  \end{align*}$$
\end{definition}

\begin{definition}[direct sum of Lie algebras]\label{def:direct-sum-of-lie-algebras}
Let $\mathfrak g$ and $\mathfrak h$ be Lie algebras~\ref{def:lie-algebra}. The \emph{direct sum} $\mathfrak g \oplus \mathfrak h$ of $\mathfrak g$ and $\mathfrak h$ is the vector space direct sum of $\mathfrak g$ and $\mathfrak h$ when considered as vector spaces~\ref{def:vector-space} endowed with the bracket $$[(X_1,X_2),(Y_1,Y_2)]=([X_1,Y_1],[X_2,Y_2]).$$
\end{definition}

\begin{definition}[ideal of a Lie algebra]\label{def:ideal-of-a-lie-algebra}
A subspace $I$ of a Lie algebra~\ref{def:lie-algebra} $L$ is an \emph{ideal} of $L$ if $x\in L$, $y\in I$ implies that $[xy]\in I$. Equivalently by antisymmetry of the bracket, $[yx] \in I$.
\end{definition}

\begin{definition}[descending central series]\label{def:descending-central-series}
The \emph{descending central series} of ideals~\ref{def:ideal-of-a-lie-algebra} of a Lie algebra~\ref{def:lie-algebra} $L$ is given by the derived algebras $$\begin{align*}L^0 &= L\\ L^1 &= [L,L]\text{ } (= L^{(1)})\\ L^2 &= [L,L^1]\\ &\text{ }\text{ } \vdots\\ L^i &= [L,L^{i-1}]\\&\text{ }\text{ } \vdots\end{align*}$$ Where $L^{(1)}$ is the second element of the derived series~\ref{def:derived-series-of-a-lie-algebra} of $L$.
\end{definition}

\begin{definition}[Lie algebra homomorphism]\label{def:lie-algebra-homomorphism}
A linear transformation~\ref{def:linear-transformation} $\phi:L \to L'$ for Lie algebras~\ref{def:lie-algebra} $L$ and $L'$ is a \emph{homomorphism} if it preserves the bracket operation, that is if $$\phi([xy]_L) = [\phi(x)\phi(y)]_{L'}.$$

A Lie algebra~\ref{def:lie-algebra} homomorphism is an \emph{isomorphism} if it is bijective~\ref{def:bijective}.
\end{definition}

\begin{definition}[module over a Lie algebra]\label{def:module-over-a-lie-algebra}
Let $\mathfrak g$ be a Lie algebra~\ref{def:lie-algebra} over the field~\ref{def:field} $F$. A $\mathfrak g$-module over a ring~\ref{def:module-over-a-ring} is a pair $(V,\cdot)$ of a vector space~\ref{def:vector-space} $V$ over $F$ and $\cdot:\mathfrak g \times V \to V$ a map such that
1. $(\lambda x + \mu y)\cdot v = \lambda (x\cdot v) + \mu(y\cdot v)$
2. $x\cdot (\lambda v + \mu w) = \lambda (x\cdot v) + \mu(y\cdot v)$
3. $[x,y]\cdot v = x\cdot (y\cdot v) - y\cdot (x\cdot v)$

for all $x,y \in \mathfrak g$, $v,w\in V$, and $\lambda,\mu \in F$.
\end{definition}

\begin{definition}[nilpotent Lie algebra]\label{def:nilpotent-lie-algebra}
A Lie algebra~\ref{def:lie-algebra} $L$ is \emph{nilpotent} if some element $L^n$ of the descending central series~\ref{def:descending-central-series} or the ascending central series is equal to $\{0\}$.
\end{definition}

\begin{definition}[norm]\label{def:norm}
A \emph{norm} on a vector space~\ref{def:vector-space} $V$ is a real-valued functional $\vert \vert \cdot\vert \vert $ such that for all $v,w \in V$ and $\alpha \in \mathbb R$, the following properties hold:
1. $\vert \vert v+w\vert \vert  \leq \vert \vert v\vert \vert  + \vert \vert w\vert \vert $;
2. $\vert \vert \alpha v\vert \vert  = \vert \alpha\vert  \cdot \vert \vert v\vert \vert $;
3. $\vert \vert v\vert \vert  \geq 0$ and $\vert \vert v\vert \vert  = 0$ if and only if $v=0$.
\end{definition}

\begin{definition}[multiplicative norm]\label{def:multiplicative-norm}
A norm~\ref{def:norm} $\vert \cdot \vert  : A\to \mathbb R$ on an algebra~\ref{def:algebra-over-a-field} $A$ is called \emph{multiplicative} if $\vert 1_A\vert  = 1$ and for all $a,b\in A$,  $\vert ab\vert  \leq \vert a\vert \cdot \vert b\vert $.
\end{definition}

\begin{definition}[normed linear space]\label{def:normed-linear-space}
A \emph{normed linear space} is a vector space~\ref{def:vector-space} together with a norm~\ref{def:norm}.
\end{definition}

\begin{definition}[dual space]\label{def:dual-space}
The \emph{dual} of a normed linear space~\ref{def:normed-linear-space} $X$ is the set of bounded linear~\ref{def:linear-transformation} functions $f: X\to \mathbb R$. This set is itself a normed linear space~\ref{def:normed-linear-space} under the linear transformation norm.
\end{definition}

\begin{definition}[orthogonal]\label{def:orthogonal}
A set $\{v_i\}_{i\in \Lambda}$ in the vector space~\ref{def:vector-space} $V$ is \emph{orthogonal} if $\langle v_i, v_j\rangle =0$ for all $i\neq j$.
\end{definition}

\begin{definition}[orthonormal]\label{def:orthonormal}
A set $S$ of vectors in an inner product space~\ref{def:inner-product-space} is \emph{orthonormal} if every element of $S$ has norm~\ref{def:norm} $1$ and each vector in the set is orthogonal~\ref{def:orthogonal} to every other vector in the set.
\end{definition}

\begin{definition}[orthogonal matrix]\label{def:orthogonal-matrix}
An $n\times n$ real matrix~\ref{def:matrix} $A$ is \emph{orthogonal} if its column vectors (equivalently, its row vectors) are orthonormal~\ref{def:orthonormal}. That is, the product~\ref{def:matrix-multiplication} of $A$ with its transpose~\ref{def:matrix-transpose} $A^T$ is equal to the identity matrix~\ref{def:identity-matrix} $I_n$: $$AA^T = A^TA = I_n$$
Equivalently, the transpose~\ref{def:matrix-transpose} of an \emph{orthogonal} matrix~\ref{def:matrix} $A$ is equal to its inverse~\ref{def:inverse-matrix}: $A^T = A^{-1}$.
\end{definition}

\begin{definition}[positive definite]\label{def:positive-definite}
An $n\times n$ Hermitian complex matrix~\ref{def:matrix} $M$ is \emph{positive definite} if $\langle x, Mx\rangle > 0$ for all $x\neq 0$ in $\mathbb C^n$ where $\langle\cdot,\cdot\rangle$ is the standard inner product~\ref{def:inner-product} on $\mathbb C^n$

If we allow the dot product to be zero, then $M$ is \emph{postive semidefinite}.
\end{definition}

\begin{definition}[absolute value]\label{def:absolute-value}
Let $F$ be a field~\ref{def:field}. An \emph{absolute value} on $F$ is a map $\vert \cdot\vert :F\to \mathbb R$ such that 
1. $\vert x\vert  \geq 0$ for all $x\in F$;
2. $\vert x\vert  = 0$ if and only if $x=0$;
3. $\vert x+y\vert  \leq \vert x\vert  + \vert y\vert$ for all $x,y\in F$;
4. $\vert x\cdot y\vert  = \vert x\vert \cdot \vert y\vert$ for all $x,y\in F$.

Together, $1$ and $2$ make up the condition that $\vert \cdot\vert$ be positive definite~\ref{def:positive-definite} while $3$ is just the triangle inequality~\ref{def:norm}.
\end{definition}

\begin{definition}[equivalent absolute values]\label{def:equivalent-absolute-values}
Two absolute values~\ref{def:absolute-value} over the same field~\ref{def:field} $F$, call them $\vert \cdot\vert _1$ and $\vert \cdot\vert _2$, are \emph{equivalent} if there exists a real number $\alpha>0$ such that $\vert x\vert _1^\alpha = \vert x\vert _2$ for all $x\in F$.
\end{definition}

\begin{definition}[non-Archimedean absolute value]\label{def:non-archimedean-absolute-value}
A \emph{non-Archimedean absolute value} on a field~\ref{def:field} $F$ is an absolute value~\ref{def:absolute-value} $\vert \cdot\vert :F\to \mathbb R$ such that $\vert x+y\vert \leq\max\{\vert x\vert ,\vert y\vert \}$ for all $x,y\in F$. This inequality is called the \emph{ultrametric inequality} and it is stronger than the triangle inequality~\ref{def:norm}.
\end{definition}

\begin{definition}[self-adjoint]\label{def:self-adjoint}
A linear transformation~\ref{def:linear-transformation} $T:V\to V$ is \emph{self-adjoint} if $T^*=T$.
\end{definition}

\begin{definition}[simple Lie algebra]\label{def:simple-lie-algebra}
If a Lie algebra~\ref{def:lie-algebra} has no ideals~\ref{def:ideal-of-a-lie-algebra} except for $\{0\}$ and itself, then it is \emph{simple}.
\end{definition}

\begin{definition}[solvable Lie algebra]\label{def:solvable-lie-algebra}
A Lie algebra~\ref{def:lie-algebra} $L$ is \emph{solvable} if some element $L^{(n)}$ of its derived series of a Lie algebra~\ref{def:derived-series-of-a-lie-algebra} is equal to 0.
\end{definition}

\begin{definition}[radical of a Lie algebra]\label{def:radical-of-a-lie-algebra}
The \emph{radical} of a Lie algebra~\ref{def:lie-algebra} $L$ is the maximal solvable~\ref{def:solvable-lie-algebra} ideal~\ref{def:ideal-of-a-lie-algebra} of $L$.
\end{definition}

\begin{definition}[semisimple Lie algebra]\label{def:semisimple-lie-algebra}
A Lie algebra~\ref{def:lie-algebra} $L$ is \emph{semisimple} if its radical~\ref{def:radical-of-a-lie-algebra} $\text{Rad}(L) = \{0\}$.
\end{definition}

\begin{definition}[span]\label{def:span}
Let $V$ be a vector space~\ref{def:vector-space}. The set $W \subseteq V$ \emph{spans} $V$ if for all $v \in V$, there exists $n \in \mathbb N$, a collection $\{\alpha_i\}_{i=1}^n \subset \mathbb R$, and a collection $\{v_i\}_{i=1}^n \subset W$ such that $$v = \sum_{i=1}^n \alpha_i v_i.$$
\end{definition}

\begin{definition}[basis]\label{def:basis}
A set $B$ is a \emph{basis} for the vector space~\ref{def:vector-space} $V$ if it consists of linearly independent~\ref{def:linearly-independent} vectors that span~\ref{def:span} $V$.
\end{definition}

\begin{definition}[complexification of a vector space]\label{def:complexification-of-a-vector-space}
Let $V$ be a vector space~\ref{def:vector-space} over the field~\ref{def:field} $\mathbb R$ of real numbers, and denote its basis~\ref{def:basis} by $\{v_k\}_{k=1}^n$. The \emph{complexification} of $V$ is the vector space~\ref{def:vector-space} $V_\mathbb C$ over the field~\ref{def:field} $\mathbb C$ of complex numbers with basis $\{v_k\}_{k=1}^n$.
\end{definition}

\begin{definition}[dimension of vector space]\label{def:dimension-of-vector-space}
Let $V$ be a vector space~\ref{def:vector-space}. The \emph{dimension} of $V$ is the cardinal number~\ref{def:cardinal-number} of any basis~\ref{def:basis} for $V$.
\end{definition}

\begin{definition}[determinant]\label{def:determinant}
The \emph{determinant} of a linear transformation~\ref{def:linear-transformation} $T: V\to V$ (where $V$ is finite-dimensional~\ref{def:dimension-of-vector-space}) is the unique antisymmetric $n$-variable multilinear~\ref{def:multilinear} map $\text{det}$ such that $$\text{det}:(\mathbb R^n)^n \to \mathbb R$$ such that $\text{det}(e_1,\dots,e_n)=1$. It is defined for a the matrix $A$ of $T$ as $$\text{det}(A) = \text{det}(A_1,\dots,A_n)$$ where the $A_i$ are the columns of $A$.
\end{definition}

\begin{definition}[characteristic polynomial]\label{def:characteristic-polynomial}
The \emph{characteristic polynomial} of a linear transformation~\ref{def:linear-transformation} $T$ is the polynomial in $\lambda$ given by the determinant~\ref{def:determinant} of the linear transformation $T-\lambda I$: $$p_T(\lambda) = \text{det}(T-\lambda I).$$
\end{definition}

\begin{definition}[field extension]\label{def:field-extension}
A field~\ref{def:field} $E$ is a \emph{field extension} of the field~\ref{def:field} $F$ if $F$ is a subfield of $E$. 

The field extension may be considered as a vector space~\ref{def:vector-space} over $F$. The dimension~\ref{def:dimension-of-vector-space} of $E$ as an $F$-vector space~\ref{def:vector-space} is denoted $[E:F]$, $\deg(E/F)$, or $\deg_F(E)$.

The most basic field extension comes from adjoining an element $\Theta$ to $F$. It is denoted $F(\Theta)$, and it is the smallest field~\ref{def:field} containing both $F$ and $\Theta$.
\end{definition}

\begin{definition}[algebraic extension]\label{def:algebraic-extension}
A field extension~\ref{def:field-extension} $F/K$ is \emph{algebraic} if every element of $K$ is algebraic over $F$.
\end{definition}

\begin{definition}[composite field extension]\label{def:composite-field-extension}
Let $F\subset L\subset E$ and $F\subset K\subset E$ be two field extensions~\ref{def:field-extension}. Then the \emph{composite} field extension~\ref{def:field-extension} $F\subset KL$ is the smallest subfield of $E$ containing both fields~\ref{def:field} $K$ and $L$.
\end{definition}

\begin{definition}[F-conjugate]\label{def:f-conjugate}
Let $E_i$ be an algebraic extension~\ref{def:algebraic-extension} over the field~\ref{def:field} $F$ for each $i\in \mathbb N$. Then $\alpha_k$ is an \emph{$F$-conjugate} of $\alpha_\ell$ if the minimal polynomial of $\alpha_k$ over $F$ is the same as the minimal polynomial of $\alpha_\ell$ over $F$.
\end{definition}

\begin{definition}[field embedding]\label{def:field-embedding}
Let $F$ be a field and let $E/F$ And $K/F$ be two field extensions~\ref{def:field}. An \emph{$F$-embedding} $\phi:E\to K$ is a ring homomorphism~\ref{def:field-extension} such that $\phi(c) = c$ for all $c\in F$.
\end{definition}

\begin{definition}[functional]\label{def:functional}
A linear~\ref{def:linear-transformation} \emph{functional} on a vector space~\ref{def:vector-space} $V$ with underlying field~\ref{def:field} $F$ is a linear transformation~\ref{def:linear-transformation} from $V$ to $F$ when considered as a one-dimensional~\ref{def:dimension-of-vector-space} vector space~\ref{def:vector-space} over itself.
\end{definition}

\begin{definition}[general linear algebra]\label{def:general-linear-algebra}
Let $(V,F)$ be a vector space~\ref{def:vector-space} and let $\text{End}(V)$ be the set of linear transformations~\ref{def:linear-transformation} $V\to V$. This is a ring~\ref{def:ring} with resepect to the usual product operation. Define a new \emph{bracket} operation by $[x,y] = xy-yx$. Then $\text{End}(V)$ is a Lie algebra~\ref{def:lie-algebra} over $F$. We call it the \emph{general linear algebra} and denote it $\mathfrak{gl}(V)$.

If $V$ is of finite dimension~\ref{def:dimension-of-vector-space}, then $\text{dim}(\text{End}(V)) = \text{dim}(V)^2$.

The general linear algebra is a Lie algebra~\ref{def:lie-algebra}.
\end{definition}

\begin{definition}[general linear group]\label{def:general-linear-group}
The \emph{general linear group~\ref{def:group}} $GL(n)$ or $GL(V)$ is the group~\ref{def:group} of automorphisms~\ref{def:automorphism} of an $n$-dimensional~\ref{def:dimension-of-vector-space} vector space~\ref{def:vector-space} $V$. Equivalently, it is the group~\ref{def:group} of invertible~\ref{def:bijective} $n\times n$ matrices.
\end{definition}

\begin{definition}[G-harmonic polynomial]\label{def:g-harmonic-polynomial}
Let $G$ be a subgroup~\ref{def:subgroup} of the general linear group~\ref{def:general-linear-group} $GL_n(\mathbb C)$ acting~\ref{def:group-action} naturally on $\mathbb C^n$. A polynomial~\ref{def:polynomial-ring} $f\in \mathbb C[x_1,\dots,x_n]$ is \emph{$G$-harmonic} if it is annihilated by all $G$-invariant~\ref{def:g-invariant-function} differential operators without constant term. That is, $f$ is $G$-harmonic if $u(f) = 0$ for all $u \in D^G_{>0}$.
\end{definition}

\begin{definition}[generate a field]\label{def:generate-a-field}
Let $F\subset E$ be a field extension~\ref{def:field-extension}. Then $F(e_1,\dots, e_n)$ is the smallest subfield of $E$ containing $F$ and all of the $e_i$. We say that the $e_i$ \emph{generate} this field.
\end{definition}

\begin{definition}[cyclotomic field]\label{def:cyclotomic-field}
The \emph{cyclotomic field of $n$th roots of unity} is generated~\ref{def:generate-a-field} over $\mathbb Q$ by a primitive~\ref{def:primitive-root-of-unity} $n$th root of unity: $\mathbb Q(\zeta_n)$.
\end{definition}

\begin{definition}[group representation]\label{def:group-representation}
Let $V$ be an $n$-dimensional~\ref{def:dimension-of-vector-space} vector space~\ref{def:vector-space}. A \emph{representation} of a group~\ref{def:group} $G$ is a group homomorphism~\ref{def:group-homomorphism} $\rho: G \to GL(V)$ from $G$ to the $n$-dimensional~\ref{def:dimension-of-vector-space} general linear group~\ref{def:general-linear-group} of invertible linear transformations~\ref{def:linear-transformation} of $V$.
\end{definition}

\begin{definition}[dimension of group representation]\label{def:dimension-of-group-representation}
Let $G$ be a group~\ref{def:group} and $V$ a vector space~\ref{def:vector-space}. The \emph{dimension} of a group representation~\ref{def:group-representation} $\rho:G\to \text{GL}(V)$ is just the usual dimension~\ref{def:dimension-of-vector-space} of $V$.
\end{definition}

\begin{definition}[character of a representation]\label{def:character-of-a-representation}
Let $(V,\rho)$ be a finite-dimensional~\ref{def:dimension-of-group-representation} representation~\ref{def:group-representation} of a group~\ref{def:group} $G$ where $V$ is a vector space~\ref{def:vector-space} over the field~\ref{def:field} $F$. The \emph{character} of this representation is the function $\chi_\rho:G\to F$ given by $\chi_\rho(g) = \text{Tr}(\rho(g))$.
\end{definition}

\begin{definition}[direct sum of group representations]\label{def:direct-sum-of-group-representations}
Let $V$ and $W$ be vector spaces~\ref{def:vector-space}. If $\phi: G\to GL(V)$ and $\psi:G\to GL(W)$ are two different representations~\ref{def:group-representation} of the group~\ref{def:group} $G$, the \emph{direct sum} $\phi \oplus \psi:G \to GL(V\oplus W)$ is also a representation~\ref{def:group-representation} of $G$ where $V\oplus W$ is given by the usual Cartesian product of vector spaces.
\end{definition}

\begin{definition}[G-invariant vectors]\label{def:g-invariant-vectors}
Let $(V,\rho)$ be a representation~\ref{def:group-representation} of the group~\ref{def:group} $G$. The \emph{$G$-invariant elements} of $V$ is $$V^G = \{v\in V\mid \rho(g)v=v\},$$ the set of vectors in $V$ left unchanged by the action of $G$.
\end{definition}

\begin{definition}[Lie algebra representation]\label{def:lie-algebra-representation}
A \emph{representation} $\phi$ of a Lie algebra~\ref{def:lie-algebra} $L$ is a homomorphism~\ref{def:lie-algebra-homomorphism} to the general linear algebra~\ref{def:general-linear-algebra} ($\phi:L \to \mathfrak{gl}(V)$) where $V$ is a vector space~\ref{def:vector-space}. Here we require the dimension~\ref{def:dimension-of-vector-space} of the vector space~\ref{def:vector-space} underlying $L$ to be finite, but admit arbitrary dimension~\ref{def:dimension-of-vector-space} for $V$.
\end{definition}

\begin{definition}[adjoint representation of a Lie algebra]\label{def:adjoint-representation-of-a-lie-algebra}
Let $\mathfrak g$ be a Lie algebra~\ref{def:lie-algebra}. Then the adjoint map $\text{ad}:\mathfrak g \to \mathfrak{gl}(\mathfrak g)$ is given by $$\text{ad}_X(Y) = [X,Y].$$ Because the adjoint map is a Lie algebra homomorphism, it is a representation~\ref{def:lie-algebra-representation} of $\mathfrak g$, and therefore we may refer to $\text{ad}$ as the \emph{adjoint representation} of $\mathfrak g$.
\end{definition}

\begin{definition}[dual representation of a Lie algebra]\label{def:dual-representation-of-a-lie-algebra}
Let $\mathfrak g$ be a Lie algebra~\ref{def:lie-algebra} and let $(V,\rho_V)$ be a representation~\ref{def:lie-algebra-representation} of $\mathfrak g$. Then the \emph{dual representation} of $\mathfrak g$ is the pair $(V^*,\rho_{V^*})$ where $V^*$ is the dual space~\ref{def:dual-space} of $V$ and $\rho_{V^*}$ is given by $$\rho_{V^*}(x) = -\rho_V(x)^
*$$ for all $x \in \mathfrak g$.
\end{definition}

\begin{definition}[inner automorphism of a Lie algebra]\label{def:inner-automorphism-of-a-lie-algebra}
An automorphism~\ref{def:automorphism} of a Lie algebra~\ref{def:lie-algebra} $L$ of the form $\text{exp}(\text{ad} (x))$ where $\text{ad}$ is the adjoint representation~\ref{def:adjoint-representation-of-a-lie-algebra} is called an \emph{inner automorphism} of $L$.
\end{definition}

\begin{definition}[morphism of group representations]\label{def:morphism-of-group-representations}
Let $V$ and $W$ be vector spaces~\ref{def:vector-space}. An \emph{intertwiner} between two representations~\ref{def:group-representation} $\phi:G\to GL(V)$ and $\psi:G\to GL(W)$ of the group~\ref{def:group} $G$ is a linear transformation~\ref{def:linear-transformation} $T:V\to W$ such that $$T\circ(\phi(g)) = (\psi(g))\circ T$$ for all $g \in G$.
\end{definition}

\begin{definition}[equivalent group representations]\label{def:equivalent-group-representations}
Two representations~\ref{def:group-representation} of a group~\ref{def:group} $G$ are \emph{equivalent} if there exists an invertible~\ref{def:bijective} intertwiner~\ref{def:morphism-of-group-representations} between them.
\end{definition}

\begin{definition}[isomorphism of group representations]\label{def:isomorphism-of-group-representations}
An \emph{isomorphism of group representations~\ref{def:group-representation}} is a morphism of group representations~\ref{def:morphism-of-group-representations} that happens to be bijective~\ref{def:bijective}.
\end{definition}

\begin{definition}[multiplicity of simple module in semisimple module]\label{def:multiplicity-of-simple-module-in-semisimple-module}
Let $A$ be an algebra~\ref{def:algebra-over-a-field} over an algebraically closed~\ref{def:algebraically-closed} field~\ref{def:field} $\mathbb k$ and let $M$ be a semisimple~\ref{def:semisimple-module} module~\ref{def:module-over-a-ring} over $A$ of finite dimension~\ref{def:dimension-of-vector-space} over $\mathbb k$. The number of times the module~\ref{def:module-over-a-ring} $V$ appears in a direct sum decomposition of $M$ is denoted $[M:V]$ and is called the \emph{multiplicity of $V$ in $M$}.
\end{definition}

\begin{definition}[nullity]\label{def:nullity}
The \emph{nullity} of a linear transformation~\ref{def:linear-transformation} $T$ is the dimension~\ref{def:dimension-of-vector-space} of the kernel~\ref{def:kernel-of-linear-transformation} of $T$.
\end{definition}

\begin{definition}[outer derivation]\label{def:outer-derivation}
An \emph{outer derivation} is any derivation~\ref{def:derivation} that is not an inner derivation~\ref{def:adjoint-representation-of-a-lie-algebra}.
\end{definition}

\begin{definition}[prime subfield]\label{def:prime-subfield}
The \emph{prime subfield} of a field~\ref{def:field} $F$ is the subfield generated~\ref{def:generate-a-field} by the multiplicative identity~\ref{def:identity-element}. It is the smallest field~\ref{def:field} contained in $F$.
\end{definition}

\begin{definition}[quaternion algebra]\label{def:quaternion-algebra}
The \emph{quaternion algebra} $\mathbb H$ is the four-dimensional~\ref{def:dimension-of-vector-space} associative algebra~\ref{def:associative-algebra} over $\mathbb R$ spanned~\ref{def:span} by the elements $1,\mathbf i,\mathbf j,\mathbf k$ where $1$ is the identity element~\ref{def:identity-element} and the relations
1. $\mathbf i^2 = \mathbf j^2 = \mathbf k^2 = -1$;
2. $\mathbf{ij} = \mathbf k$;
3. $\mathbf{ji} = -\mathbf k$;
4. $\mathbf{jk} = \mathbf i$;
5. $\mathbf{kj} =-\mathbf i$;
6. $\mathbf{ki} = \mathbf j$;
7. $\mathbf{ik} = -\mathbf j$.
\end{definition}

\begin{definition}[rank]\label{def:rank}
The \emph{rank} of a linear transformation~\ref{def:linear-transformation} $T$ is the dimension~\ref{def:dimension-of-vector-space} of the image of $T$.
\end{definition}

\begin{definition}[separable field]\label{def:separable-field}
The extension~\ref{def:field-extension} $K$ of the field~\ref{def:field} $F$ is \emph{separable} over $F$ if every element in $K$ is the root of some separable~\ref{def:separable-polynomial} polynomial~\ref{def:polynomial-ring} in $F$.
\end{definition}

\begin{definition}[sign of permutation]\label{def:sign-of-permutation}
For $\sigma$ a permutation~\ref{def:symmetric-group} in $S_n$, the \emph{sign} of $\sigma$ denoted $\text{sgn}(\sigma)$ is the determinant~\ref{def:determinant} of the linear transformation~\ref{def:linear-transformation} $P_\sigma:\mathbb R^n\to\mathbb R^n$ given by $$P_\sigma(x_1,\dots,x_n) =(x_{\sigma(1)},\dots, x_{\sigma(n)}).$$
\end{definition}

\begin{definition}[simple field extension]\label{def:simple-field-extension}
A field extension~\ref{def:field-extension} generated~\ref{def:generate-a-field} by a single element is called \emph{simple}. The single element which generates this field is called \emph{primitive} for the extension.
\end{definition}

\begin{definition}[degree of field extension]\label{def:degree-of-field-extension}
The \emph{degree} of a field extension~\ref{def:field-extension} $F \subset E$ is the dimension~\ref{def:dimension-of-vector-space} of $E$ when considered as a vector space~\ref{def:vector-space} over $F$. If $E$ is simple field extension~\ref{def:simple-field-extension} generated~\ref{def:generate-a-field} by $\alpha$, then the degree is also equal to the degree~\ref{def:degree-of-polynomial} of the minimal polynomial in $E[x]$ of $\alpha$ over $F$.
\end{definition}

\begin{definition}[quadratic field extension]\label{def:quadratic-field-extension}
A field extension~\ref{def:field-extension} is \emph{quadratic} if it is of degree~\ref{def:degree-of-field-extension} two.
\end{definition}

\begin{definition}[quadratic integers]\label{def:quadratic-integers}
The \emph{quadratic integers} are algebraic integers~\ref{def:algebraic-integer} of degree~\ref{def:degree-of-field-extension} 2. Equivalently, they are solutions to monic, integer-valued polynomials~\ref{def:polynomial-ring} of degree~\ref{def:degree-of-polynomial} 2. That is, solutions to equations of the form $$x^2+bx+c=0.$$
\end{definition}

\begin{definition}[solvability by radicals]\label{def:solvability-by-radicals}
An element algebraic over the field~\ref{def:field} $F$ can be \emph{solved for by radicals} if it is an element of a field~\ref{def:field} $K$ which can be obtained as a succession of simple~\ref{def:simple-group} extensions~\ref{def:field-extension} adjoining $n$th roots of elements of $F$. A polynomial~\ref{def:polynomial-ring} $f(x) \in F[x]$ can be \emph{solved by radicals} if all of its roots can be solved *for* by radicals.
\end{definition}

\begin{definition}[spectrum]\label{def:spectrum}
Let $A$ be an algebra~\ref{def:algebra-over-a-field} over the field~\ref{def:field} $\mathbb k$ and let $a\in A$. The \emph{spectrum} of $a$ is $$\text{spec}(a)= \{\lambda\in\mathbb k\mid \lambda - a \text{ is not invertible}\}.$$
\end{definition}

\begin{definition}[splitting field]\label{def:splitting-field}
Let $F$ be a field~\ref{def:field} and let $f$ be a monic~\ref{def:monic} polynomial~\ref{def:polynomial-ring} with coefficients in $F$. Then a field extension~\ref{def:field-extension} $F\subset E$ is a \emph{splitting field} if there exist $\{a_i\}_{i=1}^d$ such that $$f(t) = \prod_{i=1}^d (t-a_i) \in E[t]$$ and moreover $E$ is generated~\ref{def:generate-a-field} by these $a_i$, i.e. $E= F(a_1,\dots, a_d)$.
\end{definition}

\begin{definition}[normal extension]\label{def:normal-extension}
Let $E$ be an algebraic extension~\ref{def:algebraic-extension} of the field~\ref{def:field} $F$ and let $\Omega$ be an algebraically closed~\ref{def:algebraically-closed} field~\ref{def:field} containing $F$. Then $E$ is a \emph{normal extension} of $F$ if it satisfies any of the following three equivalent conditions:
1. For each $\alpha\in E$, $E$ contains a splitting field~\ref{def:splitting-field} of the minimal polynomial of $\alpha$ over $F$.
2. $\{\alpha\in E\mid E \text{ contains a splitting field of the minimal polynomial of } \alpha \text{ over } F\}$ generates~\ref{def:generate-a-field} $E$ as a field extension~\ref{def:field-extension} of $F$.
3. For every $F$-embedding~\ref{def:field-embedding} $\phi:E\to \Omega$, we have $\phi(E) \subset E$. 

Equivalently, $E$ is a \emph{normal extension} of $F$ if it is the splitting field~\ref{def:splitting-field} of a *collection* of polynomials~\ref{def:polynomial-ring} in $F[x]$.
\end{definition}

\begin{definition}[Galois extension]\label{def:galois-extension}
A field extension~\ref{def:field-extension} $E/F$ is \emph{Galois} if it is algebraic~\ref{def:algebraic-extension}, normal~\ref{def:normal-extension}, and separable~\ref{def:separable-field}. The group~\ref{def:group} of automorphisms~\ref{def:automorphism-group} of $E$ that fix $F$, denoted $\text{Gal}(E/F)$, is called the \emph{Galois group} of $E$ over $F$. 

Equivalently, $E$ is \emph{Galois} over $F$ if the order~\ref{def:order-of-a-group} of the automorphism group $\text{Aut}(E/F)$ is equal to the degree~\ref{def:degree-of-field-extension} of $E$ over $F$.
\end{definition}

\begin{definition}[abelian extension]\label{def:abelian-extension}
The field extension~\ref{def:field-extension} $K/F$ is \emph{abelian} if it is Galois~\ref{def:galois-extension} and its Galois group is abelian.
\end{definition}

\begin{definition}[Galois group of a polynomial]\label{def:galois-group-of-a-polynomial}
The \emph{Galois group} of a separable~\ref{def:separable-polynomial} polynomial~\ref{def:polynomial-ring} $f$ over the field~\ref{def:field} $F$ is the Galois group~\ref{def:galois-extension} of the splitting field~\ref{def:splitting-field} of $f$ over $F$, which exists.
\end{definition}

\begin{definition}[structure constants]\label{def:structure-constants}
Let $L$ be a Lie algebra~\ref{def:lie-algebra} with basis~\ref{def:basis} $\{x_1,\dots,x_n\}$. The \emph{structure constants} $a_{ij}^k$ are the constants in the expressions $$[x_ix_j] = \sum_{k=1}^n a_{ik}^k x_k.$$
\end{definition}

\begin{definition}[symplectic algebra]\label{def:symplectic-algebra}
Let $V$ be a finite-dimensional~\ref{def:dimension-of-vector-space} vector space~\ref{def:vector-space} with $\text{dim}(V)=2\ell$ with basis~\ref{def:basis} $\{v_1,\dots,v_{2\ell}\}$. Define a nondegenerate skew-symmetric form $f$ on $V$ with the matrix $$s = \begin{bmatrix}0 & I_\ell \\ -I_\ell & 0\end{bmatrix}.$$ Note that the dimension~\ref{def:dimension-of-vector-space} being even is required for the existence of a nondegenerate bilinear~\ref{def:multilinear} form satisfying $f(v,w) = -f(w,v)$. Then denote by $\mathfrak{sp}(V)$ the \emph{symplectic algebra} consisting of all endomorphisms~\ref{def:endomorphism} $x$ of $V$ such that $f(x(v),w) = -f(v,x(w)).$
\end{definition}

\begin{definition}[tensor product of Lie algebra representations]\label{def:tensor-product-of-lie-algebra-representations}
Let $\mathfrak g$ be a Lie algebra~\ref{def:lie-algebra} and let $(V,\rho_V)$ and $(W,\rho_W)$ be representations~\ref{def:lie-algebra-representation} of $\mathfrak g$. Then the \emph{tensor product} of these representations~\ref{def:lie-algebra-representation} is the pair $(V\otimes W, \rho_{V\otimes W})$ where $$\rho_{V\otimes W}(x) = \rho_V(x) \otimes \text{Id} + \text{Id}\otimes \rho_W(x)$$ for all $x \in \mathfrak g$.
\end{definition}

\begin{definition}[transcendental]\label{def:transcendental}
Let $F$ be a field~\ref{def:field} and let $E$ be a field extension~\ref{def:field-extension} of $F$. Let $\Theta \in E$ and consider a nonzero polynomial~\ref{def:polynomial-ring} $f\in F[x]$.  If $f(\Theta)\neq 0$, then $\Theta$ is \emph{transcendental} over $F$. 

Equivalently, an element of a field~\ref{def:field} is \emph{transcendental} if it is not algebraic.
\end{definition}

\begin{definition}[trivial group representation]\label{def:trivial-group-representation}
Let $G$ be a group~\ref{def:group} and $V$ a vector space~\ref{def:vector-space}. The \emph{trivial representation} is the representation~\ref{def:group-representation} $\rho: G\to \text{GL}(V)$ such that for all $g\in G$, $\rho(g) = \text{id}_V$.
\end{definition}

\begin{definition}[unitary transformation]\label{def:unitary-transformation}
A linear transformation~\ref{def:linear-transformation} that preserves the norm~\ref{def:norm} of vectors~\ref{def:vector-space} is called \emph{unitary}.
\end{definition}

\begin{definition}[vector subspace]\label{def:vector-subspace}
Let $V$ be a vector space~\ref{def:vector-space}. A subset $W \subseteq V$ is a \emph{vector subspace} if for all $\alpha \in F$ the ground field and $u,v \in W$, $\alpha v \in W$ and $u+v \in W$.
\end{definition}

\begin{definition}[G-invariant subspace]\label{def:g-invariant-subspace}
Let $(V,\rho)$ be a representation~\ref{def:group-representation} of the group~\ref{def:group} $G$.  A vector subspace~\ref{def:vector-subspace} $W$ of $V$ is \emph{$G$-invariant} if for all $w\in W$ and $g\in G$, $\rho(g)w \in W$.
\end{definition}

\begin{definition}[invariant subspace]\label{def:invariant-subspace}
Let $T: V\to V$ be a linear transformation~\ref{def:linear-transformation} from the vector space~\ref{def:vector-space} $V$ to itself. A vector subspace~\ref{def:vector-subspace} $W\subseteq V$ is an \emph{invariant subspace} of $T$ if the image~\ref{def:image} of $W$ is contained in $W$. That is, $T(W)\subseteq W$.
\end{definition}

\begin{definition}[invariant polynomials under group action]\label{def:invariant-polynomials-under-group-action}
Let $\mathbb k$ be a field~\ref{def:field} and consider the algebra of polynomials in $n$ variables $$P_n =\bigoplus\limits_{d\geq 0} P_n^d.$$ Given a representation~\ref{def:group-representation} of a group~\ref{def:group} $G$ in $\mathbb k^n$, we get an action~\ref{def:group-action} of $G$ on $P_n$: for $f\in P_n$, $$g:f\mapsto g^*f.$$ Then each homogeneous subspace~\ref{def:vector-subspace} $P_n^d$ of $P_n$ is a $G$-invariant subspace~\ref{def:invariant-subspace}.
\end{definition}

\begin{definition}[orthogonal complement]\label{def:orthogonal-complement}
The \emph{orthogonal complement} of $W$ a subspace~\ref{def:vector-subspace} of the vector space~\ref{def:vector-space} $V$ is the set $$W^\perp = \{v \in V \mid \forall w \in W, \text{ } \langle v, w\rangle = 0\}.$$ This is itself a vector subspace~\ref{def:vector-subspace}:
write_proof
\end{definition}

\begin{definition}[projection]\label{def:projection}
A \emph{projection} is a linear transformation~\ref{def:linear-transformation} $P:V\to W$ between vector spaces~\ref{def:vector-space} $V$ and $W$ such that $P^2 = P$. That is, $P$ is idempotent~\ref{def:idempotent} when considered as an element of the ring~\ref{def:ring} of linear transformations~\ref{def:linear-transformation}.

Equivalently, for $V'$ a subspace~\ref{def:vector-subspace} of a vector space~\ref{def:vector-space} $V$, a linear transformation~\ref{def:linear-transformation} $P:V\to V$ is a \emph{projection} onto $V'$ if 
1. $P_{\vert V'} = \text{id}_V'$. That is, if the restriction of $P$ to $V'$ is equal to the identity function on $V'$, and
2. $\text{im}(P) \subseteq V'$.
\end{definition}

\begin{definition}[quotient vector space]\label{def:quotient-vector-space}
Define an equivalence relation~\ref{def:equivalence-relation} $\sim$ on a vector space~\ref{def:vector-space} $V$ with subspace~\ref{def:vector-subspace} $W$ by $$v_1 \sim v_2 \text{ if } v_1-v_2 \in W.$$ Then the \emph{quotient} of $V$ by $W$ is the set of equivalence classes~\ref{def:equivalence-class} of $\sim$ in $V$.
\end{definition}

\begin{definition}[quotient of Lie algebra by ideal]\label{def:quotient-of-lie-algebra-by-ideal}
Let $L$ be a Lie algebra~\ref{def:lie-algebra} and $I\subset L$ an ideal~\ref{def:ideal-of-a-lie-algebra}. Then the \emph{quotient} $L / I$ is constructed as follows:
- The underlying vector space~\ref{def:vector-space} is the usual quotient vector space~\ref{def:quotient-vector-space} $L/I$ when $L$ and $I$ are considered as vector spaces~\ref{def:vector-space}. 
- Lie multiplication is given by $[x+I,y+I] = [xy]+I$.
\end{definition}

\begin{definition}[real projective space]\label{def:real-projective-space}
The \emph{real projective space} $\mathbb RP^n$ of dimension $n$ is the quotient~\ref{def:quotient-vector-space} of $\mathbb R^{n+1}$ by the equivalence relation~\ref{def:equivalence-relation} $\sim$ given by $x\sim y$ if and only if $x=-y$. This amounts to considering $\mathbb RP^n$ as the set of lines through the origin in $\mathbb R^{n+1}$, or the unit sphere $S^n$ where we identify antipodal points $u$ and $-u$.
\end{definition}

\begin{definition}[Schur polynomials]\label{def:schur-polynomials}
Let $\Lambda=\{\mu = (\mu_1,\dots, \mu_n)\in\mathbb Z^n \mid \mu_1\geq \cdots\geq \mu_n \geq 0\} \subseteq \mathbb Z^n_{\geq 0}$. For each $\mu \in \Lambda$, define $\nabla_\mu \in \mathbb Z[x_1,\dots, x_n]^\text{sign}$ to be the antisymmetric~\ref{def:invariant-polynomials-under-group-action} polynomial~\ref{def:polynomial-ring} under the sign representation of the symmetric group~\ref{def:symmetric-group} given by $$\nabla_\mu = \sum_{s\in S_n} \text{sign}(s) x^{s(\mu)} = \det\left\vert \begin{matrix}x_1^{\mu_1}  & x_1^{\mu_2} & \cdots & x_1^{\mu_n} \\ x_2^{\mu_1} & x_2^{\mu_2} & \cdots & x_2^{\mu_n} \\ \vdots &\vdots & \ddots &\vdots \\ x_n^{\mu_1} & x_n^{\mu_2}&\cdots & x_n^{\mu_n}\end{matrix} \right\vert .$$ The second equality follows from the Leibniz formula for determinants~\ref{def:determinant}. 

Set $\rho \in \Lambda$ to $\rho = (n-1,n-2,\dots, 1, 0)$. Up to sign, $\nabla_\rho$ is the Vandermonde determinant. 

Because $\mathbb Z[x_1,\dots, x_n]^\text{sign}$ is a rank~\ref{def:rank-of-a-module}-one free~\ref{def:free-module} module~\ref{def:module-over-a-ring} over $\mathbb Z[x_1,\dots, x_n]^{S_n}$ with generator~\ref{def:generate-an-r-module} $\Delta_n$, it follows that $\nabla_\rho$ is also a generator of this module~\ref{def:module-over-a-ring}. Thus for every $\mu\in \Lambda$ there exists a unique polynomial~\ref{def:polynomial-ring} $s_\mu\in\mathbb Z[x_1,\dots, x_n]^{S_n}$ called the \emph{Schur polynomial} of $\mu$ such that $\nabla_{\mu + \rho} = s_\mu \cdot\nabla_\rho$. 

Explicitly, $s_\mu = \nabla_{\mu + \rho}/\nabla_\rho$.
\end{definition}

\begin{definition}[subalgebra]\label{def:subalgebra}
Let $L$ be a Lie algebra~\ref{def:lie-algebra}. A vector subspace~\ref{def:vector-subspace} $K$ of $L$ is a \emph{subalgebra} of $L$ if $[xy] \in K$ whenever $x,y \in K$. Moreover, $K$ is itself a Lie algebra~\ref{def:lie-algebra}.
\end{definition}

\begin{definition}[direct sum decomposition of a Lie algebra]\label{def:direct-sum-decomposition-of-a-lie-algebra}
Let $\mathfrak g$ be a Lie algebra~\ref{def:lie-algebra} with subalgebras~\ref{def:subalgebra} $\mathfrak g_1$ and $\mathfrak g_2$. Then $\mathfrak g$ decomposes as the direct sum of Lie algebras~\ref{def:direct-sum-of-lie-algebras} $\mathfrak g_1$ and $g_2$ if it is the direct sum of vector spaces $\mathfrak g_1\oplus \mathfrak g_2$ and $[X_1,X_2]=0$ for all $X_1\in\mathfrak g_1$ and $X_2\in \mathfrak g_2$.
\end{definition}

\begin{definition}[left algebra ideal]\label{def:left-algebra-ideal}
Let $A$ be an algebra~\ref{def:algebra-over-a-field} over the ring~\ref{def:ring} $R$. A \emph{left ideal} $I$ of $A$ is a subalgebra~\ref{def:subalgebra} of $A$ such that $ax\in I$ for all $a\in A$, $x\in I$. 

Equivalently (if we want to elide the subalgebra~\ref{def:subalgebra} condition), a \emph{left ideal} $I$ of $A$ is a subset $I\subset A$ such that for all $x,y\in I$, $a\in A$, and $r\in R$,
1. $x-y\in I$;
2. $rx\in I$;
3. $ax\in I$.
\end{definition}

\begin{definition}[linear Lie algebra]\label{def:linear-lie-algebra}
Any subalgebra~\ref{def:subalgebra} of a general linear algebra~\ref{def:general-linear-algebra} $\mathfrak{gl}(V)$ is called a \emph{linear Lie algebra}.
\end{definition}

\begin{definition}[right algebra ideal]\label{def:right-algebra-ideal}
Let $A$ be an algebra~\ref{def:algebra-over-a-field} over the ring~\ref{def:ring} $R$. A \emph{right ideal} $I$ of $A$ is a subalgebra~\ref{def:subalgebra} of $A$ such that $xa\in I$ for all $a\in A$, $x\in I$. 

Equivalently (if we want to elide the subalgebra~\ref{def:subalgebra} condition), a \emph{right ideal} $I$ of $A$ is a subset $I\subset A$ such that for all $x,y\in I$, $a\in A$, and $r\in R$,
1. $x-y\in I$;
2. $rx\in I$;
3. $xa\in I$.
\end{definition}

\begin{definition}[algebra ideal]\label{def:algebra-ideal}
An \emph{ideal} $I$ of an algebra~\ref{def:algebra-over-a-field} $A$ is a subset $I\subset A$ that is both a left algebra ideal~\ref{def:left-algebra-ideal} and a right algebra ideal~\ref{def:right-algebra-ideal}. In the case that $A$ is a commutative~\ref{def:commutative} algebra, these notions coincide and every right ideal is also a left ideal is also a two-sided ideal.

As long as the ring~\ref{def:ring} $R$ is not a rng~\ref{def:rng}, $A$ will inherit the definition of ideal~\ref{def:ring-ideal} from its underlying ring structure.
\end{definition}

\begin{definition}[generate an algebra ideal]\label{def:generate-an-algebra-ideal}
Let $A$ be an algebra~\ref{def:algebra-over-a-field}. Then its ideals~\ref{def:algebra-ideal} are \emph{generated} in the same way that ring ideals~\ref{def:ring-ideal} are generated~\ref{def:generate-a-ring-ideal} by the inheritance following from the fact that algebras are rings.
\end{definition}

\begin{definition}[subrepresentation]\label{def:subrepresentation}
Let $\phi:G\to GL(V)$ be a representation~\ref{def:group-representation} of the group~\ref{def:group} $G$. A \emph{sub-representation} is a representation~\ref{def:group-representation} $\psi:G\to GL(W)$ for $W$ a subspace~\ref{def:vector-subspace} of $V$.
\end{definition}

\begin{definition}[irreducible representation]\label{def:irreducible-representation}
A group representation~\ref{def:group-representation} is \emph{irreducible} if it has no nontrivial sub-representations~\ref{def:subrepresentation}.
\end{definition}

\begin{definition}[character table]\label{def:character-table}
The \emph{character table} of a group~\ref{def:group} $G$ is a table that displays the value taken by the characters~\ref{def:character-of-a-representation} of the irreducible representations~\ref{def:irreducible-representation} of $G$ on the different conjugacy classes~\ref{def:conjugacy-classes} of $G$. Recall that characters~\ref{def:character-of-a-representation} are class functions~\ref{def:class-function}
\end{definition}

\begin{definition}[complete reducibility]\label{def:complete-reducibility}
A representation~\ref{def:group-representation} $\rho$ of a group~\ref{def:group} $G$ is called \emph{completely reducible} or is said to have the \emph{complete reducibility property} if it can be written as a direct sum~\ref{def:direct-sum-of-group-representations} of irreducible~\ref{def:irreducible-representation} representations~\ref{def:group-representation}.
\end{definition}

\begin{definition}[isotypic component]\label{def:isotypic-component}
Let $S$ be a simple~\ref{def:irreducible-representation} Lie algebra representation~\ref{def:lie-algebra-representation} $\mathfrak g$. The \emph{isotypic component of type $S$} of a representation~\ref{def:lie-algebra-representation} $V$ is the direct sum~\ref{def:direct-sum-of-group-representations} of subrepresentations~\ref{def:subrepresentation} of $V$ isomorphic to $S$.
\end{definition}

\begin{definition}[semisimple representation of a Lie algebra]\label{def:semisimple-representation-of-a-lie-algebra}
Let $\mathfrak g$ be a Lie algebra~\ref{def:lie-algebra}. A \emph{semisimple representation} of $\mathfrak g$ is a linear~\ref{def:linear-transformation} representation~\ref{def:lie-algebra-representation} of $\mathfrak g$ that may be written as the direct sum~\ref{def:direct-sum-of-group-representations} of irreducible representations~\ref{def:irreducible-representation}.
\end{definition}

\begin{definition}[tensor product of vector spaces]\label{def:tensor-product-of-vector-spaces}
Given two vector spaces $V_1$ and $V_2$, $V_1 \otimes V_2$ is called the \emph{tensor product} of $V_1$ and $V_2$ and it is defined as follows: 

Let $F(V_1 \times V_2)$ be the free vector space generated by $V_1 \times V_2$, and let $K \subset F(V_1 \times V_2)$ be the subspace of $F(V_1 \times V_2)$ generated by the elements 
$$\begin{align*} &(v_1 + v_1', v_2) - (v_1,v_2) - (v_1', v_2)\\ &(v_1, v_2+v_2') - (v_1,v_2)- (v_1, v_2') \\ &(\lambda v_1, v_2) - \lambda (v_1, v_2) \\ &(v_1, \lambda v_2) - \lambda(v_1,v_2)\end{align*}$$ for all $v_1,v_1'\in V_1$, $v_2, v_2' \in V_2$, and $\lambda \in \mathbb R$. The \emph{tensor product} $V_1 \otimes V_2$ is the quotient vector space~\ref{def:quotient-vector-space} $$V_1 \otimes V_2 = F(V_1 \times V_2)/K.$$

There is also a definition in terms of a universal property.
\end{definition}

\begin{definition}[exterior power of a vector space]\label{def:exterior-power-of-a-vector-space}
Let $V$ be a vector space~\ref{def:vector-space}. The $m$-fold \emph{exterior power} of $V$ is the quotient vector space~\ref{def:quotient-vector-space} of the $m$-fold tensor product~\ref{def:tensor-product-of-vector-spaces} of $V$ by the subspace~\ref{def:vector-subspace} generated~\ref{def:span} by the pure tensors $$v_1\otimes\cdots\otimes v_m$$ for $v_i \in V$ such that $v_j = v_k$ for some $j\neq k$. It is denoted $\wedge^m(V)$ and the elements are written $v_1\wedge\cdots\wedge v_m$.
\end{definition}

\begin{definition}[adjugate]\label{def:adjugate}
Let $V$ be an $n$-dimensional~\ref{def:dimension-of-vector-space} vector space~\ref{def:vector-space} and let $T:V\to V$ be a linear transformation~\ref{def:linear-transformation}. The \emph{adugate} of $T$ is $$\text{adj}(T):V\to V$$ the unique linear transformation~\ref{def:linear-transformation} such that for all $v \in V$ and all $w \in \wedge^{n-1}(V)$ the $n-1$-fold exterior product~\ref{def:exterior-power-of-a-vector-space} of $V$, $$\text{adj}(T)(v) \wedge w = v \wedge (\wedge^{n-1}(T)(w)) \in \wedge^n(V).$$
\end{definition}

\begin{definition}[exterior power of group representations]\label{def:exterior-power-of-group-representations}
Let $(V,\rho)$ be a representation~\ref{def:group-representation} of the group~\ref{def:group} $G$. The $n$th \emph{exterior power} of $V_1$ and $V_2$ is a representation~\ref{def:group-representation} of $G$ over the exterior power of a vector space~\ref{def:exterior-power-of-a-vector-space} $\wedge^nV$ under the linear transformation~\ref{def:linear-transformation} $\rho:G\to \wedge^n(V)$ such that $$\rho(g)(v_1\wedge v_2\wedge \cdots \wedge v_n) = (\rho(g)v_1\wedge \rho(g)v_2\wedge\cdots\wedge\rho(g)v_n).$$
\end{definition}

\begin{definition}[exterior power of linear transformation]\label{def:exterior-power-of-linear-transformation}
Let $V$ be a vector space~\ref{def:vector-space} and let $T:V\to V$ be a linear transformation~\ref{def:linear-transformation}. For integer $m$, there exists a multilinear~\ref{def:multilinear} transformation $$\wedge^m(T): \wedge^m(V) \to \wedge^m(V)$$ from the $m$-fold exterior product~\ref{def:exterior-power-of-a-vector-space} to itself given by $$v_1\wedge\cdots\wedge v_m \mapsto T(v_1)\wedge\cdots\wedge T(v_m).$$
\end{definition}

\begin{definition}[minor]\label{def:minor}
Let $V$ be a vector space~\ref{def:vector-space} and let $T: V\to V$ be a linear transformation~\ref{def:linear-transformation} with corresponding matrix $A$. The matrix~\ref{def:matrix} coefficients of the matrix corresponding to the exterior product~\ref{def:exterior-power-of-linear-transformation} $\wedge^m(T)$ of $T$ are called the \emph{minors} of $A$.
\end{definition}

\begin{definition}[symmetric power of vector spaces]\label{def:symmetric-power-of-vector-spaces}
Let $V$ be a vector space~\ref{def:vector-space}. The $m$-fold \emph{symmetric product} of $V$ is the quotient vector space~\ref{def:quotient-vector-space} of the $m$-fold tensor product~\ref{def:tensor-product-of-vector-spaces} $V^{\otimes m}$ by the subspace generated~\ref{def:span} by the tensors $$v_1\otimes\cdots\otimes v_m - v_{\sigma(1)}\otimes \cdots\otimes v_{\sigma(m)}$$ for all $v_i \in V$ and all $\sigma \in S_m$ (the set of permutations~\ref{def:symmetric-group} of $\{1,\dots,m\}$). It is denoted $\text{Sym}^m(V)$.
\end{definition}

\begin{definition}[symmetric power of group representations]\label{def:symmetric-power-of-group-representations}
Let $(V,\rho)$ be a representation~\ref{def:group-representation} of the group~\ref{def:group} $G$. The $n$th \emph{symmetric power} of $V_1$ and $V_2$ is a representation~\ref{def:group-representation} of $G$ over the symmetric power of vector spaces~\ref{def:symmetric-power-of-vector-spaces} $\text{Sym}^n(V)$ under the linear transformation~\ref{def:linear-transformation} $\rho:G\to \text{Sym}^n(V)$ such that $$\rho(g)(v_1\odot v_2\odot \cdots \odot v_n) = (\rho(g)v_1\odot \rho(g)v_2\odot\cdots\odot\rho(g)v_n).$$
\end{definition}

\begin{definition}[tensor product of group representations]\label{def:tensor-product-of-group-representations}
Let $(V_1,\rho_1)$ and $(V_2,\rho_2))$ be representations~\ref{def:group-representation} of the group~\ref{def:group} $G$. The \emph{tensor product} of $V_1$ and $V_2$ is a representation~\ref{def:group-representation} of $G$ over the tensor product of vector spaces~\ref{def:tensor-product-of-vector-spaces} $V_1\otimes V_2$ under the linear transformation~\ref{def:linear-transformation} $\rho_1\otimes \rho_2 = \rho:G\to V_1\otimes V_2$ such that $\rho(g)(v_1\otimes v_2) = (\rho_1(g)v_1)\otimes (\rho_2(g)v_2)$.
\end{definition}

\begin{definition}[trace]\label{def:trace}
Let $T:V\to V$ be a linear transformation~\ref{def:linear-transformation} from the $n$-dimensional~\ref{def:dimension-of-vector-space} vector space~\ref{def:vector-space} $V$ to itself. The \emph{trace} of $T$ is $$\text{tr}(T) = \sum_{i=1}^n \langle T(e_i), e_i\rangle$$ where $\{e_i\}_{i=1}^n$ is an orthonormal basis for $V$.

Equivalently, if $V^*$ is the dual~\ref{def:dual-space} ov $V$, then the \emph{trace} of a linear transformation~\ref{def:linear-transformation} $T: V\to V$ is the functional~\ref{def:linear-transformation} $$\text{tr}:V^*\otimes V \to \mathbb R$$ from the tensor product~\ref{def:tensor-product-of-vector-spaces} of of $V^*$ and $V$ determined by $\phi \otimes v \mapsto \phi(v)$ for $\phi \in V^*$ and $v \in V$.
\end{definition}

\begin{definition}[special linear algebra]\label{def:special-linear-algebra}
Let $V$ be a finite-dimensional~\ref{def:dimension-of-vector-space} vector space~\ref{def:vector-space} with $\text{dim}(V) = \ell +1$. Let $\mathfrak{sl}(V)$ be the set of endomorphisms~\ref{def:endomorphism} of $V$ with trace~\ref{def:trace} equal to zero. Since $\text{Tr}(xy)= \text{Tr}(yx)$ and $\text{Tr}(x+y) =\text{Tr}(x)+\text{Tr}(y)$, $\mathfrak{sl}(V)$ is a subalgebra~\ref{def:subalgebra} of the general linear algebra~\ref{def:general-linear-algebra} $\mathfrak{gl}(V)$.
\end{definition}

\begin{definition}[volume]\label{def:volume}
The \emph{volume} of a rectangle~\ref{def:rectangle} $Q$ in $\mathbb R^n$ is the product of the lengths of the component intervals~\ref{def:component-interval} of $Q$: $$\text{vol}(Q) = (b_1-a_1)\cdots(b_n-a_n).$$
\end{definition}

\begin{definition}[Lebesgue outer measure]\label{def:lebesgue-outer-measure}
For any subset $A \subset\mathbb R^n$, the \emph{Lebesgue outer measure~\ref{def:outer-measure}} is given by $\mu^*(A) = \inf\{\sum\limits_{R\in\mathcal C}$vol~\ref{def:volume}$(R) \mid \mathcal C$ is a countable~\ref{def:countable} collection of rectangles~\ref{def:rectangle} covering $A\}$.
\end{definition}

\begin{definition}[Lebesgue measure]\label{def:lebesgue-measure}
The \emph{Lebesgue measure} on $\mathbb R^n$ is constructed as follows: 

A set $A \subset\mathbb R^n$ is \emph{Lebesgue measurable} if for every $S\subset \mathbb R^n$, we have the following relation in the Lebesgue outer measure~\ref{def:lebesgue-outer-measure} $\lambda^*$: $$\lambda^*(S) = \lambda^*(S\cap A) + \lambda^*(S\setminus A).$$
\end{definition}

\begin{definition}[measure zero]\label{def:measure-zero}
A subset $A \subset\mathbb R^n$ has \emph{measure zero} if for every $\varepsilon > 0$, there exists a countable covering of $A$ by rectangles~\ref{def:rectangle} $\{Q_i\}_{i\in\mathbb N}$ such that the sum of the volumes~\ref{def:volume} of the $Q_i$ is less than $\varepsilon$, i.e. that $$\sum_{i\in\mathbb N} \text{vol}(Q_i) < \varepsilon.$$
\end{definition}

\begin{definition}[weak convergence]\label{def:weak-convergence}
Let $X$ be a normed linear space~\ref{def:normed-linear-space} with dual space~\ref{def:dual-space} $X^*$. A sequence $\{x_i\}_{i\in\mathbb N}$ in $X$ converges \emph{weakly} to $x \in X$ if $\lim\limits_{i\to\infty} f(x_i)=f(x)$ for all $f\in X^*$.
\end{definition}

\begin{definition}[Weyl group]\label{def:weyl-group}
Let $G$ be a group~\ref{def:group} acting~\ref{def:group-action} on a set $S$. For a subgroup~\ref{def:subgroup} $H$ of $G$, denote by $NH$ the normalizer~\ref{def:normalizer-of-a-group} of $H$ in $G$ and let $WH = NH/H$. The quotient~\ref{def:quotient-by-normal-subgroup} $WH$ is the  \emph{Weyl group}.
\end{definition}

\begin{definition}[width]\label{def:width}
The \emph{width} of a rectangle~\ref{def:rectangle} $Q$ in $\mathbb R^n$ is the maximum of the lengths of the component intervals~\ref{def:component-interval} of $Q$: $$\text{width}(Q)= \max_{1\leq i \leq n} (b_i-a_i).$$
\end{definition}

\begin{definition}[Young subgroup]\label{def:young-subgroup}
Let $\lambda = \{\lambda_1,\dots,\lambda_k\}$ be a partition~\ref{def:partition-of-an-integer}. Let $\alpha_i = \{\lambda_i, \lambda_i + 1, \dots, \lambda_{i+1} - 1\}$ so that $\{1,\dots,n\} = \coprod_{i=1}^k \alpha_i$ as disjoint subsets. Corresponding to this decomposition is the \emph{Young subgroup~\ref{def:subgroup}} of the symmetric group~\ref{def:symmetric-group} $S_n$ given by $$S_{\alpha_1} \times \cdots \times S_{\alpha_k}$$ where $S_{\alpha_i} = \{\sigma\in S_n \mid \sigma(j) = j \text{ for all } j\notin \alpha_i\}$.
\end{definition}

\begin{definition}[zero divisor]\label{def:zero-divisor}
Let $R$ be a ring~\ref{def:ring}. A nonzero element $r\in R$ is a \emph{zero divisor} if there exists nonzer $s\in R$ such that $sr = 0$ or $rs = 0$.
\end{definition}

\begin{definition}[integral domain]\label{def:integral-domain}
A ring~\ref{def:ring} $R$ is an \emph{integral domain} if it is commutative~\ref{def:commutative} and has no zero divisors~\ref{def:zero-divisor}.
\end{definition}

\begin{definition}[annihilator of a submodule]\label{def:annihilator-of-a-submodule}
Let $R$ be an integral domain~\ref{def:integral-domain} and let $M$ be module~\ref{def:module-over-a-ring} over the ring $R$. Then for any submodule~\ref{def:submodule} $N$ of $M$, the \emph{annihilator} of $N$ is defined as $$\text{Ann}(N) = \{r\in R\mid rn=0\text{ for all } n\in N\}.$$
\end{definition}

\begin{definition}[associate elements in an integral domain]\label{def:associate-elements-in-an-integral-domain}
Let $R$ be an integral domain~\ref{def:integral-domain}. Then $a,b\in R$ are \emph{associate} if the ideals they generate~\ref{def:generate-a-ring-ideal} are equal. That is, they are associate if there exists a unit~\ref{def:unit-of-a-ring} $u\in R$ such that $ua=b$.
\end{definition}

\begin{definition}[irreducible element of an integral domain]\label{def:irreducible-element-of-an-integral-domain}
An element $x$ of an integral domain~\ref{def:integral-domain} is \emph{irreducible} if it is not a unit~\ref{def:unit-of-a-ring} and if whenever $x=ab$ for some $a,b\in R$, then either $a$ or $b$ is a unit~\ref{def:unit-of-a-ring}.
\end{definition}

\begin{definition}[prime element of an integral domain]\label{def:prime-element-of-an-integral-domain}
An element $x$ of an integral domain~\ref{def:integral-domain} is \emph{prime} if the ideal~\ref{def:ring-ideal} generate a generated~\ref{def:generate-a-ring-ideal} by $x$ is a prime ideal~\ref{def:prime-ideal}. That is, for all $a,b \in R$ such that $x$ divides~\ref{def:division-in-a-ring} $ab$, then either $x$ divides~\ref{def:division-in-a-ring} $a$ or $x$ divides~\ref{def:division-in-a-ring} $b$.
\end{definition}

\begin{definition}[prime number]\label{def:prime-number}
A \emph{prime number} is a prime element~\ref{def:prime-element-of-an-integral-domain} of $\mathbb N$. Equivalently, it is a natural number that is not divisible~\ref{def:division-in-a-ring} by any number other than $1$ and itself.
\end{definition}

\begin{definition}[cyclotomic polynomial]\label{def:cyclotomic-polynomial}
Let $p$ be a prime~\ref{def:prime-number}. Then the $p$th \emph{cyclotomic polynomial~\ref{def:polynomial-ring}} is the polynomial~\ref{def:polynomial-ring} over $\mathbb Z$ whose roots are the primitive~\ref{def:primitive-root-of-unity} $p$th roots of unity. It is denoted $\Phi_p (x) = x^{p-1} + x^{p-2} + \cdots + x + 1$.
\end{definition}

\begin{definition}[elementary abelian p-group]\label{def:elementary-abelian-p-group}
An \emph{elementary abelian $p$-group} is an abelian group~\ref{def:group} $G$ in which every nontrivial~\ref{def:identity-element} element has order~\ref{def:order-of-a-group-element} $p$ where $p$ is prime~\ref{def:prime-number}.
\end{definition}

\begin{definition}[linearly independent group elements]\label{def:linearly-independent-group-elements}
Let $A$ be an elementary abelian p-group~\ref{def:elementary-abelian-p-group} and let $\mathbf a\in(a_1,\dots,a_r)\in A^r$. Let $f_\mathbf a:(\mathbb Z/p\mathbb Z)\to A$ be the unique homomorphism~\ref{def:group-homomorphism} satisfying $$f_\mathbf a(1,0,\dot,0) = a_1, \dots, f_\mathbf a(0,\dots, 1) = a_r.$$ If the function $f_\mathbf a$ is injective~\ref{def:injective}, then the $a_i$ are \emph{linearly independent}.
\end{definition}

\begin{definition}[basis for abelian group]\label{def:basis-for-abelian-group}
Let $A$ be an elementary abelian p-group~\ref{def:elementary-abelian-p-group} and let $\mathbf a\in(a_1,\dots,a_r)\in A^r$. Then the $a_i$ form a \emph{basis} for $A$ if they generate~\ref{def:generating-set-of-a-group} $A$ and if they are linearly independent~\ref{def:linearly-independent-group-elements} in the sense that
\end{definition}

\begin{definition}[principal ideal domain]\label{def:principal-ideal-domain}
An integral domain~\ref{def:integral-domain} $R$ is a \emph{principal ideal domain} if every ideal~\ref{def:ring-ideal} of $R$ is principal~\ref{def:principal-ideal}.
\end{definition}

\begin{definition}[torsion free module]\label{def:torsion-free-module}
Let $R$ be an integral domain~\ref{def:integral-domain} and let $M$ be an $R$-module~\ref{def:module-over-a-ring}. Then $M$ is \emph{torsion-free} if its torsion submodule~\ref{def:submodule} is trivial.
\end{definition}

\begin{definition}[unique factorization domain]\label{def:unique-factorization-domain}
An integral domain~\ref{def:integral-domain} $R$ is a \emph{unique factorization domain} if
1. Every nonzero, non-unit~\ref{def:unit-of-a-ring} $r\in R$ can be written as a product of irreducibles~\ref{def:irreducible-element-of-an-integral-domain}; and
2. Whenever $p_1\cdots p_n=q_1\cdots q_m$ for $p_i,q_j$ irreducible~\ref{def:irreducible-element-of-an-integral-domain} and $m,n\in\mathbb N$, then $m=n$ and there is a relabeling/reordering of the $q_j$ such that $q_i$ is associate~\ref{def:associate-elements-in-an-integral-domain} to $p_i$ for all $1\leq i\leq n$.

In particular, for every association~\ref{def:associate-elements-in-an-integral-domain} class~\ref{def:equivalence-class} of irreducibles~\ref{def:irreducible-element-of-an-integral-domain}, we can choose a representative, and let $P$ be the collection of such representatives so that every $r\in R$ can be written uniquely as $$r= u\prod_{p\in P} p^{e_p}$$ for $u$ a unit~\ref{def:unit-of-a-ring} and all but finitely many of the $e_p$ equal to zero.
\end{definition}

\begin{definition}[σ-algebra]\label{def:σ-algebra}
A \emph{$\sigma$-algebra} on the set $X$ is a subset $\mathcal M \subset \mathcal P(X)$ of the power set~\ref{def:power-set} of $X$ satisfying the following properties:
1. $X \in \mathcal M$;
2. $\mathcal M$ is closed under complementation, i.e. if $A \in \mathcal M$, then $X\setminus A \in \mathcal M$;
3. $\mathcal M$ is closed under countable~\ref{def:countable} unions, i.e. if $A_i \in \mathcal M$ for $i \in \mathbb N$, then $\bigcup\limits_{i\in\mathbb N} A_i \in \mathcal M$.

!meme
\end{definition}

\begin{definition}[generate a σ-algebra]\label{def:generate-a-σ-algebra}
The σ-algebra~\ref{def:σ-algebra} \emph{generated} by a collection of subsets is the σ-algebra~\ref{def:σ-algebra} containing all of the subsets such that any σ-algebra~\ref{def:σ-algebra} containing the sets also contains it.
\end{definition}

\begin{definition}[commutator subgroup]\label{def:commutator-subgroup}
Let $u,v$ be elements of the group~\ref{def:group} $G$. Define $[u,v] = uvu^{-1}v^{-1}$. The subgroup~\ref{def:subgroup} of $G$ generated~\ref{def:generate-a-σ-algebra} by $[u,v]$ for all pairs $u,v\in G$ is denoted $[G,G]$ and is called the \emph{commutator subgroup} of $G$.
\end{definition}

\begin{definition}[abelianization]\label{def:abelianization}
The quotient~\ref{def:quotient-by-normal-subgroup} of the group $G$ by its commutator subgroup~\ref{def:commutator-subgroup} $[G,G]$ (which is normal~\ref{def:normal-subgroup}) is called the \emph{abelianization} of $G$. It is an abelian group.
\end{definition}

\begin{definition}[derived series of a group]\label{def:derived-series-of-a-group}
The \emph{derived series} of a group~\ref{def:group} $G$ is the descending series of consecutive derived subgroups~\ref{def:commutator-subgroup} of $G$: $$G = G^{(0)} \supseteq G^{(1)} = [G,G] \supseteq \cdots \supseteq G^{(i)}\supseteq \cdots$$ where $G^{(i+1)} = [G^{(i)},G^{(i)}]$. Note that this series may not reach the trivial subgroup.
\end{definition}

\begin{definition}[lower central series]\label{def:lower-central-series}
Let $G$ be a group~\ref{def:group}. The \emph{lower central series} for $G$ is defined as follows: Let $G_0=G$, let $G_1 = [G,G]$, and for each $i \geq 1$, define $G_{i+1} = [G,G_{i}]$ where the bracket notation denotes a particular subgroup~\ref{def:subgroup} of the commutator subgroup~\ref{def:commutator-subgroup}: $[G,H] = \{ghg^{-1}h^{-1}\mid g\in G, h\in H\}$. Note that this series may or may not reach the trivial~\ref{def:identity-element} subgroup~\ref{def:subgroup}.
\end{definition}

\begin{definition}[measurable]\label{def:measurable}
A \emph{measurable space} is the pair $(X,\mathcal M)$ consisting of a set $X$ and a σ-algebra~\ref{def:σ-algebra} $\mathcal M$ of subsets of $X$. A subset $E$ of $X$ is \emph{measurable} if it is an element of $\mathcal M$.
\end{definition}

\begin{definition}[Borel space]\label{def:borel-space}
A \emph{Borel space} is a measurable space~\ref{def:measurable} $(X, B)$ given by the Borel σ-algebra on $X$.
\end{definition}

\begin{definition}[measurable function]\label{def:measurable-function}
Let $(X,\Sigma)$ and $(Y, T)$ be measurable spaces~\ref{def:measurable}. A function $f:X\to Y$ is \emph{measurable} if the preimage under $f$ of every set in $T$ is in $\Sigma$.
\end{definition}

\begin{definition}[Borel function]\label{def:borel-function}
A \emph{Borel function} is a measurable function~\ref{def:measurable-function} between Borel spaces.
\end{definition}

\begin{definition}[measure space]\label{def:measure-space}
A \emph{measure} $\mu$ on a measurable~\ref{def:measurable} space $(X,\mathcal M)$ is a function $\mu:\mathcal M \to [0,\infty]$ valued on the extended real numbers for which $\mu(\emptyset) = 0$ for any countable~\ref{def:countable}, pairwise disjoint collection $\{E_i\}_{i\in\mathbb N}$ of measurable sets, $$\mu\left(\bigcup_{i\in\mathbb N} E_i\right) = \sum_{i\in\mathbb N} \mu(E_i).$$

A \emph{measure space} is a measurable~\ref{def:measurable} space $(X,\mathcal M)$ together with a \emph{measure} $\mu$ on $\mathcal M$.
\end{definition}

\begin{definition}[absolute continuity of measure]\label{def:absolute-continuity-of-measure}
Let $(X,\Sigma)$ be a measurable space~\ref{def:measurable} and let $\mu$ and $\nu$ be two measures~\ref{def:measure-space} on $(X,\Sigma)$. Then $\mu$ is \emph{absolutely continuous with respect to $\nu$} if $\mu(A) =0$ for every $A\in \Sigma$ for which $\nu(A)=0$.
\end{definition}

\begin{definition}[almost everywhere]\label{def:almost-everywhere}
Let $(X,\Sigma,\mu)$ be a measure space~\ref{def:measure-space}. A property $P$ holds \emph{almost everywhere} in $X$ if there exists a set $N \in \Sigma$ such that $\mu(N) = 0$ and $P$ is true for all $x \in X\setminus N$.
\end{definition}

\begin{definition}[almost uniform convergence]\label{def:almost-uniform-convergence}
Let $(X,\Sigma,\mu)$ be a measure space~\ref{def:measure-space} and let $\{f_n\}_{n\in\mathbb N}$ be a sequence of (measurable) functions $f_n: X\to \mathbb R$. The sequence converges \emph{almost uniformly} to $f$ if for all $\varepsilon > 0$, there exists a set $A\in \Sigma$ of measure~\ref{def:measure-space} less than $\varepsilon$ such that $f_n$ converges uniformly to $f$ on $X\setminus A$.
\end{definition}

\begin{definition}[bounded measure space]\label{def:bounded-measure-space}
A measure space~\ref{def:measure-space} $(X,\Sigma, \mu)$ is \emph{bounded} if $\mu(X) < \infty$.
\end{definition}

\begin{definition}[essential infimum]\label{def:essential-infimum}
Let $(X,\Sigma, \mu)$ be a measure space~\ref{def:measure-space}. The \emph{essential supremum} $\alpha$ is the largest real number $\alpha$ such that $\mu(\{x\mid f(x) < \alpha\}) = 0$.
\end{definition}

\begin{definition}[essential supremum]\label{def:essential-supremum}
Let $(X,\Sigma, \mu)$ be a measure space~\ref{def:measure-space}. The \emph{essential supremum} $\alpha$ of a function $f:X\to \mathbb R$ is the smallest real number $\alpha$ such that $\mu(\{x\mid f(x) > \alpha\}) = 0$.
\end{definition}

\begin{definition}[Hardy-Littlewood maximal operator]\label{def:hardy-littlewood-maximal-operator}
Let $(X,\Sigma, \mu)$ be a measure space~\ref{def:measure-space}. For $f \in L_1(X,\mu)$, the \emph{Hardy-Littlewood maximal operator} $Mf$ is defined as follows: $$(Mf)(x)=\sup_{r>0}\frac{1}{\mu(B(x,r))}\int_{B(x,r)}\vert f\vert \text d\mu.$$
\end{definition}

\begin{definition}[Lebesgue integral of characteristic function]\label{def:lebesgue-integral-of-characteristic-function}
Let $(X, \Sigma, \mu)$ be a measure space~\ref{def:measure-space}. The \emph{Lebesgue integral of the characteristic function} of a measurable~\ref{def:measurable} set $A \in \Sigma$ is defined to be $$\int \chi_Ad\mu = \mu(A).$$
\end{definition}

\begin{definition}[Lusin function]\label{def:lusin-function}
Let $(X,\Sigma, \mu)$ and $(Y, T, \nu)$ be measure spaces~\ref{def:measure-space}. A function $f:X\to Y$ is \emph{Lusin} or \emph{Luzin} if the image of every set of $\mu$-measure~\ref{def:measure-space} zero has zero $\nu$-measure~\ref{def:measure-space}.
\end{definition}

\begin{definition}[measurable hull]\label{def:measurable-hull}
Let $(X,\Sigma, \mu)$ be a measure space~\ref{def:measure-space} and let $A \subset X$. The \emph{measurable hull} $H$ of $A$ is the measurable~\ref{def:measurable} set $H\supseteq A$ such that for any measurable~\ref{def:measurable} set $F$ satisfying $A\subseteq F\subseteq H$, we have $\mu(F) = \mu(H)$.
\end{definition}

\begin{definition}[outer regular]\label{def:outer-regular}
Let $(X,\Sigma,\mu)$ be a measure space~\ref{def:measure-space}. A measurable~\ref{def:measurable} subset $A$ of $X$ is \emph{outer regular} if $$\mu(A) = \inf\{\mu(G) \mid G\supseteq A, G\text{ open and measurable}\}.$$
\end{definition}

\begin{definition}[p-norm]\label{def:p-norm}
Let $(X,\Sigma,\mu)$ be a measure space~\ref{def:measure-space}. The \emph{$p$-norm~\ref{def:norm}} of a measurable function~\ref{def:measurable-function} $f \to \mathbb R$ is $$\vert \vert f\vert \vert _p = \left(\int_X \vert f\vert ^p\text d\mu\right)^{1/p}.$$
\end{definition}

\begin{definition}[product measure]\label{def:product-measure}
Let $(X_1,\Sigma_1,\mu_1)$ and $(X_2,\Sigma_2,\mu_2)$ be two measure spaces~\ref{def:measure-space}, and let $\Sigma_1\otimes \Sigma_2$ be the σ-algebra generated~\ref{def:generate-a-σ-algebra} by sets of the form $B_1\times B_2$ for $B_1\in \Sigma_1$, $B_2\in \Sigma_2$. A \emph{product measure} $\mu_1\times \mu_2$ is a measure~\ref{def:measure-space} on the measurable space~\ref{def:measurable} $(X_1\times X_2, \Sigma_1\otimes\Sigma_2)$ such that $$(\mu_1\times\mu_2)(B_1\times B_2) = \mu_1(B_1)\mu_2(B_2)$$ for all $B_1\in \Sigma_1$, $B_2\in \Sigma_2$.
\end{definition}

\begin{definition}[regular]\label{def:regular}
A measure~\ref{def:measure-space} is \emph{regular} if it is both inner regular and outer regular~\ref{def:outer-regular}.
\end{definition}

\begin{definition}[signed measure]\label{def:signed-measure}
A \emph{signed measure} on the measurable space~\ref{def:measurable} $(X,\Sigma)$ is a function $$\mu: \Sigma \to \mathbb R \cup \{-\infty\} \text{ or } \mathbb R \cup \{\infty\}$$ such that $\mu(\emptyset) =0$ and $\mu$ is σ-additive. Note that $\mu$ cannot take on both $-\infty$ and $\infty$ as values.
\end{definition}

\begin{definition}[positive variation]\label{def:positive-variation}
Let $(X,\Sigma)$ be a measurable space~\ref{def:measurable} with $\mu$ a signed measure~\ref{def:signed-measure}. The \emph{positive variation} of $\mu$ is the function $\pi$ given by $$\pi(A) = \sup_{S\subset A} \mu(S).$$
\end{definition}

\begin{definition}[negative variation]\label{def:negative-variation}
Let $(X,\Sigma)$ be a measurable space~\ref{def:measurable} with $\mu$ a signed measure~\ref{def:signed-measure}. The \emph{negative variation} of $\mu$ is the function $\nu$ given by $$\nu(A) = \pi(A)-\mu(A)$$ where $\pi$ is the positive variation~\ref{def:positive-variation} of $\mu$.
\end{definition}

\begin{definition}[simple function]\label{def:simple-function}
Let $(X,\Sigma)$ be a measurable space~\ref{def:measurable}, and let $\{A_i\}_{i=1}^n$ be a sequence of disjoint measurable~\ref{def:measurable} sets, and let $\{a_i\}_{i=1}^n$ be a sequence of real numbers. A \emph{simple function} is a function $f:X\to \mathbb R$ of the form $$f(x) = \sum_{k=1}^n a_k\chi_{A_k}(x)$$ where $\chi_A$ is the characteristic function of $A$.
\end{definition}

\begin{definition}[Lebesgue integral of a simple function]\label{def:lebesgue-integral-of-a-simple-function}
Let $(X,\Sigma, \mu)$ be a measure space~\ref{def:measure-space}. Let $f = \sum\limits_{k=1}^n a_k \chi_{A_k}$ be a simple function~\ref{def:simple-function} such that $\mu(A_k) < \infty$ for $a_k \neq 0$.  The \emph{Lebesgue integral of the simple function~\ref{def:simple-function}} $f$ is defined as $$\int\left(\sum_{k=1}^n a_k\chi_{A_k}\right)d\mu = \sum_{k=1}^n a_k \int\chi_{A_k}d\mu = \sum_{k=1}^n a_k\mu(A_k),$$ consistent with the Lebesgue integral of characteristic functions~\ref{def:lebesgue-integral-of-characteristic-function}.

If $B$ is a measurable~\ref{def:measurable} subset of $X$ and $g = \sum\limits_{k=1}^n b_k\chi_{B_k}$ is a measurable~\ref{def:measurable-function} simple function~\ref{def:simple-function} on $B$, define $$\int_B gd\mu = \int\chi_Bgd\mu = \sum_{k=1}^nb_k\mu(B_k\cap B).$$
\end{definition}

\begin{definition}[Lebesgue integral]\label{def:lebesgue-integral}
To define the \emph{Lebesgue integral} for a general real-valued function, we must first define it for nonnegative functions.

Let $(X,\Sigma,\mu)$ be a measure space~\ref{def:measure-space} and let $f$ be a nonnegative, measurable~\ref{def:measurable-function}, extended-real-valued function on $X$ (that is, $f(x) \geq 0$ for all $x \in X$, and it may be the case that $f(x) = +\infty$). Define $$\int_X fd\mu = \sup\left\{\int_X sd\mu \mid 0\leq s\leq f, s \text{ simple}\right\}$$ where the supremum~\ref{def:supremum} is taken over the Lebesgue integrals~\ref{def:lebesgue-integral-of-a-simple-function} of nonnegative simple functions~\ref{def:simple-function} $s$ less than $f$.

Now let $f$ be a measurable~\ref{def:measure-space} extended-real-valued function on $X$. Then $f$ may be written as $f = f^+ - f^-$ where $$f^+(x) =\begin{cases}f(x) & f(x)>0 \\0 &\text{otherwise} \end{cases} \quad\text{and}\quad f^-(x) = \begin{cases}-f(x) & f(x)<0 \\ 0 &\text{otherwise} \end{cases}$$ so that $f^+$ and $f^-$ are both nonnegative and measurable~\ref{def:measure-space}. Note that $\vert f\vert  = f^+ + f^-$. The Lebesgue integral of the measurable function~\ref{def:measurable-function} $f$ \emph{exists} if at least one of $\int f^+d\mu$ and $\int f^- d\mu$ is finite. In this case, define $$\int f d\mu = \int f^+d\mu - \int f^-d\mu.$$ If $\int\vert f\vert d\mu < \infty$, then $f$ is \emph{Lebesgue integrable}.
\end{definition}

\begin{definition}[Lp space]\label{def:lp-space}
Let $1\leq p < \infty$ and let $(X,\Sigma, \mu)$ be a measure space~\ref{def:measure-space}. Let $F$ be the set of functions $f: X\to \mathbb R$ such that $$\vert \vert f\vert \vert _p = \left(\int_X\vert f\vert ^p\text d\mu\right)^{1/p}\ < \infty$$ where $\vert \vert f\vert \vert _p$ is the p-norm~\ref{def:p-norm}. Equivalently, $F$ is the set of functions such that the $p$th power of their absolute value is integrable~\ref{def:lebesgue-integral}. Let $N = \{f \in F \mid f = 0 \text{almost everywhere}\}$. Because the integral of a positive function is zero if and only if it is zero almost everywhere, we have that $N$ is the kernel~\ref{def:kernel-of-linear-transformation} of the p-norm~\ref{def:p-norm}. The normed linear space~\ref{def:normed-linear-space} given by the quotient vector space~\ref{def:quotient-vector-space} of $F$ by $N$ is $L_p(X,\mu)$.

For $p = \infty$, consider the set $G$ of functions $f:X \to \mathbb R$  for which there exists $c \in\mathbb R$ such that $\mu(\{\vert f\vert  > c\}) = 0$ and take the quotient vector space~\ref{def:quotient-vector-space} of this set by the equivalence relation~\ref{def:equivalence-relation} $\sim$ on $G$ given by $f\sim g$ if $f(x) = g(x)$ almost everywhere~\ref{def:almost-everywhere} on $X$.  This space with the essential supremum~\ref{def:essential-supremum} of $f$ as the norm~\ref{def:norm} is $L_\infty(X,\mu)$.
\end{definition}

\begin{definition}[essentially bounded]\label{def:essentially-bounded}
The elements of L~\ref{def:lp-space}$_\infty$ are called \emph{essentially bounded}. That is, their absolute values have finite essential supremum~\ref{def:essential-supremum}.
\end{definition}

\begin{definition}[Radon-Nikodym derivative]\label{def:radon-nikodym-derivative}
Let $(X,\Sigma)$ be a measurable space~\ref{def:measurable} with measures~\ref{def:measure-space} $\mu$ and $\nu$ such that $\nu$ is absolutely continuous~\ref{def:absolute-continuity-of-measure} with respect to $\mu$. The \emph{Radon-Nikodym derivative} $\frac{\text d\nu}{\text d\mu}$ is the function $f \in$ L~\ref{def:lp-space}$_1(\mu)$ such that $\nu(A) = \int_A f\text d\mu$.
\end{definition}

\begin{definition}[singular measure]\label{def:singular-measure}
Let $(X,\Sigma)$ be a measurable space~\ref{def:measurable} and let $\mu$ and $\nu$ both be measures~\ref{def:measure-space} on $(X,\Sigma)$. Then $\mu$ and $\nu$ are \emph{singular}, written $\mu \perp \nu$, if there exist disjoint sets $A,B\subset X$ such that $\mu \equiv 0$ outside of $A$ and $\nu \equiv 0$ outside of $B$. That is, $\mu(A') = 0$ for all $A'\subset X\setminus A$ and $\nu(B') =0$ for all $B'\subset X\setminus B$.
\end{definition}

\begin{definition}[total variation]\label{def:total-variation}
Let $(X,\Sigma)$ be a measurable space~\ref{def:measurable} with $\mu$ a signed measure~\ref{def:signed-measure}. The \emph{total variation} of $\mu$ is the function $\tau$ given by $$\tau(A) = \pi(A) + \nu(A)$$ where $\pi$ is the positive variation~\ref{def:positive-variation} of $\mu$ and $\nu$ is the negative variation~\ref{def:negative-variation} of $\mu$.
\end{definition}

\begin{definition}[σ-finite]\label{def:σ-finite}
Let $(X,\Sigma, \mu)$ be a measure space~\ref{def:measure-space}. The measure~\ref{def:measure-space} $\mu$ is \emph{σ-finite} if it satisfies one of the following equivalent criteria:
1. $X$ can be covered with at most countably many (pairwise disjoint, this is not required) measurable~\ref{def:measurable} sets of finite measure~\ref{def:measure-space};
2. $X$ can be covered with a monotone sequence of measurable~\ref{def:measurable} sets of finite measure~\ref{def:measure-space};
3. There exists a positive measurable function~\ref{def:measurable-function} $f$ whose integral~\ref{def:lebesgue-integral} is finite.
\end{definition}

\begin{definition}[absolute convergence]\label{def:absolute-convergence}
Let $A$ be a Banach algebra~\ref{def:banach-algebra} with norm~\ref{def:norm} $\vert \cdot \vert$. Given a sequence $\{a_i\}_{i\in \mathbb N}$ of elements of $A$, the series $$\sum_{i=1}^\infty a_i$$ is \emph{absolutely convergent} if there exists $L \in\mathbb R$ such that  $$\sum_{i=1}^\infty \vert a_i\vert =L.$$

This means that the partial sums $$\sum_{i=1}^n a_i$$ form a Cauchy sequence~\ref{def:cauchy-sequence} in $A$, and since $A$ is complete~\ref{def:complete-metric-space} by definition, the sequence must have a limit in $A$. We denote this limit by the infinite sum.
\end{definition}

\begin{definition}[adjoint representation of a Lie group]\label{def:adjoint-representation-of-a-lie-group}
Let $G$ be a Lie group~\ref{def:lie-group} with associated~\ref{def:lie-correspondence} Lie algebra~\ref{def:lie-algebra} $\mathfrak g$. Then the adjoint map $\text{Ad}:G \to \text{GL}(\mathfrak g)$ is given by $$\text{Ad}_A(X) = AXA^{-1}.$$ Because the adjoint map is a Lie algebra homomorphism, it is a representation~\ref{def:lie-algebra-representation} of $\mathfrak g$, and therefore we may refer to $\text{ad}$ as the \emph{adjoint representation of  $G$} acting~\ref{def:group-action} on $\mathfrak g$ when considered as a vector space~\ref{def:vector-space}.
\end{definition}

\begin{definition}[Alexandroff space]\label{def:alexandroff-space}
An \emph{Alexandroff space} is a topological space~\ref{def:topological-space} whose topology~\ref{def:topological-space} is closed under arbitrary intersections, not just finite ones.
\end{definition}

\begin{definition}[analytic function]\label{def:analytic-function}
A function $f:\mathbb U\to \mathbb C$ is \emph{analytic} on the open~\ref{def:open} set $U\subset\mathbb C$ if for any $z_0\in U$, there exists a power series~\ref{def:power-series} $$f(z) = \sum_{n=0}^\infty a_n(z-z_0)^n$$ with $a_n \in \mathbb C$ for all $n$ that converges to $f(z)$ in a neighborhood of $z_0$.
\end{definition}

\begin{definition}[atlas]\label{def:atlas}
An \emph{atlas} for a topological space~\ref{def:topological-space} $M$ is a collection of charts~\ref{def:chart} on $M$ that covers $M$.
\end{definition}

\begin{definition}[Baire class]\label{def:baire-class}
The \emph{Baire classes} of continuous~\ref{def:continuous} real-valued functions on a topological space~\ref{def:topological-space} $X$ are defined as follows:
- continuous~\ref{def:continuous} functions are Baire-0;
- functions that are the pointwise limit of a sequence of continuous~\ref{def:continuous} functions are Baire-1;
- for all $\alpha \in \mathbb N$, functions that are the pointwise limit of a sequence functions of Baire class less than $\alpha$ are Baire-$\alpha$.
\end{definition}

\begin{definition}[Baire property]\label{def:baire-property}
A subset $A$ of a topological space~\ref{def:topological-space} $X$ has the \emph{Baire property} if there exists an open~\ref{def:open} set $G$ such that the symmetric difference~\ref{def:symmetric-difference} of $A$ and $G$ ($A \triangle G$) is meager~\ref{def:meager}.
\end{definition}

\begin{definition}[Baire space]\label{def:baire-space}
A topological space~\ref{def:topological-space} is a \emph{Baire space} if every nonempty open~\ref{def:open} set in $X$ is nonmeager~\ref{def:nonmeager}.

Equivalently, a topological space~\ref{def:topological-space} is a \emph{Baire space} if every countable~\ref{def:countable} intersection of open~\ref{def:open} dense~\ref{def:dense} sets is also dense~\ref{def:dense}.
\end{definition}

\begin{definition}[Banach algebra]\label{def:banach-algebra}
A \emph{Banach algebra} is an algebra~\ref{def:algebra-over-a-field} equipped with a complete~\ref{def:complete-metric-space} multiplicative~\ref{def:multiplicative-norm} norm~\ref{def:norm}.
\end{definition}

\begin{definition}[Banach space]\label{def:banach-space}
A normed linear space~\ref{def:normed-linear-space} $(X, \vert \vert \cdot\vert \vert )$ is \emph{Banach} if it is complete~\ref{def:complete-metric-space} under the metric~\ref{def:metric-space} induced by the norm.
\end{definition}

\begin{definition}[Banach-Mazur game]\label{def:banach-mazur-game}
A \emph{Banach-Mazur game} $MB(X,Y, \mathcal W)$ is a topological game~\ref{def:topological-game} played as follows:

Let $X$ be a topological space~\ref{def:topological-space}, $Y$ a fixed subset of $X$, and $\mathcal W$ a family of subsets of $X$ such that
1. each element of $\mathcal W$ has nonempty interior~\ref{def:interior}.
2. every nonempty subset of $X$ contains an element of $\mathcal W$. 

Players 1 and 2 take turns picking elements of $\mathcal W$ in a sequence $W_1\supset W_2\supset\cdots$. Player 1 wins if and only if $$
Y\cap \left(\bigcap_{n<\omega} W_n\right)$$ is nonempty. If Player 1 loses, Player 2 wins.
\end{definition}

\begin{definition}[based topological space]\label{def:based-topological-space}
A \emph{based} or \emph{pointed topological space} is a topological space~\ref{def:topological-space} $X$ together with a point $x_0\in X$ called the \emph{basepoint}. Unless the basepoint is clear, write $(X, x_0)$ for the based space.
\end{definition}

\begin{definition}[Borel σ-algebra]\label{def:borel-σ-algebra}
The \emph{Borel σ-algebra} on a topological space~\ref{def:topological-space} $X$ is the σ-algebra~\ref{def:σ-algebra} generated~\ref{def:generate-a-σ-algebra} by the open~\ref{def:open} sets in $X$.
\end{definition}

\begin{definition}[boundary]\label{def:boundary}
Let $X$ be a topological space~\ref{def:topological-space}. The \emph{boundary} $\partial A$ of a set $A \subset X$ is the intersection of the closure~\ref{def:closure} of $A$ and the closure~\ref{def:closure} of the complement of $A$. That is, $$\partial A = \overline A \cap \overline{(X\setminus A)}.$$
\end{definition}

\begin{definition}[boundary points of a manifold]\label{def:boundary-points-of-a-manifold}
Let $M$ be a manifold with boundary~\ref{def:topological-m-dimensional-manifold-with-boundary}. The \emph{boundary points} of $M$ are the points $p\in M$ for which there exists a chart~\ref{def:chart} $(U,\phi)$ such that $\phi(p) \in \partial H^m$, the boundary of the closed half-space.
\end{definition}

\begin{definition}[bounded]\label{def:bounded}
A subset $E$ of a metric space~\ref{def:metric-space} $(X,\rho)$ is \emph{bounded} if it has finite diamter~\ref{def:diameter-of-a-set}.
\end{definition}

\begin{definition}[Cantor set]\label{def:cantor-set}
The \emph{Cantor set} $\mathcal C$ is constructed from the unit interval~\ref{def:interval} as follows:

Let $C_n = \frac{C_{n-1}}{3} \cup \left(\frac{2}{3} +\frac{C_{n-1}}{3}\right)$ with $C_0 = [0,1]$. Then $$\mathcal C = \bigcap_{n\in\mathbb N} C_n.$$

A set $E\subset \mathbb R$ is called a \emph{Cantor set} if it is compact~\ref{def:compact}, has empty interior~\ref{def:interior}, and has no isolated points~\ref{def:isolated-point}.
\end{definition}

\begin{definition}[category of topological spaces]\label{def:category-of-topological-spaces}
The \emph{category of topological spaces} is the category~\ref{def:category} whose objects are topological spaces~\ref{def:topological-space} and whose morphisms are continuous~\ref{def:continuous} functions.
\end{definition}

\begin{definition}[Cauchy sequence]\label{def:cauchy-sequence}
A sequence $\{x_n\}_{n\in\mathbb N}$ in a metric space~\ref{def:metric-space} $(X,\rho)$ is \emph{Cauchy} if for all $\varepsilon > 0$, there exists $n_0 > 0$ such that $\rho(x_n, x_m) < \varepsilon$ for all $m,n \geq n_0$
\end{definition}

\begin{definition}[centroid]\label{def:centroid}
Let $S \subset\mathbb R^n$ be rectifiable~\ref{def:rectifiable} with nonzero volume~\ref{def:volume-of-a-rectifiable-set}. The \emph{centroid} of $S$ is the point $c(S) \in \mathbb R^n$ whose $k$th coordinate is given by the integral~\ref{def:riemann-integrable} $$c_k(S) = \frac{1}{\text{vol}(S)} \int_S \pi_k,$$ where $\pi_k:\mathbb R^n \to\mathbb R$ is such that $(x_1,\dots,x_n)\mapsto x_k$.
\end{definition}

\begin{definition}[characteristic map of open simplex]\label{def:characteristic-map-of-open-simplex}
Let $X$ be a ∆-complex. An open simplex $e_\alpha^n$ of dimension $n$ in $X$ has a canonical map $\sigma_\alpha: \Delta^n \to X$ that restricts to a homeomorphism~\ref{def:homeomorphism} from the interior~\ref{def:interior} of the standard n-simplex $\Delta^n$ onto $e_\alpha^n$.
\end{definition}

\begin{definition}[chart]\label{def:chart}
A \emph{chart} or \emph{coordinate patch} on an $m$-dimensional manifold~\ref{def:manifold} $M$ is an open~\ref{def:open} set $U\subset M$ together with a homeomorphism~\ref{def:homeomorphism} $\phi: U\to \mathbb R^m$ onto an open~\ref{def:open} subset of the closed half-space~\ref{def:closed-half-space} $H^m$.
\end{definition}

\begin{definition}[class]\label{def:class}
Let $U \subset \mathbb R^n$ be open~\ref{def:open}. A function $f: U\to \mathbb R^m$ is of \emph{class} $C^k$ if its derivatives~\ref{def:differentiable} $Df, D^2f,\dots, D^kf$ all exist and are continuous~\ref{def:continuous} on $U$.
\end{definition}

\begin{definition}[closed]\label{def:closed}
Let $X$ be a topological space~\ref{def:topological-space}. The subset $A \subset X$ is \emph{closed} if its complement $X\setminus A$ is open~\ref{def:open}. Equivalently, if $X$ is a metric space~\ref{def:metric-space}, then a subset $A \subset X$ is closed~\ref{def:closed} if and only if it contains all of its limit points~\ref{def:limit-point}.
\end{definition}

\begin{definition}[closed half-space]\label{def:closed-half-space}
The \emph{closed~\ref{def:closed} half-space} $H^n \subset\mathbb R^n$ is $$H^n = \{(x_1,\dots,x_n)\in\mathbb R^n\mid x_i \in \mathbb R, x_n \geq 0\}.$$
\end{definition}

\begin{definition}[closure]\label{def:closure}
Let $X$ be a topological space~\ref{def:topological-space}. The \emph{closure} of the set $A \subset X$ is the smallest closed~\ref{def:closed} subset of $X$ containing $A$, or the union of $A$ and its limit points~\ref{def:limit-point}.
\end{definition}

\begin{definition}[comeager]\label{def:comeager}
Let $X$ be a topological space~\ref{def:topological-space}. A set $A\subset X$ is \emph{comeager} if its complement is meager~\ref{def:meager}.
\end{definition}

\begin{definition}[compact]\label{def:compact}
A topological space~\ref{def:topological-space} $X$ is \emph{compact} if every open cover~\ref{def:open-cover} of $X$ has a finite subcover.
\end{definition}

\begin{definition}[compact Lie group]\label{def:compact-lie-group}
A Lie group~\ref{def:lie-group} $G\subseteq \text{GL}(n,\mathbb C)$ is \emph{compact} if it is compact~\ref{def:compact} in the usual sense when considered as a subset of $M_n(\mathbb C) \cong \mathbb R^{2n^2}.$
\end{definition}

\begin{definition}[compact linear transformation]\label{def:compact-linear-transformation}
A linear transformation~\ref{def:linear-transformation} $T:V\to V$ from a normed linear space~\ref{def:normed-linear-space} $V$ to itself is \emph{compact} if the closure~\ref{def:closure} of the image~\ref{def:image} of the closed~\ref{def:closed} unit ball $B = \{v \in V \mid \vert \vert v\vert \vert  \leq 1\}$ is compact~\ref{def:compact} in the usual sense.
\end{definition}

\begin{definition}[compact support]\label{def:compact-support}
A function $f: X\to\mathbb R$ has \emph{compact support~\ref{def:support}} if the closure~\ref{def:closure} of its support~\ref{def:support} $S$ is compact~\ref{def:compact} as a subset of $X$.
\end{definition}

\begin{definition}[compact symplectic group]\label{def:compact-symplectic-group}
The \emph{compact symplectic group} $\text{Sp}(n)$ is the group~\ref{def:group} intersection of the complex symplectic group~\ref{def:complex-symplectic-group} $\text{Sp}(2n,\mathbb C)$ and the unitary group~\ref{def:unitary-group} $\text{U}(2n)$: $$\text{Sp}(n) = \text{Sp}(2n,\mathbb C)\cap\text{U}(2n).$$ It is a matrix Lie group~\ref{def:matrix-lie-group}.
\end{definition}

\begin{definition}[compactly closed]\label{def:compactly-closed}
A subspace~\ref{def:subspace-topology} $A$ of a topological space~\ref{def:topological-space} $X$ is \emph{compactly closed} if the inverse image $g^{-1}(A)$ is closed~\ref{def:closed} in $K$ for any map $g:K\to X$ from a compact~\ref{def:compact} Hausdorff~\ref{def:hausdorff} space~\ref{def:topological-space}  $K$ into $X$.
\end{definition}

\begin{definition}[complete metric space]\label{def:complete-metric-space}
A metric space~\ref{def:metric-space} $(X,\rho)$ is \emph{complete} if every Cauchy sequence~\ref{def:cauchy-sequence} is also convergent~\ref{def:sequence-convergence}.
\end{definition}

\begin{definition}[complex exponential]\label{def:complex-exponential}
The \emph{complex exponential} is given by the power series~\ref{def:power-series} $$e^z = \sum_{n=0}^\infty \frac{z^n}{n!}$$ which has infinite radius of convergence~\ref{def:radius-of-convergence}.
\end{definition}

\begin{definition}[complex integral]\label{def:complex-integral}
A complex~\ref{def:complex-numbers}-valued function $f:[a,b]\to\mathbb C$ is \emph{integrable} if the component functions $f(t) = u(t)+iv(t)$ are Riemann integrable~\ref{def:riemann-integrable} as real-valued functions. The \emph{complex integral} is then given by $$\int_a^b u(t)\text dt + i\int_a^b v(t)\text dt$$
\end{definition}

\begin{definition}[complex Lie group]\label{def:complex-lie-group}
A matrix Lie group~\ref{def:matrix-lie-group} $\mathfrak g$ is \emph{complex} if it is a complex subspace~\ref{def:vector-subspace} of $M_n(\mathbb C)$. That is, if $iX\in \mathfrak g$ for all $X\in \mathfrak g$.
\end{definition}

\begin{definition}[complex line integral]\label{def:complex-line-integral}
Let $U$ be an open~\ref{def:open} subset of $\mathbb C$ and $f:U \to  \mathbb C$ a continuous~\ref{def:continuous} function and $\gamma:[a,b]\to U$ a piecewise continuously~\ref{def:class} differentiable~\ref{def:differentiable} path~\ref{def:path} contained in $U$. The \emph{complex line integral} of $f$ over $\gamma$ is the complex integral~\ref{def:complex-integral} $$\int_\gamma f(z)\text dz = \int_a^b f(\gamma(t))\gamma'(t)\text dt$$

todo expand definition in measure
\end{definition}

\begin{definition}[complex orthogonal group]\label{def:complex-orthogonal-group}
The \emph{complex orthogonal group} $\text{O}(n,\mathbb C)$ is the set of $n\times n$ complex matrices~\ref{def:matrix} whose transposes~\ref{def:matrix-transpose} are equal to their inverses~\ref{def:inverse-element}. 

This definition mirrors the definition of the real orthogonal group~\ref{def:orthogonal-group}, but instead we consider complex-valued matrices.

It is a matrix Lie group~\ref{def:matrix-lie-group} because it is a subgroup~\ref{def:subgroup} of the general linear group~\ref{def:general-linear-group} $\text{GL}(n,\mathbb C)$.
\end{definition}

\begin{definition}[complex symplectic group]\label{def:complex-symplectic-group}
The \emph{complex symplectic group} $\text{Sp}(2n,\mathbb C)$ is the group~\ref{def:group} of $2n\times 2n$ complex-valued matrices~\ref{def:matrix} that preserve the form $\omega$ as in the definition of the real symplectic group~\ref{def:real-symplectic-group}.

This group~\ref{def:group} is a matrix Lie group~\ref{def:matrix-lie-group} because it is a subgroup~\ref{def:subgroup} of the general linear group~\ref{def:general-linear-group} $\text{GL}(n,\mathbb C)$.
\end{definition}

\begin{definition}[component function]\label{def:component-function}
Let $U \subset \mathbb R^n$ be open~\ref{def:open}. A function $f:\mathbb U \to \mathbb R^n$ may be written as an $m$-tuple of functions $f_i:U\to \mathbb R$ so that $$f=(f_1,\dots,f_m).$$
\end{definition}

\begin{definition}[connected]\label{def:connected}
A topological space~\ref{def:topological-space} is \emph{connected} if it cannot be expressed as the disjoint union of two nonempty open~\ref{def:open} subsets.
\end{definition}

\begin{definition}[continuous]\label{def:continuous}
Let $f: X\to Y$ be a function between topological spaces~\ref{def:topological-space} $X$ and $Y$. Then $f$ is \emph{continuous} if and only if the preimage of an open~\ref{def:open} set in $Y$ is open~\ref{def:open} in $X$.

A function $f: X\to Y$ between metric spaces~\ref{def:metric-space} $(X,\rho)$ and $(Y,\sigma)$ is \emph{continuous} at $x\in X$ if and only if for every $\varepsilon > 0$, there exists $\delta> 0$ such that $\sigma(f(x),f(x')) < \varepsilon$ whenever $\rho(x,x') < \delta$.

A function $f: X\to Y$ between metric spaces~\ref{def:metric-space} $(X,\rho)$ and $(Y,\sigma)$ is \emph{continuous} at $x\in X$ if and only if for any sequence $\{x_n\}_{n\in\mathbb N}$ in $X$ converging~\ref{def:sequence-convergence} to $x$, $\{f(x_n)\}_{n\in \mathbb N}$ converges to $f(x)$.

These three definitions are all equivalent in their respective contexts.
\end{definition}

\begin{definition}[continuous group representation]\label{def:continuous-group-representation}
A representation~\ref{def:group-representation} of a topological group~\ref{def:topological-group} $G$ (written $\rho:G\to GL(V)$ for $GL(V)$  the general linear group~\ref{def:general-linear-group} of some vector space~\ref{def:vector-space} $V$) is called \emph{continuous} if it is of finite dimension of dimension~\ref{def:dimension-of-vector-space} and if $\rho$ is continuous~\ref{def:continuous} in the usual sense as a function between topological spaces~\ref{def:topological-space}.
\end{definition}

\begin{definition}[contractible]\label{def:contractible}
A topological space~\ref{def:topological-space} is \emph{contractible} if it is homotopy equivalent to a point.
\end{definition}

\begin{definition}[contracting sequence of sets in a metric space]\label{def:contracting-sequence-of-sets-in-a-metric-space}
A descending sequence $\{E_n\}_{n\in\mathbb N}$ of nonempty subsets of a metric space~\ref{def:metric-space} $X$ is \emph{contracting} if the limit of the diameters~\ref{def:diameter-of-a-set} goes to zero. That is, $$\lim_{n\to \infty} \text{diam}(E_n) = 0.$$
\end{definition}

\begin{definition}[contraction mapping]\label{def:contraction-mapping}
Let $(X,\rho)$ be a metric space~\ref{def:metric-space}. Then a function $f:X\to X$ is a \emph{contraction mapping} if there exists $r \in [0,1)$ such that for all $x,y \in x$, $$\rho(f(x),f(y)) \leq r\rho(x,y).$$
\end{definition}

\begin{definition}[convergence in measure]\label{def:convergence-in-measure}
Let $(X,\Sigma,\mu)$ be a measure space~\ref{def:measure-space} and let $\{f_n\}_{n\in\mathbb N}$ be a sequence of functions $f_n: X\to \mathbb R$. The sequence \emph{converges in measure} to $f$ if for all $\varepsilon > 0$, the measure~\ref{def:measure-space} of the set $\mu(\{x\in X\mid \vert f_n(x) -f(x)\vert >\varepsilon\})$ converges~\ref{def:sequence-convergence} to zero as $n$ goes to infinity.
\end{definition}

\begin{definition}[critical point]\label{def:critical-point}
Let $U\subset\mathbb R^n$ be an open~\ref{def:open} subset and let $f:U\to\mathbb R$ be differentiable~\ref{def:differentiable} on $U$. The point $a \in U$ is a \emph{critical point} of $f$  if $Df(a) =0$.
\end{definition}

\begin{definition}[critical value]\label{def:critical-value}
Let $U\subset\mathbb R^n$ be an open~\ref{def:open} set and let $f:U\to\mathbb R^n$ be of class~\ref{def:class} $C^1$. A point $a\in \mathbb R^n$ is a \emph{critical value} of $f$ if there exists a critical point~\ref{def:critical-point} $c \in U$ such that $f(c) =a$.
\end{definition}

\begin{definition}[dense]\label{def:dense}
Let $X$ be a topological space~\ref{def:topological-space}. A set $A \subset X$ is \emph{dense} in $X$ if every point in $x$ is either in $A$ or is a limit point~\ref{def:limit-point} of $A$. Equivalently, $A$ is dense in $X$ if its closure~\ref{def:closure} is all of $X$.
\end{definition}

\begin{definition}[determined]\label{def:determined}
A game~\ref{def:game} is \emph{determined} if one player has a winning strategy.
\end{definition}

\begin{definition}[diameter of a set]\label{def:diameter-of-a-set}
The \emph{diameter} of a nonempty subset $E$ of a metric space~\ref{def:metric-space} $(X,\rho)$ is given by the supremum~\ref{def:supremum} $$\text{diam}(E) = \sup\{\rho(x,y)\mid x,y \in E\}.$$
\end{definition}

\begin{definition}[diffeomorphism]\label{def:diffeomorphism}
Let $U,V\subset\mathbb R^n$ be open~\ref{def:open}. A \emph{class~\ref{def:class} $C^k$ diffeomorphism} is a bijective~\ref{def:bijective} $C^k$ function $g:U\to V$ such that the inverse $g^{-1}$ is also of class~\ref{def:class} $C^k$.
\end{definition}

\begin{definition}[differentiability in closed half-space]\label{def:differentiability-in-closed-half-space}
Let $U\subset\mathbb H^n$ be open~\ref{def:open} in the subspace topology~\ref{def:subspace-topology} on $H^n$ considered as a subset of $\mathbb R^n$. Then a function $f:U\to\mathbb R^m$ is \emph{differentiable} if there is an open set $V \subset\mathbb R^n$ such that $U = V\cap H^n$ onto which we may extend $f$ to a differentiable~\ref{def:differentiable} function $f_V:V\to\mathbb R^n$.
\end{definition}

\begin{definition}[differentiable]\label{def:differentiable}
Let $U \subset \mathbb R^n$ be open~\ref{def:open}. A function $f: U\to \mathbb R^m$ is \emph{differentiable} at $a \in U$ if there exists a linear transformation~\ref{def:linear-transformation} $A:\mathbb R^n \to \mathbb R^m$ such that $$\lim_{h\to 0} \frac{f(a+h)-f(a)-A(h)}{\vert \vert h\vert \vert } = 0.$$ The transformation $A$ is denoted $Df(a)$ and is called the \emph{total derivative} of $f$ at $a$.

Equivalently, let $V$ and $V$' be vector spaces~\ref{def:vector-space} with norms $\vert \cdot\vert $ and $\vert \cdot\vert '$, respectively. A function $f:U\to V'$ from $U$ an open~\ref{def:open} subset of $V$ is \emph{differentiable} at $x\in U$ if there exists a linear transformation~\ref{def:linear-transformation} $d_x f :V\to V'$ called the \emph{differential} of $f$ at $x$ and a small neighborhood~\ref{def:neighborhood} $U_x\subseteq U$ of $x$ such that for all $v\in U$, we have that $$f(x+v) = f(x) + d_x(f) + o(v)$$ where $o: U_x\to V'$ is a map such that $$\lim_{\vert v\vert \to 0} \frac{\vert o(v)\vert '}{\vert v\vert } = 0.$$
\end{definition}

\begin{definition}[differential k-form]\label{def:differential-k-form}
Let $U\subset\mathbb R^n$ be an open~\ref{def:open} set. A \emph{differential $k$-form} on $U$ is a map to the $k$th exterior product~\ref{def:exterior-power-of-a-vector-space} of the dual space~\ref{def:dual-space} of $\mathbb R^n$: $$\phi:U\to\wedge^k((\mathbb R^n)^*).$$

A $0$-form is just a function. Wedging a $0$-form with something amounts to multiplication by the function.

Let $U\subset\mathbb R^n$ be an open~\ref{def:open} set. Then we write $\mathcal A^k(U)$ for the vector space~\ref{def:vector-space} of k-forms~\ref{def:differential-k-form} on $U$.
\end{definition}

\begin{definition}[directional derivative]\label{def:directional-derivative}
Let $U \subset \mathbb R^n$ be open~\ref{def:open}. Let $v \in \mathbb R^n$. The \emph{directional derivative} of a function $f:U\to\mathbb R^m$ in the direction of $v$ at a point $a \in U$ is $$D_vf(a) = \frac{\partial f}{\partial v}(a) = \lim_{h\to 0}\frac{f(a+hv)-f(a)}{h}$$ if this limit exists.
\end{definition}

\begin{definition}[discrete topology]\label{def:discrete-topology}
Any set $X$ can be equipped with the \emph{discrete} topology~\ref{def:topological-space}, namely the topology in which every set is both open~\ref{def:open} and closed~\ref{def:closed}.
\end{definition}

\begin{definition}[distance from a point to a set]\label{def:distance-from-a-point-to-a-set}
Let $X$ be a metric space~\ref{def:metric-space} and $K \subset X$. For $x \in X$, the \emph{distance} from $x$ to $K$ is given by the infimum~\ref{def:infimum} $$\text{dist}(x,K) = \inf_{y\in K}(\rho(x, y)).$$
\end{definition}

\begin{definition}[embedded m-dimensional manifold]\label{def:embedded-m-dimensional-manifold}
Let $M\subset\mathbb R^n$. Then $M$ is an \emph{embedded $m$-dimensional manifold} if for all $x\in M$, there exists a neighborhood~\ref{def:neighborhood} $U \subset\mathbb R^n$ of $x$ and a smooth~\ref{def:smooth} function $f:U\to\mathbb R^{n-m}$  such that $M\cap U = f^{-1}(0)$ and $Df(y):\mathbb R^n\to\mathbb R^{n-m}$ is surjective~\ref{def:surjective} for all $y \in U$.
\end{definition}

\begin{definition}[embedded m-dimensional manifold with boundary]\label{def:embedded-m-dimensional-manifold-with-boundary}
Let $m< n$. Then an \emph{$m$-dimensional smooth manifold embedded in $\mathbb R^n$ with boundary} is a subset $M \subset\mathbb R^n$ such that for all points $p \in M$, there exists an open~\ref{def:open} set $U \subset H^k$ the closed half-space~\ref{def:closed-half-space} and an open~\ref{def:subspace-topology} neighborhood $V \subset M$ of $p$ and a map $\phi: U \to V$ such that 
1. $\phi$ is a homeomorphism~\ref{def:homeomorphism}
2. $\phi$ is smooth~\ref{def:smooth} in $H^k$ ( todo: what is this exactly?)
3. the total derivative~\ref{def:differentiable} $D\phi(x)$ has rank $m$ for all $x \in U$.
\end{definition}

\begin{definition}[embedding]\label{def:embedding}
Let $W_1,W_2 \subset \mathbb R^n$ be open~\ref{def:open}. A function $g: W_1\to W_2$ is a \emph{$C^k$-embedding} if $g$ is injective~\ref{def:injective}, $g$ is of class~\ref{def:class} $C^k$, and the total derivative~\ref{def:differentiable} $Dg(x)$ is invertible for all $x \in W_1$.
\end{definition}

\begin{definition}[equivalent metrics]\label{def:equivalent-metrics}
Let $X$ be a metric space~\ref{def:metric-space}. Two metrics~\ref{def:metric-space} $\rho_1$, $\rho_2$ are \emph{equivalent} if there exists $c>0$ such that $$\frac{1}{c}\rho_1(x,y) \leq \rho_2(x,y) \leq c\rho_1(x,y).$$
\end{definition}

\begin{definition}[exterior]\label{def:exterior}
Let $X$ be a topological space~\ref{def:topological-space}. For any $A\subset X$, the \emph{exterior} $\text{ext}(A)$ of $A$ is the complement of the closure~\ref{def:closure} of $A$.
\end{definition}

\begin{definition}[fundamental group]\label{def:fundamental-group}
Let $X$ be a topological space~\ref{def:topological-space} and let $x \in X$. Denote by $\pi_1(X,x)$ the set of equivalence classes~\ref{def:equivalence-class} of loops~\ref{def:loop} at $x$ under homotopy equivalence of paths. This is called the \emph{fundamental group of $X$ at $x$}.
\end{definition}

\begin{definition}[G-space]\label{def:g-space}
Let $G$ be a group~\ref{def:group} and let $X$ be a topological space~\ref{def:topological-space}. Then $X$ is a \emph{$G$-space} if it is a G-set~\ref{def:group-action} and the group action~\ref{def:group-action} $G\times X\to X$ is continuous~\ref{def:continuous}.
\end{definition}

\begin{definition}[Gale-Stewart game]\label{def:gale-stewart-game}
A \emph{Gale-Stewart game} is a two-player game~\ref{def:game} in which both players have perfect information, make an infinite sequence of moves, and there are no draws.
\end{definition}

\begin{definition}[game]\label{def:game}
A \emph{game} is something which may be represented as a normal-form game~\ref{def:normal-form-game} or an extended-form game.
\end{definition}

\begin{definition}[generalized orthogonal group]\label{def:generalized-orthogonal-group}
The \emph{generalized orthogonal group~\ref{def:orthogonal-group}} $\text{O}(n,k)$ is the group of $(n+k)\times (n+k)$ real-valued matrices~\ref{def:matrix} that preserve the symmetric~\ref{def:symmetric} bilinear~\ref{def:multilinear} form $$[x,y]_{n,k} = x_1y_1 + \cdots + x_ny_n - x_{n+1}y_{n+1} - \cdots - x_{n+k}y_{n+k}.$$That is, $A \in \text{O}(n,k)$ if and only if $[x,y]_{n,k} = [Ax,Ay]_{n,k}$ for all $x,y \in \mathbb R^n$. 

This group~\ref{def:group} is a matrix Lie group~\ref{def:matrix-lie-group} because it is a subgroup~\ref{def:subgroup} of the general linear group~\ref{def:general-linear-group} $\text{GL}(n,\mathbb C)$.
\end{definition}

\begin{definition}[generate a topology]\label{def:generate-a-topology}
The topology~\ref{def:topological-space} \emph{generated}  by a basis~\ref{def:topological-basis} $\mathcal B$ is the collection $\tau$ of subsets $U$ such that for every such that for every $x\in U$, there exists a $B\in \mathcal B$ such that $x\in B\subset U$.

Equivalently, $U\in \tau$ if and only if it is a union of sets in $\beta$.
\end{definition}

\begin{definition}[gradient]\label{def:gradient}
Let $U\subset\mathbb R^n$ be an open~\ref{def:open} subset and $f:U\to \mathbb R$ be differentiable~\ref{def:differentiable} on $U$. The \emph{gradient} of $f$ at $a \in U$ is the vector of the partial derivatives~\ref{def:partial-derivative} of $f$: $$\nabla f(a) = \left(\frac{\partial f}{\partial x_1}(a),\cdots, \frac{\partial f}{\partial x_n}(a) \right).$$
\end{definition}

\begin{definition}[graph of a function]\label{def:graph-of-a-function}
Let $U \subset \mathbb R^n$ be open~\ref{def:open}. The set $U$ is the \emph{graph} of a some function if for all $x \in U$, there exists unique $y \in \mathbb R^m$ such that $(x,y) \in A$.

Alternatively, let $f:A\to B$ be a set function. The set $$\Gamma = \{(a,b) \in A\times B \mid b=f(a)\}$$ is the \emph{graph} of $f$.
\end{definition}

\begin{definition}[Haar integral]\label{def:haar-integral}
We define the \emph{Haar integral} over a topological space~\ref{def:topological-space} $X$ follows: 

Any Borel~\ref{def:borel-σ-algebra} measure~\ref{def:measure-space} $\mu$ on $X$ gives a linear~\ref{def:linear-transformation} functional~\ref{def:functional} $$\int:C_c(X) \to \mathbb C, f\mapsto \int_X(f) := \int_X f(x)\text d\mu(x)$$ where $C_c(X)$ is the collection of continuous~\ref{def:continuous} complex-valued functions on $X$ with compact support~\ref{def:compact-support} and the integral~\ref{def:lebesgue-integral} is defined according to that measure.

The Haar integral is this $\mathbb C$-linear function $\int:C_c(X) \to \mathbb C$  such that
1. for all $f\in C_c(X)$, if $f(x) \geq 0$ for all $x\in X$ then $\int_X f \geq 0$. Moreover, $\int_X f= 0$ implies that $f=0$.  (when $f(x)$ is real)
2. For any compact~\ref{def:compact} $K\subset X$, there exists some $c_K\geq 0$ such that for all continuous~\ref{def:continuous} functions with support~\ref{def:support} in $K$, $\left\vert \int_X f\right\vert  \leq c_K\cdot \max\limits_{x\in K} \vert f(x)\vert $.

If we want the integral to be invariant~\ref{def:g-invariant-function} with respect to the action of a group~\ref{def:group-action}, we choose the measure~\ref{def:measure-space} to be the Haar measure.

Also keep in mind that this is in fact a special case of the Lebesgue integral~\ref{def:lebesgue-integral}!
\end{definition}

\begin{definition}[Hamel basis]\label{def:hamel-basis}
An orthonormal sequence~\ref{def:orthonormal} of vectors $\{v_i\}_{i\in\mathbb N}$ in a Hilbert space~\ref{def:hilbert-space} $H$ is an \emph{orthonormal (Hamel) basis} for $H$ if $$h = \sum_{i\in\mathbb N} \langle v_i, h\rangle v_i$$ for all $h \in H$.
\end{definition}

\begin{definition}[harmonic function]\label{def:harmonic-function}
Let $D\subset\mathbb R^n$ be connected~\ref{def:connected} and open~\ref{def:open}. A function $f:D\to \mathbb R$ is \emph{harmonic} if for all $z\in D$ and every $\varepsilon > 0$ with $\text{dist}(z,\partial D) > \varepsilon$, then $f(z) = MV(f,z,\varepsilon)$ where $\partial D$ denotes the boundary~\ref{def:boundary} of $D$ and $MV$ denotes the circular mean value. 

Equivalently, a function $f$ is harmonic if its Laplacian~\ref{def:laplacian} is 0.
\end{definition}

\begin{definition}[harmonic polynomial]\label{def:harmonic-polynomial}
A \emph{harmonic polynomial} is a harmonic function~\ref{def:harmonic-function} that is also a polynomial~\ref{def:polynomial-ring}.
\end{definition}

\begin{definition}[Hausdorff]\label{def:hausdorff}
A topological space~\ref{def:topological-space} $X$ is \emph{Hausdorff} if for every pair of points $x,y\in X$, there exist disjoint open~\ref{def:open} sets $U,V\subset X$ such that $x\in U$ and $y\in V$.
\end{definition}

\begin{definition}[Hessian matrix]\label{def:hessian-matrix}
Let $U\subset\mathbb R^n$ be an open~\ref{def:open} subset and let $f:U\to\mathbb R$ be of class~\ref{def:class} at least $C^2$. The \emph{Hessian matrix} of $f$ at the point $a \in U$ is the matrix 
$$H = \begin{bmatrix} \frac{\partial^2 f}{\partial x_i\partial x_j} \end{bmatrix}_{i,j=1}^n$$ of the mixed second partial derivatives~\ref{def:partial-derivative} of $f$.
\end{definition}

\begin{definition}[Hilbert space]\label{def:hilbert-space}
A \emph{Hilbert space} $(H, \langle\cdot,\cdot\rangle)$ is an inner product space~\ref{def:inner-product-space} that is complete~\ref{def:complete-metric-space} as a metric space~\ref{def:metric-space} with metric induced by the norm from the inner product.
\end{definition}

\begin{definition}[homeomorphism]\label{def:homeomorphism}
A \emph{homeomorphism} between two topological spaces~\ref{def:topological-space} $X$ and $Y$ is a function $f: X\to Y$ such that 
1. $f$ is continuous~\ref{def:continuous}.
2. $f$ is bijective~\ref{def:bijective}.
3. $f^{-1}:Y\to X$ is continuous~\ref{def:continuous}.

Two topological spaces~\ref{def:topological-space} $X$ and $Y$ are \emph{homeomorphic} if there exists a homeomorphism between them.
\end{definition}

\begin{definition}[homomorphism of Lie groups]\label{def:homomorphism-of-lie-groups}
Let $G$ and $H$ be Lie groups~\ref{def:lie-group}. A map $\Phi$ from $G$ to $H$ is a \emph{Lie group homomorphism} if it is a continuous~\ref{def:continuous} group homomorphism~\ref{def:group-homomorphism}.
\end{definition}

\begin{definition}[identity component]\label{def:identity-component}
Let $G$ be a matrix Lie group~\ref{def:matrix-lie-group}. The \emph{identity component} $G_0$ of $G$ is the set of $A\in G$ for which there exists a continuous~\ref{def:continuous} path $\gamma:[0,1]\to G$ such that $\gamma(0) = I$ and $\gamma(1) = A$. Because connectedness and path-connectedness are equivalent, in this case, $G_0$ amounts to the connected component~\ref{def:connected-component} of $G$ that contains the identity matrix~\ref{def:identity-matrix} $I$.
\end{definition}

\begin{definition}[induced homomorphism on fundamental group]\label{def:induced-homomorphism-on-fundamental-group}
Let $p:X\to Y$ be a continuous~\ref{def:continuous} map between topological spaces~\ref{def:topological-space}. Define $p_*:\pi_1(X,x) \to \pi_1(Y,p(x))$ to be the map between the fundamental groups~\ref{def:fundamental-group} of $X$ at $x$ and $Y$ at $p(x)$ given by $p_*[f] = [p\circ f]$. This map is a homomorphism~\ref{def:group-homomorphism} as $p_*[f\cdot g] = [p\circ (f\cdot g)] = p_*[f]p_*[g]$.
\end{definition}

\begin{definition}[interior]\label{def:interior}
Let $X$ be a topological space~\ref{def:topological-space}. For any $A\subset X$, the \emph{interior} $\text{int}(A)$ of $A$ is the largest open~\ref{def:open} set contained in $A$.
\end{definition}

\begin{definition}[interior points of a manifold]\label{def:interior-points-of-a-manifold}
Let $M$ be a manifold with boundary~\ref{def:topological-m-dimensional-manifold-with-boundary}. The \emph{interior points} of $M$ are the points $p\in M$ for which there exists a chart~\ref{def:chart} $(U,\phi)$ such that $\phi(p) \in H^m_+$, the interior~\ref{def:interior} of the closed half-space~\ref{def:closed-half-space}.
\end{definition}

\begin{definition}[isolated point]\label{def:isolated-point}
A point $x$ is an \emph{isolated point} of a subset $S$ of a topological space~\ref{def:topological-space} $X$ if it is an element of $S$ but there exists a neighborhood~\ref{def:neighborhood} of $x$ containing no other points of $S$.
\end{definition}

\begin{definition}[isolated singularity]\label{def:isolated-singularity}
Let $f:U\setminus\{a\}\to \mathbb C$ be holomorphic~\ref{def:holomorphic} on the open~\ref{def:open} set $C\setminus \{a\}$ except at $a$. Then $a$ is an \emph{isolated singularity} of $f$ and can be one of the following:
1. $a$ is a \emph{removable singularity} if there exists a holomorphic~\ref{def:holomorphic} function $g:U\to \mathbb C$ such that $f(z) = g(z)$ for all $z\in U\setminus a$. Such a function $g$ is a continuous~\ref{def:continuous} and holomorphic~\ref{def:holomorphic} extension of $f$ over $a$ as in Riemann's theorem on removable singularities.
2. $a$ can be a \emph{pole} if there exists a holomorphic~\ref{def:holomorphic} function $g:U\to\mathbb C$ with $g(a)\neq 0$ and $n\in\mathbb N$ such that $f(z) = \frac{g(z)}{(z-a)^n}$ for all $z\in U\setminus \{a\}$. The number $n$ is the \emph{order} of the pole.
3. $a$ can be an \emph{essential singularity} of $f$ if it is neither a removable singularity nor a pole. This occurs if and only if the Laurent series~\ref{def:laurent-series} for $f$ has infinitely many nonzero terms in negative degree. That is, if the principal part~\ref{def:principal-part} is infinite.
\end{definition}

\begin{definition}[isometry]\label{def:isometry}
A function $f: X\to Y$ between metric spaces~\ref{def:metric-space} $(X, \rho)$ and $(Y,\sigma)$ is an \emph{isometry} if for all $a,b \in X$, $$\sigma(f(a), f(b)) = \rho(a,b).$$ Two metric spaces~\ref{def:metric-space} are \emph{isometric} if there exists an isometry between them.
\end{definition}

\begin{definition}[isomorphism of Lie groups]\label{def:isomorphism-of-lie-groups}
Let $G$ and $H$ be Lie groups~\ref{def:lie-group}. A map $\phi$ from $G$ to $H$ is an \emph{isomorphism of Lie groups} if it is a homeomorphic~\ref{def:homeomorphism} group homomorphism~\ref{def:group-homomorphism}. Equivalently, $\Phi$ is an \emph{isomorphism of Lie groups} if it is a bijective~\ref{def:bijective} homomorphism of Lie groups~\ref{def:homomorphism-of-lie-groups} with continuous~\ref{def:continuous} inverse~\ref{def:inverse-function}.
\end{definition}

\begin{definition}[Jacobian determinant]\label{def:jacobian-determinant}
Let $U\subset\mathbb R^n$ be an open~\ref{def:open} subset and let $g: U \to\mathbb R^n$ be  differentiable~\ref{def:differentiable}. Then \emph{Jacobian determinant~\ref{def:determinant}} of $g$ is the function $Jg:U\to\mathbb R$ given by the absolute value of thedeterminant~\ref{def:determinant} of the Jacobian matrix~\ref{def:jacobian-matrix} of $g$: $$Jg(x)= \vert \text{det}(Dg(x))\vert .$$
\end{definition}

\begin{definition}[Jacobian matrix]\label{def:jacobian-matrix}
Let $U \subset \mathbb R^n$ be open~\ref{def:open} and let $f:U\to\mathbb R^n$ be differentiable~\ref{def:differentiable}. Then its total derivative~\ref{def:differentiable} $Df(a)$ at a point $a\in U$ is given by the $m\times n$ matrix $$\left[ \frac{\partial f_j}{\partial x_i}(a)\right]_{i=1,j=1}^{n,m}$$ where $f = (f_1,\dots, f_m)$ are the component functions~\ref{def:component-function} of $f$. This matrix is called the \emph{Jacobian} matrix of $f$ at $a$.
\end{definition}

\begin{definition}[L2 inner product on group functions]\label{def:l2-inner-product-on-group-functions}
Let $G$ be a compact~\ref{def:compact} topological group~\ref{def:topological-group} and let $dg$ be a Haar measure on $G$. The \emph{$L^2$-inner product} on $C(G)$, the space of continuous~\ref{def:continuous} functions $G\to \mathbb C$, is a Hermitian inner product~\ref{def:hermitian-inner-product} given by $$\langle f_1,f_2\rangle = \frac{1}{\text{vol}(G)}\int_G f_1(g)\overline{f_2(g)}\text dg$$ for any $f_1,f_2\in C(G)$. The integral is the Haar integral~\ref{def:haar-integral} and the volume~\ref{def:volume-of-compact-topological-space} of $G$ is the integral of the constant function $1$ over $G$.
\end{definition}

\begin{definition}[Laplacian]\label{def:laplacian}
Let $D\subset\mathbb R^2$ be connected~\ref{def:connected} and open~\ref{def:open} and let $f:D\to \mathbb R$ be of class~\ref{def:class} $C^2$. Then the \emph{Laplacian} of $f$ is given by the sums of the second partial derivatives~\ref{def:partial-derivative} $$\Delta f(z) = f_{xx}(z) + f_{yy}(z).$$
\end{definition}

\begin{definition}[Laurent series]\label{def:laurent-series}
Let $f:A\to \mathbb C$ be a function holomorphic~\ref{def:holomorphic} on the annulus $A$. Then $f(z)$ can be written as the \emph{Laurent series} $$f(z) = \sum_{n=-\infty}^\infty a_n (z-c)^n$$ for $a_n,c\in \mathbb C$ and $a_n$ given by $$a_n = \frac{1}{2\pi i} \int_\gamma \frac{f(z)}{(z-c)^{n+1}}\text dz$$ for $\gamma$ a positively-oriented piecewise continuously differentiable~\ref{def:class} loop~\ref{def:loop} in $A$.
\end{definition}

\begin{definition}[Lebesgue measurable function]\label{def:lebesgue-measurable-function}
A function $f: \mathbb R^n \to \mathbb R$ is \emph{Lebesgue measurable} if for every open~\ref{def:open} $G \in \mathbb R$, the preimage $f^{-1}(G)$ is Lebesgue measurable~\ref{def:lebesgue-measure}.
\end{definition}

\begin{definition}[Lie correspondence]\label{def:lie-correspondence}
Let $G$ be a matrix Lie group~\ref{def:matrix-lie-group}. Then the \emph{Lie algebra~\ref{def:lie-algebra} of $G$} is the set of all matrices~\ref{def:matrix} $X \in M_n(\mathbb C)$ such that the matrix exponential $e^{tX} \in G$ for all $t \in \mathbb R$.
\end{definition}

\begin{definition}[Lie group]\label{def:lie-group}
A \emph{Lie group} is a group~\ref{def:group} that is also a manifold~\ref{def:embedded-m-dimensional-manifold}.
\end{definition}

\begin{definition}[limit point]\label{def:limit-point}
A point $x \in X$ a topological space~\ref{def:topological-space} is a \emph{limit point} of the subset $A\subset X$ if it is the limit of some sequence~\ref{def:sequence-convergence} of elements of $A$.
\end{definition}

\begin{definition}[local maximum]\label{def:local-maximum}
Let $U\subset\mathbb R^n$ be an open~\ref{def:open} subset and let $f:U\to\mathbb R$. The point $a \in U$ is a \emph{local maximum} of $f$  if there is a neighborhood~\ref{def:neighborhood} $V$ of $a$ such that for all $x \in V$, $f(a)\geq f(x)$.
\end{definition}

\begin{definition}[local minimum]\label{def:local-minimum}
Let $U\subset\mathbb R^n$ be an open~\ref{def:open} subset and let $f:U\to\mathbb R$. The point $a \in U$ is a \emph{local minimum} of $f$  if there is a neighborhood~\ref{def:neighborhood} $V$ of $a$ such that for all $x \in V$, $f(a)\leq f(x)$.
\end{definition}

\begin{definition}[locally compact group]\label{def:locally-compact-group}
A topological group~\ref{def:topological-group} $G$ is \emph{locally compact} if the underlying topological space~\ref{def:topological-space} is both locally compact and Hausdorff~\ref{def:hausdorff}.
\end{definition}

\begin{definition}[locally finite cover]\label{def:locally-finite-cover}
An open cover~\ref{def:open-cover} $\{U\}$ of a topological space~\ref{def:topological-space} $B$ is \emph{locally finite} if every point $b\in B$ is in only finitely many $U$.
\end{definition}

\begin{definition}[locally finite graph]\label{def:locally-finite-graph}
A graph $X$ is \emph{locally finite} if each vertex is a boundary~\ref{def:boundary} point of only finitely many eges or if $X$ is a locally compact topological space~\ref{def:topological-space}.
\end{definition}

\begin{definition}[locally finite measure]\label{def:locally-finite-measure}
Let $(X,\mathcal T)$ be a Hausdorff~\ref{def:hausdorff} topological space~\ref{def:topological-space} and let $\Sigma$ be a σ-algebra~\ref{def:σ-algebra}  containing $\mathcal T$ (so that $\Sigma$ is at least as fine as the Borel σ-algebra~\ref{def:borel-σ-algebra} on $X$). A measure~\ref{def:measure-space} $\mu$ on $\Sigma$ is \emph{locally finite} if for all $p \in X$, there exists an open~\ref{def:open} neighborhood $N_p$ of $p$ such that the $\mu$-measure of $N_p$ is finite. That is, $$\forall p \in X, \exists N_p \in T \text{ s.t. } p \in N_p \text{ and } \vert \mu(N_p)\vert  < \infty.$$
\end{definition}

\begin{definition}[locally path-connected]\label{def:locally-path-connected}
A topological space~\ref{def:topological-space} $X$ is \emph{locally path-connected} if for any $x\in X$ and any neighborhood~\ref{def:neighborhood} $U$ of $x$, there is a smaller neighborhood~\ref{def:neighborhood} $V$ of $x$, each of whose points can be connected to $x$ by a path~\ref{def:path} in $U$.
\end{definition}

\begin{definition}[logarithmic derivative]\label{def:logarithmic-derivative}
Let $f$ be a differentiable~\ref{def:differentiable} function in any sense. The \emph{logarithmic derivative} of $f$ is $f'/f$. This follows from the chain rule as the derivative of the logarithm of $f$.
\end{definition}

\begin{definition}[loop]\label{def:loop}
A \emph{loop} is a path~\ref{def:path} $p$ such that $p(0) = p(1)$. This point is called the \emph{basepoint} of the loop.
\end{definition}

\begin{definition}[Lorentz group]\label{def:lorentz-group}
The \emph{Lorentz group~\ref{def:group}} is the generalized orthogonal group~\ref{def:generalized-orthogonal-group} $\text{O}(3,1)$ (or, alternatively, $\text{O}(1,3)$ depending on the metric~\ref{def:metric-space} signature convention being used).
\end{definition}

\begin{definition}[lower integral]\label{def:lower-integral}
Let $f:Q\to\mathbb R$ be a bounded~\ref{def:bounded} function defined on the rectangle~\ref{def:rectangle} $Q$ in $\mathbb R^n$. The \emph{lower integral} of $f$ over $Q$ is the supremum~\ref{def:supremum} over all partitions~\ref{def:partition-of-a-set} of $Q$ of the lower sum~\ref{def:lower-sum} of $Q$: $$\underline{\int_Q}f = \sup_P L(f,P)$$
\end{definition}

\begin{definition}[lower sum]\label{def:lower-sum}
Let $f: Q\to \mathbb R$ be a bounded~\ref{def:bounded} function defined on the rectangle~\ref{def:rectangle} $Q$ and let $P$ be a partition of a set~\ref{def:partition-of-a-set} of $Q$. For each subrectangle~\ref{def:subrectangle} $R$ determined by $P$, let $m_R(f) = \inf\limits_{x\in R} f(x)$. The \emph{lower sum} for $f$ determined by $P$ is $$L(f,P) = \sum_{R} m_R(f)\text{vol}(R)$$ where $\text{vol}$ is the volume~\ref{def:volume} of the rectangle~\ref{def:rectangle} $R$.
\end{definition}

\begin{definition}[manifold]\label{def:manifold}
A \emph{manifold} is a topological space~\ref{def:topological-space} that is second-countable~\ref{def:second-countable}, Hausdorff~\ref{def:hausdorff}, and locally homeomorphic~\ref{def:homeomorphism} to $\mathbb R^n$.
\end{definition}

\begin{definition}[matrix Lie group]\label{def:matrix-lie-group}
A \emph{matrix Lie group~\ref{def:lie-group}} is a closed subgroup of the general linear~\ref{def:general-linear-group} group~\ref{def:group} $\text{GL}(n,\mathbb C)$.
\end{definition}

\begin{definition}[matrix logarithm]\label{def:matrix-logarithm}
The \emph{logarithm} of an $n\times n$ matrix~\ref{def:matrix} $A$ is given by $$\ln(A) = \sum_{m=1}^\infty (-1)^{m+1}\frac{(A-I)^m}{m}$$ whenever the sequence converges~\ref{def:sequence-convergence}.
\end{definition}

\begin{definition}[meager]\label{def:meager}
Let $X$ be a topological space~\ref{def:topological-space}. A set $A\subset X$ is \emph{meager} if it can be expressed as a countable~\ref{def:countable} union of nowhere dense~\ref{def:nowhere-dense} sets
\end{definition}

\begin{definition}[mean value property]\label{def:mean-value-property}
A function $f:\mathbb C\to \mathbb C$ has the \emph{mean value property} if for every $z\in \mathbb C$ and every $r>0$, $f(z)$ is equal to the complex integral~\ref{def:complex-integral} $$f(z) = \frac{1}{2\pi}\int_0^{2\pi} f(z + r\cos\theta +ir\sin\theta)\text d\theta.$$
\end{definition}

\begin{definition}[metric completion]\label{def:metric-completion}
Let $X$ be a metric space~\ref{def:metric-space}. The \emph{completion} of $X$ is the complete metric space~\ref{def:complete-metric-space} $(\tilde X, \tilde \rho)$ such that $X$ is a dense~\ref{def:dense} subset of $\tilde X$ and $\rho(x,y) = \tilde\rho(x,y)$ for all $x,y \in X$ as constructed here.
\end{definition}

\begin{definition}[metric space]\label{def:metric-space}
A \emph{metric space} is a set $X$ together with a function $\rho: X\times X \to \mathbb R$ such that for all $x,y,z \in X$, the following properties are satisfied:
1. $\rho(x,y) = 0$ if and only if $x=y$.
2. $\rho(x,y) = \rho(y,x)$.
3. $\rho(x,z) \leq \rho(x,y) + \rho(y,z)$. 

Property 3 is called the \emph{triangle inequality}and the function $\rho$ is called a \emph{metric}. The metric generates a topology~\ref{def:topological-space} on $X$ where the open~\ref{def:open} sets $G$ are those for which there exists $\varepsilon > 0$  for all $x \in G$ such that $\{y\in X\mid \rho(x,y) < \varepsilon\}\subset G$.
\end{definition}

\begin{definition}[multiplicative character]\label{def:multiplicative-character}
A \emph{multiplicative character} on a topological group~\ref{def:topological-group} $G$ is a homomorphism~\ref{def:group-homomorphism} from $G$ to the group~\ref{def:group} of units~\ref{def:unit-of-a-ring} of the field~\ref{def:field} of complex numbers $\mathbb C^\times$.
\end{definition}

\begin{definition}[neighborhood]\label{def:neighborhood}
A \emph{neighborhood} of a point $x$ in a topological space~\ref{def:topological-space} $X$ is a subset $N$ of $X$ containing $x$ which contains an open~\ref{def:open} set $U\subseteq N$ that also contains $x$.
\end{definition}

\begin{definition}[nonmeager]\label{def:nonmeager}
Let $X$ be a topological space~\ref{def:topological-space}. A set $A\subset X$ is \emph{nonmeager} if it is not meager~\ref{def:meager}.
\end{definition}

\begin{definition}[normal-form game]\label{def:normal-form-game}
A game~\ref{def:game} in \emph{normal form} is a triplet $$G = (I, \{S_i\}_{i\in I}, \{h_i\}_{i\in I})$$ where 
1. $I$ is the set of _players_. $I = \{1,\dots, n\}$ for some $n \in\mathbb N$.
2. for each $i \in I$, $S_i$ is player $i$'s set of strategies. $\prod\limits_{i\in I} S_i$ is the set of strategy profiles
3. for each $i\in I$, the function $h_i: S_i \to \mathbb R$ is the payoff function for player $i$.
\end{definition}

\begin{definition}[nowhere dense]\label{def:nowhere-dense}
Let $X$ be a topological space~\ref{def:topological-space}. A set $A\subset X$ is \emph{nowhere dense} in $X$ if for every nonempty open~\ref{def:open} subset $G\subset X$, there exists a nonempty open~\ref{def:open} subset $H\subset G$ such that $A\cap H=\emptyset$.

Equivalently, a set $A\subset X$ id \emph{nowhere dense} in $X$ if its closure~\ref{def:closure} has empty interior~\ref{def:interior}.
\end{definition}

\begin{definition}[open]\label{def:open}
A subset $G$ of a topological space~\ref{def:topological-space} $(X,\mathcal T)$ is \emph{open} if it is an element of $\mathcal T$.

When $X$ is a metric space~\ref{def:metric-space}, a subset $A \subset X$ is \emph{open} if and only if for all $a \in A$, there exists $\varepsilon > 0$ such that $B(a,\varepsilon) = \{x \in X \mid \rho(a,x) < \varepsilon\}\subset A$.
\end{definition}

\begin{definition}[open cover]\label{def:open-cover}
Let $X$ be a topological space~\ref{def:topological-space}. An \emph{open cover} of $A$ is a collection of open~\ref{def:open} subsets of $X$ such that $A$ lies in the union of the elements of the collection. That is, an open covering is a collection $\{G_\lambda\}_{\lambda\in\Lambda}$ of subsets of $X$ such that $A\subset \bigcup\limits_{\lambda\in\Lambda} G_\lambda$.
\end{definition}

\begin{definition}[opposite topology]\label{def:opposite-topology}
Let $(X,\tau)$ be a topological space~\ref{def:topological-space}. The \emph{opposite topology} on  $X$ has open sets given by the sets that are closed~\ref{def:closed} under the original topology~\ref{def:topological-space}.
\end{definition}

\begin{definition}[orientable]\label{def:orientable}
Let $M\subset\mathbb R^n$ be an embedded m-dimensional manifold with boundary~\ref{def:embedded-m-dimensional-manifold-with-boundary}. Then $M$ is \emph{orientable} if there exists an open~\ref{def:open} set $U \subset\mathbb R^n$ containing $M$ and an m-form~\ref{def:differential-k-form} $\phi$ on $U$ such that for every local parametrization $f:W\to M$ with $W \subset H^+$ open~\ref{def:open} in $H^+$, the pullback $f^*\phi$ is an m-form~\ref{def:differential-k-form} on $W$ which is nowhere zero.
\end{definition}

\begin{definition}[orientation]\label{def:orientation}
The pair $(U,\phi)$ as in the definition of orientable~\ref{def:orientable} defines an \emph{orientation} on an embedded m-dimensional manifold with boundary~\ref{def:embedded-m-dimensional-manifold-with-boundary} $M$.

Two pairs $(U,\phi)$ and $(V,\psi)$ determine the same orientation if for every local parametrization $f:W\to M$ where $W \subset H^m$ is open~\ref{def:open} in $H^m$, there exists a positive function $g:W\to\mathbb R$ such that $f^*(\phi) = gf^*\phi$.

An orientation on $M$ is an equivalence class~\ref{def:equivalence-class} if pairs that determine the same orientation on $M$.
\end{definition}

\begin{definition}[orientation-preserving]\label{def:orientation-preserving}
Let $U, V \subset\mathbb R^n$ be open~\ref{def:open} and let $f:U\to V$ be a $C^k$ diffeomorphism~\ref{def:diffeomorphism}. Then $f$ is \emph{orientation-preserving} if the determinant~\ref{def:determinant} of its total derivative~\ref{def:differentiable} is always positive: $$\text{det}(Df(x)) > 0$$ for all $x \in U$.
\end{definition}

\begin{definition}[oriented manifold]\label{def:oriented-manifold}
An embedded m-dimensional manifold with boundary~\ref{def:embedded-m-dimensional-manifold-with-boundary} that admits an orientation~\ref{def:orientation} is an \emph{oriented manifold}.
\end{definition}

\begin{definition}[orthogonal group]\label{def:orthogonal-group}
The \emph{orthogonal~\ref{def:orthogonal-matrix} group~\ref{def:group}} $\text{O}(n)$ is the group~\ref{def:group} of $n\times n$ orthogonal matrices~\ref{def:orthogonal-matrix}.

This group~\ref{def:group} is a matrix Lie group~\ref{def:matrix-lie-group} because it is a subgroup~\ref{def:subgroup} of the general linear group~\ref{def:general-linear-group} $\text{GL}(n,\mathbb C)$.
\end{definition}

\begin{definition}[outward pointing]\label{def:outward-pointing}
Let $M \subset \mathbb R^n$ be an oriented manifold~\ref{def:oriented-manifold} of dimension $m$. Then the boundary~\ref{def:boundary-points-of-a-manifold} $\partial M$ is an $(m-1)$-dimensional oriented manifold~\ref{def:oriented-manifold} (with the induced orientation). Let $\phi U\to M$ be a local parametrization of $M$ near $p \in \partial M$ so that $0 \in U$ and $\phi(0) = p$, where $U$ is open~\ref{def:subspace-topology} in the half space~\ref{def:closed-half-space} $H^m_+$. A tangent vector~\ref{def:tangent-space} $$v \in T_pM = D\phi(0)(\mathbb R^m)\subset\mathbb R^n$$ is \emph{outward pointing} if it is the image of a vector $w \in \mathbb R^m$ which does *not* lie in $H_+^m$.
\end{definition}

\begin{definition}[partial derivative]\label{def:partial-derivative}
Let $U \subset \mathbb R^n$ be open~\ref{def:open}. The \emph{partial derivative} of a function $f:U\to\mathbb R^m$ with respect to the $i$th coordinate $x_i$ at $a = (a_1,\dots,a_n) \in U$ is defined as $$\frac{\partial f}{\partial x_i}(a) = \lim_{h\to 0} \frac{f(a_1,\dots,a_i+h,\dots,a_n) - f(a)}{h}.$$
\end{definition}

\begin{definition}[partition of a rectangle]\label{def:partition-of-a-rectangle}
A \emph{partition} $P$ of a closed~\ref{def:closed} interval~\ref{def:interval} $[a,b]$ is a finite collection of points in $[a,b]$ such that $a,b \in P$.

More generally, a \emph{partition} $P$ of a rectangle~\ref{def:rectangle} $Q$ in $\mathbb R^n$ is a collection $P = \{P_i\}_{i=1}^n$ of partitions of the component intervals~\ref{def:component-interval} $[a_i,b_i]$ of $Q$.
\end{definition}

\begin{definition}[partition of unity]\label{def:partition-of-unity}
Let $\mathcal A$ be a collection of open~\ref{def:open} sets in $\mathbb R^n$ and let $A$ be their union. A \emph{partition of unity subordinate to $\mathcal A$} is a countable collection of smooth~\ref{def:smooth} functions $\{\phi_i\}_{i\in I}$ such that $\phi_i:\mathbb R^n \to \mathbb R$ such that 
1. each $\phi_i \geq 0$.
2. each $\phi_i$ has compact support~\ref{def:compact-support} contained in some $U \in \mathcal A$.
3. for all points $a$ of $A$, there exists an open~\ref{def:open} neighborhood $V$ of $a$ such that all but finitely many of the $\phi_i$ are zero on $V$.
4. for all points $a$ in $A$, $\sum\limits_{i\in I} \phi_i(a) = 1$.
\end{definition}

\begin{definition}[path]\label{def:path}
Let $X$ be a topological space~\ref{def:topological-space}. A \emph{path} in $X$ is a continuous~\ref{def:continuous} function from the unit interval~\ref{def:interval} $I$ to $X$.
\end{definition}

\begin{definition}[path-connected]\label{def:path-connected}
A topological space~\ref{def:topological-space} $X$ is \emph{path-connected} if for every pair of points $x,y \in X$, there exists a continuous~\ref{def:continuous} function $\gamma:[0,1]\to X$ such that $\gamma(0)=x$ and $\gamma(1)=y$.
\end{definition}

\begin{definition}[perfect set]\label{def:perfect-set}
A subset $S$ of a topological space~\ref{def:topological-space} is \emph{perfect} if it is closed~\ref{def:closed} and contains no isolated points~\ref{def:isolated-point}. 

Equivalently, a set $S$ is \emph{perfect} if it is equal to the set of its limit points~\ref{def:limit-point}.
\end{definition}

\begin{definition}[Pontryagin dual]\label{def:pontryagin-dual}
Let $G$ be a locally compact~\ref{def:locally-compact-group} abelian topological~\ref{def:topological-group} group~\ref{def:group}. The \emph{Pontryagin dual} of $G$, denoted $\widehat G$, is the group~\ref{def:group} $G$ of continuous~\ref{def:continuous} homomorphisms~\ref{def:group-homomorphism} from $G$ to the circle group~\ref{def:circle-group} $\mathbb T$.
\end{definition}

\begin{definition}[principal part]\label{def:principal-part}
The \emph{principal part} of a Laurent series~\ref{def:laurent-series} $$f(z) = \sum_{n=-\infty}^\infty a_n(z-c)^n$$ is the series of terms of negative degree: $$\sum_{n=-\infty}^{-1}a_n(z-c)^k.$$
\end{definition}

\begin{definition}[product topology]\label{def:product-topology}
Let $X$ and $Y$ be topological spaces~\ref{def:topological-space}. The topology~\ref{def:topological-space} on their Cartesian product $X\times Y$ has basis~\ref{def:topological-basis} the products $U\times V$ of open~\ref{def:open} sets $U\subset X$ and $V\subset Y$.

Given a (possibly infinite) set $\{X_i\}_{i\in I}$ of topological spaces~\ref{def:topological-space}, the \emph{product topology} on $\prod_{i\in I} X_i$ is generated~\ref{def:generate-a-topology} by the basis consisting of all products $\prod_{i\in I} U_i$ for $U$ open~\ref{def:open} in $X$ and $U_i = X_i$ for all but finitely many $i\in I$.
\end{definition}

\begin{definition}[quotient of topological spaces]\label{def:quotient-of-topological-spaces}
The quotient of a topological space~\ref{def:topological-space} $X$ by a subspace~\ref{def:subspace-topology} $Y\subseteq X$ is the quotient by~\ref{def:quotient-by-equivalence-relation} the equivalence relation~\ref{def:equivalence-relation} $\sim$ given by $y\sim y'$ if $y$ and $y'\in Y$. 

Let $q: X\to X/Y$ be given by $x\mapsto [x]$. This function is (well-defined and) surjective~\ref{def:surjective}. The topology~\ref{def:topological-space} on the quotient $X/Y$ is the collection of $U\subset X/Y$ such that the inverse image $q^{-1}(U)$ is open~\ref{def:open} in $X$.
\end{definition}

\begin{definition}[radius of convergence]\label{def:radius-of-convergence}
The \emph{radius of convergence} of a power series~\ref{def:power-series} $$f(x) = \sum_{n=0}^\infty a_n(x-c)^n$$over a Banach algebra~\ref{def:banach-algebra} with norm~\ref{def:norm} $\vert \cdot\vert $ is a real number $r$ such that the series converges to $f(x)$ for all $\vert x-a\vert <r$.
\end{definition}

\begin{definition}[Radon measure]\label{def:radon-measure}
Let $X$ be a Hausdorff~\ref{def:hausdorff} topological space~\ref{def:topological-space} and let $(X, \Sigma, \mu)$ be a measure space~\ref{def:measure-space} where $\Sigma$ is the Borel σ-algebra~\ref{def:borel-σ-algebra} on $X$. The measure~\ref{def:measure-space} $\mu$ is a \emph{Radon measure} if it is inner regular and locally finite~\ref{def:locally-finite-measure}.
\end{definition}

\begin{definition}[real symplectic group]\label{def:real-symplectic-group}
The \emph{real symplectic group} $\text{Sp}(2n,\mathbb R)$ is the group~\ref{def:group} of real $2n\times 2n$ matrices~\ref{def:matrix} that preserve the skew-symmetric bilinear~\ref{def:multilinear} form $$\omega(x,y) = \langle x,\Omega y\rangle$$ where $\langle\cdot,\cdot\rangle$ is the Euclidean inner product~\ref{def:euclidean-inner-product} and $$\Omega = \begin{bmatrix}0 & I_n \\ -I_n & 0\end{bmatrix}$$ for $I_n$ the $n\times n$ identity matrix~\ref{def:identity-matrix}. That is, $A\in \text{Sp}(2n, \mathbb R)$ if and only if $\omega(Ax,Ay) = \omega(x,y)$ for all $x,y \in \mathbb R^n$. Equivalently,  $A\in \text{Sp}(2n, \mathbb R)$ if and only if $$-\Omega A^T \Omega = A^{-1}$$ where , as usual, $A^T$ is the transpose~\ref{def:matrix-transpose} of $A$ and $A^{-1}$ is its inverse~\ref{def:inverse-matrix}.

This group~\ref{def:group} is a matrix Lie group~\ref{def:matrix-lie-group} because it is a subgroup~\ref{def:subgroup} of the general linear group~\ref{def:general-linear-group} $\text{GL}(n,\mathbb C)$.
\end{definition}

\begin{definition}[rectifiable]\label{def:rectifiable}
Let $S \subset\mathbb R^n$ be a bounded~\ref{def:bounded} set. Then $S$ is \emph{rectifiable} if the constant function $1$ is Riemann integrable~\ref{def:riemann-integrable} over $S$.
\end{definition}

\begin{definition}[regular family]\label{def:regular-family}
A collection $\mathcal V$ of measurable~\ref{def:measurable} subsets of $\mathbb R^n$ is a \emph{regular family} if there exists $C > 0$ such that the diameter~\ref{def:diameter-of-a-set} of each $V \in \mathcal V$ satisfies the following inequality: $$\text{diam}(V)^n \leq C\mu(V).$$
\end{definition}

\begin{definition}[residue]\label{def:residue}
The \emph{residue} of a meromorphic function $f: \mathbb C\setminus \{a_i\}_{i\in I}\to \mathbb C$ at $a_i$ is the $-1$st coefficient of the Laurent series~\ref{def:laurent-series} for $f$.
\end{definition}

\begin{definition}[Riemann integrable]\label{def:riemann-integrable}
A bounded~\ref{def:bounded} function $f:Q\to\mathbb R$ defined on the rectangle~\ref{def:rectangle} $Q$ in $\mathbb R^n$ is \emph{integrable} if the lower integral~\ref{def:lower-integral} of $f$ is equal to the upper integral~\ref{def:upper-integral} of $f$. The \emph{integral} of $f$ over $Q$ is defined to be $$\int_Q f = \underline{\int_Q} f = \overline{\int_Q} f.$$

Alternatively, let $f: S\to \mathbb R$ be a bounded~\ref{def:bounded} function from the bounded~\ref{def:bounded} set $S$ to $\mathbb R$, and let $f_S$ be the extension of $S$ in the usual sense to $\mathbb R$. Choose $Q$ a rectangle~\ref{def:rectangle} containing $S$. The \emph{integral} of this function over $S$ is $$\int_S f = \int_Q f_S,$$ if it exists.
\end{definition}

\begin{definition}[Schwartz space]\label{def:schwartz-space}
Let $F$ be a local field~\ref{def:field} with a non-Archimedean~\ref{def:non-archimedean-absolute-value} absolute value~\ref{def:absolute-value}. Then the \emph{Schwartz space}, or the collection of \emph{Schwartz-Bruhat functions}, consists of the functions from $F$ to $\mathbb C$ that are locally constant and have compact support~\ref{def:compact-support}.

It is often either denoted $C_c^\infty(F)$ or $\mathcal S(F)$.
\end{definition}

\begin{definition}[second-countable]\label{def:second-countable}
A set $X$ is \emph{second-countable} if there exists a countable collection $\mathcal A$ of open~\ref{def:open} sets in $X$ such that every open~\ref{def:open} set is a union of sets of $A$.
\end{definition}

\begin{definition}[semi-locally simply connected]\label{def:semi-locally-simply-connected}
A topological space~\ref{def:topological-space} $X$ is \emph{semi-locally simply connected} if every point $b\in B$ has a neighborhood~\ref{def:neighborhood} $U$ such that the map $\pi_1(U,b)\to \pi_1(B,b)$ between fundamental groups~\ref{def:fundamental-group} is the trivial group homomorphism~\ref{def:group-homomorphism}.
\end{definition}

\begin{definition}[separable metric space]\label{def:separable-metric-space}
A metric space~\ref{def:metric-space} is \emph{separable} if it contains a countable, dense~\ref{def:dense} subset.
\end{definition}

\begin{definition}[sequence convergence]\label{def:sequence-convergence}
A sequence $\{x_i\}_{i\in\mathbb N}$ in a topological space~\ref{def:topological-space} $X$ \emph{converges} to $x \in X$ if for every open~\ref{def:open} $U\subset X$ such that $x \in U$, there exists $n_0 \in\mathbb N$ such that for every $n \geq n_0$, we have $x_n \in U$.

In a general topological space~\ref{def:topological-space}, (i.e. in some non-Hausdorff spaces) the limit of a sequence need not be unique!
\end{definition}

\begin{definition}[sequentially compact]\label{def:sequentially-compact}
A topological space~\ref{def:topological-space} $X$ is \emph{sequentially compact} if every sequence in $X$ has a convergent~\ref{def:sequence-convergence} subsequence.
\end{definition}

\begin{definition}[simple region]\label{def:simple-region}
Let $C$ be a compact~\ref{def:compact}, rectifiable~\ref{def:rectifiable} set in $\mathbb R^{n-1}$. Let $\phi,\psi: C \to\mathbb R$ be continuous~\ref{def:continuous} functions such that $\phi(x) \leq \psi(x)$ for all $x \in C$. Then the set $$S = \{(x,t)\mid x\in C, \phi(x) \leq t \leq \psi(x)\} \subset \mathbb R^{n-1}\times \mathbb R$$ is called a \emph{simple region} in $\mathbb R^n$.
\end{definition}

\begin{definition}[simply connected]\label{def:simply-connected}
A topological space~\ref{def:topological-space} $X$ is \emph{simply connected} if it is path-connected~\ref{def:path-connected} and for every two continuous~\ref{def:continuous} maps $\gamma_1:[0,1]\to X$ and $\gamma_2:[0,1]\to X$ there exists a homotopy $F:[0,1]\times[0,1]\to X$ such that $F(x,0)= \gamma_1(x)$ and $F(x,1) = \gamma_2(x)$.  

A topological space~\ref{def:topological-space} $X$ is \emph{simply connected} if it is path-connected~\ref{def:path-connected} and has trivial fundamental group~\ref{def:fundamental-group}.
\end{definition}

\begin{definition}[smooth]\label{def:smooth}
A function of class~\ref{def:class} $C^{\infty}$ is \emph{smooth}. That is, a smooth function has continuous~\ref{def:continuous} total derivative~\ref{def:differentiable} of all orders.
\end{definition}

\begin{definition}[smooth topological m-dimensional manifold with boundary]\label{def:smooth-topological-m-dimensional-manifold-with-boundary}
A \emph{smooth topological $m$-dimensional manifold with boundary} is a topological m-dimensional manifold with boundary~\ref{def:topological-m-dimensional-manifold-with-boundary} $M$ together with a collection $\mathcal A$ of charts~\ref{def:chart} on $M$ such that 
1. for all $p\in M$, there exists a chart~\ref{def:chart} $(U, \phi)\in \mathcal A$ such that $p \in U$;
2. for any two charts~\ref{def:chart} $(U,\phi)$ and $(V, \psi) \in \mathcal A$, the transition function~\ref{def:transition-function} $$\psi\circ \phi^{-1}:\phi(U\cap V)\to\psi(U\cap V)$$ is a smooth~\ref{def:smooth} diffeomorphism~\ref{def:diffeomorphism}.
3. the set $\mathcal A$ (called an atlas~\ref{def:atlas}) is the maximal set of collections of charts~\ref{def:chart} satisfying 1. and 2.
\end{definition}

\begin{definition}[special linear group]\label{def:special-linear-group}
The \emph{special linear group~\ref{def:group}} $\text{SL}(n,\mathbb C)$ (or, over the reals, $\text{SL}(n,\mathbb R)$) is the group~\ref{def:group} of $n\times n$ complex (or real) valued matrices with determinant~\ref{def:determinant} $1$ together with the operation~\ref{def:binary-operation} of matrix multiplication~\ref{def:matrix-multiplication}.

This group~\ref{def:group} is a matrix Lie group~\ref{def:matrix-lie-group} because it is a subgroup~\ref{def:subgroup} of the general linear group~\ref{def:general-linear-group} $\text{GL}(n,\mathbb C)$.
\end{definition}

\begin{definition}[special orthogonal group]\label{def:special-orthogonal-group}
The \emph{special orthogonal group~\ref{def:orthogonal-group}} $\text{SO}(n)$ is the group~\ref{def:group} of $n\times n$ orthogonal~\ref{def:orthogonal-matrix} matrices~\ref{def:matrix} with unit determinant~\ref{def:determinant}.

This group~\ref{def:group} is a matrix Lie group~\ref{def:matrix-lie-group} because it is a subgroup~\ref{def:subgroup} of the general linear group~\ref{def:general-linear-group} $\text{GL}(n,\mathbb C)$.
\end{definition}

\begin{definition}[special unitary group]\label{def:special-unitary-group}
The \emph{special unitary group~\ref{def:unitary-group}} $\text{SU}(n)$ is the group~\ref{def:group} of unitary~\ref{def:unitary-matrix} matrices~\ref{def:matrix} with unit determinant~\ref{def:determinant}.

This group~\ref{def:group} is a matrix Lie group~\ref{def:matrix-lie-group} because it is a subgroup~\ref{def:subgroup} of the general linear group~\ref{def:general-linear-group} $\text{GL}(n,\mathbb C)$.
\end{definition}

\begin{definition}[subbasis]\label{def:subbasis}
Let $(X,\tau)$ be a topological space~\ref{def:topological-space}. A \emph{subbasis} $S$ for the topology~\ref{def:topological-space} $\tau$ is a subset of $\tau$ such that every  set in $\tau$ is a union of finite intersection of sets in $S$.
\end{definition}

\begin{definition}[subspace topology]\label{def:subspace-topology}
Let $A\subset X$ be a subset of the topological space~\ref{def:topological-space} $X$. The \emph{subspace topology} on $A$ is the topology in which the open~\ref{def:open} sets are given by all sets $A \cap U$ where $U\subset X$ is open~\ref{def:open} in $X$.
\end{definition}

\begin{definition}[symmetric derivative]\label{def:symmetric-derivative}
Let $\nu$ be a Borel~\ref{def:borel-σ-algebra} measure~\ref{def:measure-space} on $\mathbb R^n$ (i.e. the underlying measure space~\ref{def:measure-space} is $(\mathbb R^n, \mathcal B, \nu)$). The \emph{symmetric derivative} of $\nu$ at $x \in \mathbb R^n$ is $$(D\nu)(x) = \lim_{r\to 0} \frac{\nu(B(x,r))}{\mu(B(x,r))}$$ if this limit exists. This is the derivative ( todo: link to 208 notions?) in the symmetric differential basis.
\end{definition}

\begin{definition}[T0 space]\label{def:t0-space}
A topological space~\ref{def:topological-space} is $T_0$ if if for any two points $x,y \in X$, there exists an open~\ref{def:open} neighborhood~\ref{def:neighborhood} of one that does not contain the other.
\end{definition}

\begin{definition}[T1 space]\label{def:t1-space}
A topological space~\ref{def:topological-space} is $T_1$ if singletons are closed~\ref{def:closed}.
\end{definition}

\begin{definition}[tangent space]\label{def:tangent-space}
Let $M$ be an embedded m-dimensional manifold~\ref{def:embedded-m-dimensional-manifold}. The \emph{tangent space} of $M$ at $x$, denoted $T_xM$, is the vector space~\ref{def:vector-space} given by the kernel~\ref{def:kernel-of-linear-transformation} of the total derivative~\ref{def:differentiable} of the function given in the definition of the kind of manifold~\ref{def:manifold}.
\end{definition}

\begin{definition}[topological basis]\label{def:topological-basis}
A \emph{basis} for a topology~\ref{def:topological-space} $\tau$ on a set $X$ is a collection $\mathcal B$ of subsets of $X$ such that 
1. All points are contained in at least one element of $\mathcal B$, and
2. If $x\in B'\cap B''$ then there exists at least one $B\in \mathcal B$ such that $x\in B \subset B'\cap B''$.
\end{definition}

\begin{definition}[topological field]\label{def:topological-field}
A \emph{topological field} is a field~\ref{def:field} $F$ that is a topological ring~\ref{def:topological-ring} and also whose inverse maps $F\times F^\times\to F$ for both binary operations~\ref{def:binary-operation} are also continuous~\ref{def:continuous}.
\end{definition}

\begin{definition}[topological game]\label{def:topological-game}
A \emph{topological game} $G(X, Y, \mathcal S)$ is a game~\ref{def:game} of perfect information played between two players on a topological space~\ref{def:topological-space} $X$ where $Y$ is a fixed subset of $Y$, and $\mathcal S$ is some family of subsets of $X$.
\end{definition}

\begin{definition}[topological group]\label{def:topological-group}
A \emph{topological group} is a topological space~\ref{def:topological-space} $G$ that is also a group~\ref{def:group} such that the binary operation~\ref{def:binary-operation} on the group~\ref{def:group} and the inverse map are both continuous~\ref{def:continuous}. That is, the two maps $$\cdot:G\times G\to G \text{ given by } (x,y)\mapsto xy$$$$\text{inv}:G\to G \text{ given by } x\mapsto x^{-1}$$ are both continuous~\ref{def:continuous} where $G\times G$ is given the product topology~\ref{def:product-topology}.
\end{definition}

\begin{definition}[topological m-dimensional manifold with boundary]\label{def:topological-m-dimensional-manifold-with-boundary}
A \emph{topological $m$-dimensional manifold with boundary} is a topological space~\ref{def:topological-space} $M$ such that 
1. for all $p \in M$, there exists an open~\ref{def:open} neighborhood $U$ of $p$ homeomorphic~\ref{def:homeomorphism} to an open~\ref{def:open} set in the closed half-space~\ref{def:closed-half-space} $H^m$; 
2. $M$ is Hausdorff~\ref{def:hausdorff};
3. $M$ is second-countable~\ref{def:second-countable}.
\end{definition}

\begin{definition}[topological ring]\label{def:topological-ring}
A \emph{topological ring} is a ring~\ref{def:ring} $R$ together with a topology~\ref{def:topological-space} on the underlying set where the addition map $+:R\times R\to R$ (given by $+(a,b) = a+b$ for all $a,b\in R$) and the multiplication map $\cdot:R\times R\to R$ (given by $\cdot(a,b)= ab$ for all $a,b\in R$) are both continuous~\ref{def:continuous}.
\end{definition}

\begin{definition}[topological space]\label{def:topological-space}
A \emph{topological space} is a set $X$ together with a collection $T$ of subsets of $X$. The elements of $T$ are called the open~\ref{def:open} sets and they satisfy the following properties:
1. $X\in T$ and $\emptyset \in T$.
2. Arbitrary unions of open sets are open.
3. Finite intersections of open sets are open.

$T$ is called a \emph{topology} on $X$.
\end{definition}

\begin{definition}[topology of the union]\label{def:topology-of-the-union}
Let $X$ and $Y$ be topological spaces~\ref{def:topological-space}. The topology~\ref{def:topological-space} on their disjoint union $X\amalg Y$ has as open~\ref{def:open} sets the (disjoint) unions of sets that are open~\ref{def:open} in $X$ and $Y$.
\end{definition}

\begin{definition}[totally bounded]\label{def:totally-bounded}
A subset $E$ of the metric space~\ref{def:metric-space} $(X,\rho)$ is \emph{totally bounded} if for all $\varepsilon > 0$, there exists a finite set $\{x_i\}_{i=1}^n$ of points in $X$ such that $$E \subseteq \bigcup_{i=1}^n B(x_i, \varepsilon),$$ where $B(x_i, \varepsilon) = \{x \in X \mid \rho(x_i,x)< \varepsilon\}.$ That is, for every $\varepsilon > 0$, $E$ can be covered by a finite set of balls of radius $\varepsilon$.
\end{definition}

\begin{definition}[totally disconnected]\label{def:totally-disconnected}
A topological space~\ref{def:topological-space} $X$ is \emph{totally disconnected} if the connected components~\ref{def:connected-component} of $X$ are exactly the singletons.
\end{definition}

\begin{definition}[transition function]\label{def:transition-function}
The \emph{transition function} of two charts~\ref{def:chart} $(U,\phi)$ and $(V, \psi)$ on a manifold~\ref{def:manifold} $M$ is the composition $\psi \circ \phi^{-1}: \phi(U\cap V)\to \psi(U\cap V)$.
\end{definition}

\begin{definition}[typical]\label{def:typical}
A property is \emph{typical} if it holds for every element of a residual~\ref{def:comeager} set.
\end{definition}

\begin{definition}[uniformly continuous]\label{def:uniformly-continuous}
A function $f:X\to Y$ between metric spaces~\ref{def:metric-space} $(X,\rho)$ and $(Y, \sigma)$ is \emph{uniformly continuous} if for all $x,y \in X$ and all $\varepsilon > 0$, there exists $\delta > 0$ such that $\sigma(f(x), f(y)) < \varepsilon$ whenever $\rho(x,y) < \delta$.
\end{definition}

\begin{definition}[unitary group]\label{def:unitary-group}
The \emph{unitary~\ref{def:unitary-matrix} group~\ref{def:group}} $\text{U}(n)$ is the group~\ref{def:group} of $n\times n$ unitary~\ref{def:unitary-matrix} matrices together with the operation of matrix multiplication~\ref{def:matrix-multiplication}.

This group~\ref{def:group} is a matrix Lie group~\ref{def:matrix-lie-group} because it is a subgroup~\ref{def:subgroup} of the general linear group~\ref{def:general-linear-group} $\text{GL}(n,\mathbb C)$.
\end{definition}

\begin{definition}[unitary group representation]\label{def:unitary-group-representation}
A complex group representation~\ref{def:group-representation} $\rho: G\to GL(V)$ is \emph{unitary} if $\rho(g)$ is unitary~\ref{def:unitary-group} for all $g\in G$.  Explicitly, it is \emph{unitary} if the Hermitian adjoint~\ref{def:adjoint-of-a-linear-transformation} $\rho(g)^* =  \rho(g)^{-1} = \rho(g^{-1})$ for all $g\in G$.
\end{definition}

\begin{definition}[upper integral]\label{def:upper-integral}
Let $f:Q\to\mathbb R$ be a bounded~\ref{def:bounded} function defined on the rectangle~\ref{def:rectangle} $Q$ in $\mathbb R^n$. The \emph{upper integral} of $f$ over $Q$ is the infimum~\ref{def:infimum} over all partitions~\ref{def:partition-of-a-set} of $Q$ of the upper sum~\ref{def:upper-sum} of $Q$: $$\overline{\int_Q}f = \inf_P U(f,P)$$
\end{definition}

\begin{definition}[upper sum]\label{def:upper-sum}
Let $f: Q\to \mathbb R$ be a bounded~\ref{def:bounded} function defined on the rectangle~\ref{def:rectangle} $Q$ and let $P$ be a partition of a set~\ref{def:partition-of-a-set} of $Q$. For each subrectangle~\ref{def:subrectangle} $R$ determined by $P$, let $M_R(f) = \sup\limits_{x\in R} f(x)$. The \emph{lower sum} for $f$ determined by $P$ is $$U(f,P) = \sum_{R} M_R(f)\text{vol}(R)$$ where $\text{vol}$ is the volume~\ref{def:volume} of the rectangle~\ref{def:rectangle} $R$.
\end{definition}

\begin{definition}[Vitali covering]\label{def:vitali-covering}
Let $X$ be a metric space~\ref{def:metric-space} with $A \subset X$. Let $\mu$ be the Lebesgue measure~\ref{def:lebesgue-measure}. A collection $\mathcal V$ of sets covering $A$ is a \emph{Vitali covering} if for each $x \in A$ and $\varepsilon > 0$, there exists $U \in \mathcal V$ such that $x \in U$, and the diameter~\ref{def:diameter-of-a-set} of $U$ is nonzero and less than $\varepsilon$.
\end{definition}

\begin{definition}[volume of a rectifiable set]\label{def:volume-of-a-rectifiable-set}
The volume~\ref{def:volume} of a rectifiable~\ref{def:rectifiable} set $S$ is $\text{vol}(S) = \int_S 1$, or the integral~\ref{def:riemann-integrable} of the constant function over $S$.
\end{definition}

\begin{definition}[volume of compact topological space]\label{def:volume-of-compact-topological-space}
Let $X$ be a compact~\ref{def:compact} topological space~\ref{def:topological-space} and let $1_X$ be the constant function $1$ on this space. Then for a fixed Haar integral~\ref{def:haar-integral} $\int_X$, define the \emph{volume} of $X$ to be $$\text{vol}(X) = \int_X 1_X.$$
\end{definition}

\begin{definition}[weak Hausdorff]\label{def:weak-hausdorff}
A topological space~\ref{def:topological-space} $X$ is \emph{weak Hausdorff} if for every map $g:K\to X$ from a compact~\ref{def:compact} Hausdorff~\ref{def:hausdorff} space~\ref{def:topological-space} $K$ is closed~\ref{def:closed} in $X$. 

Note that $g(K)$ is itself Hausdorff~\ref{def:hausdorff}, ( todo: why?) so $g(K)$ is a compact~\ref{def:compact} subspace~\ref{def:subspace-topology} of $X$.
\end{definition}

\begin{definition}[wedge sum]\label{def:wedge-sum}
Let $(X,x_0)$ and $(Y,y_0)$ be based topological spaces~\ref{def:based-topological-space}. The \emph{wedge sum} $X\vee Y$ is the quotient~\ref{def:quotient-of-topological-spaces} of the disjoint union $X\amalg Y$ given by identifying $x_0$ and $y_0$ to a single point.
\end{definition}

\begin{definition}[winding number]\label{def:winding-number}
Let $z_0\in \mathbb C$ and let $\gamma:[0,1]\to \mathbb C\setminus \{z_0\}$ be a continuous~\ref{def:continuous} loop~\ref{def:loop} that does not pass through $z_0$. The \emph{winding number} of $\gamma$ around $\alpha$ is the complex line integral~\ref{def:complex-line-integral} $$\text{Ind}_\gamma(\alpha) = \frac{1}{2\pi i}\int_\gamma \frac{1}{\zeta - z_0} \text d\zeta,$$
a special case of the Cauchy integral formula.
\end{definition}

\end{document}